% Integration fragment. Root may renumber AG.* into Part A.
\subsection{Correspondences, character orbits, and integral group-ring tests}

\begin{claim}{AG.1}{Proved; a precise rational correspondence statement}
Let $L$ be a finite field of characteristic $p$, let $C/L$ be smooth,
proper and geometrically connected of genus at least two, and let
$P\ne Q$ in $C(L)$ have conjugate sections in the geometrically
prime-to-$p$ arithmetic fundamental group. Let
$\pi:D\to C$ be a connected finite etale cover with prime-to-$p$ geometric
Galois closure, let $P'\in D(L)$ lie above $P$, and let
$\phi:D\to X_0$ be an $L$-morphism to a smooth proper geometrically
connected curve of genus at least one. Assume
$J_{X_0}(L)[p]=0$. Then a lift $Q'\in D(L)$ supplied by R1 satisfies
$\phi(Q')=\phi(P')$. Moreover,
\[
 [\phi_*\pi^*(Q)-\phi_*\pi^*(P)]=0\quad\hbox{in }J_{X_0}(L).
\]
If $\pi$ is a torsor under a constant finite group $G$ over $L$,
then $P'$ makes the entire $P$-fibre rational, the $Q$-fibre is also
rational, and the stronger equality of effective divisors holds:
\[
             \phi_*\pi^*(Q)=\phi_*\pi^*(P).
\]
For a general cover the same divisor equality holds if the geometric
$p$-rank of $J_{X_0}$ is zero, or if $\deg(\pi)$ is strictly less
than the geometric gonality of $X_0$.
\end{claim}
\begin{proof}
R1 gives $[Q'-P']\in J_D(L)[p^\infty]$. The homomorphism
$\phi_*:J_D\to J_{X_0}$, defined over $L$ by pushforward of
degree-zero divisors, sends this class to zero, since a finite
abelian group with no element of order $p$ has no nonzero
$p$-primary element. On a smooth proper curve of positive genus,
the equality of the degree-one divisor classes of two points
forces the points to coincide: otherwise a function with divisor
their difference gives a degree-one map to $\mathbf P^1$.
This proves the first assertion.

On the base curve, R1 already gives $[Q-P]\in J_C(L)[p^\infty]$.
Apply the $L$-homomorphism $\phi_*\pi^*:J_C\to J_{X_0}$ to
obtain the asserted equality of divisor classes. For the constant
$G$-torsor, choose the lifts $P',Q'$ with equal section classes.
Their translates $gP',gQ'$ are $L$-rational and have $p$-primary
difference for every $g\in G$. Applying the preceding argument
to each pair and summing gives equality of effective divisors.
For geometric $p$-rank zero we may instead pass to a finite field
over which the Galois closure and all lifts are rational; the
target still has no geometric $p$-primary points. Equality of
the two geometric divisors descends to $L$.
Finally, distinct effective linearly equivalent divisors of degree
$N$ give, after cancelling their common part, a nonconstant function
of degree at most $N$. This proves the asserted gonality alternative.

The distinction between divisor classes and effective divisors
is necessary in the proof. One rational lift and the condition
$J_{X_0}(L)[p]=0$ do not assert that
$J_{X_0}(L')[p]=0$ for residue extensions $L'/L$ of the other
lifts. We therefore do not use that argument for the unrestricted
divisor assertion in lead L5.

The map $P\mapsto\phi_*\pi^*P$ to $\operatorname{Sym}^{\deg\pi}X_0$
is a morphism: etale locally on $C$, the finite etale fibre splits
into sections, whose images give an ordered tuple; quotienting by
the permutation action descends it. If $\phi$ is nonconstant, this
map is nonconstant, since a constant effective divisor would contain
every point of $\phi(D)$ in its finite support. Its equal-fibre locus
is consequently a proper closed subset of $C\times C$. This alone
does not exclude a hypothetical collision lying in that subset.
\end{proof}

\paragraph{Consequences of point separation and the finite spectral test.}
By AG40--AG42, A1 holds, so the unrestricted collision-dependent
implications in L5 hold as well: there is no pair of distinct
points satisfying their premise. In particular the unqualified
effective-divisor conclusion is no longer an open assertion.
The rationality qualifications above describe what the direct
correspondence argument proves without that theorem; they do
not weaken the now-established implication.

For clarity, the proposed spectral test has the following exact
meaning. Let $D/\mathbf F_q$ be smooth, proper and geometrically
connected of genus $h\ge1$, and write
\[
 P_D(T)=\det(1-TF\mid H^1_{\mathrm{et}}
       (D_{\overline{\mathbf F}_q},\mathbf Q_\nu)),
 \qquad
 H_D(X)=X^{2h}P_D(X^{-1}),\quad \nu\ne p.
\]
Then $J_D$ has a nonzero supersingular isogeny factor if and
only if one root $\alpha$ of $H_D$ has
$\alpha/\sqrt q$ a root of unity. Indeed, on a simple
isogeny factor this is exactly the supersingular
Newton-slope criterion; see
\cite[\S\S21.20--21.21(3)]{AC_Oort2008}. If one root has
this property, every conjugate in its simple factor does,
so that factor is supersingular. Conversely, a
supersingular factor supplies such roots.

This is a finite exact test: it is equivalent to
\[
 \deg\gcd\bigl(H_D(X),X^{2n}-q^n\bigr)>0
 \quad\hbox{for some }1\le n\le8h^2.
\]
To prove the bound, suppose $\beta=\alpha^2/q$ has exact
root-of-unity order $n$. Then
$\varphi(n)=[\mathbf Q(\beta):\mathbf Q]\le2h$.
For $n=\prod_\ell\ell^{a_\ell}$,
\[
 \frac{\varphi(n)^2}{n}
       =\prod_{\ell\mid n}\ell^{a_\ell-2}(\ell-1)^2
       \ge\frac12.
\]
Every odd-prime factor is at least one, and the factor
at two is at least $1/2$. Thus $n\le8h^2$.
Conversely, any nonconstant displayed greatest common
divisor gives $\alpha^{2n}=q^n$, which is precisely the
required root-of-unity condition.

For a constant abelian cover $D\to C$, the factorization
$P_D=\prod_\chi L(C,\chi,T)$ makes this the character
test in L5, with the trivial character included when
testing the whole Jacobian. Omitting it tests only the
new factor. A constant nonabelian cover requires all
irreducible Artin factors with their multiplicities;
a non-Galois cover requires its actual permutation
representation, or simply the polynomial $P_D$ itself.
The fixed-cover test supplies no existence theorem for
an etale cover having such a factor and no assertion
that the resulting maps jointly separate points.
Universal prime-to-$p$ supersingular correspondences
therefore remain the separate existence problem in A5.

\begin{claim}{AG.2}{Proved; the closure operations in L6}
Let $\mathcal N$ be the class of smooth proper curves of genus at
least two over $\overline{\mathbf F}_p$ for which the point-section
map is injective over every finite field of definition and every
finite extension thereof. It has the following closure properties.
If $D\to C$ is a connected finite etale cover with prime-to-$p$
geometric Galois closure and $D\in\mathcal N$, then
$C\in\mathcal N$. If $C$ has a family of morphisms to curves
in $\mathcal N$ which jointly separates geometric points, then
$C\in\mathcal N$. One may also use positive-genus targets
with geometric $p$-rank zero.
\end{claim}
\begin{proof}
Suppose a collision exists on $C$ over a finite field. The finite
cover and a chosen geometric lift of $P$ are defined over a finite
extension $L$. Equality of the section classes restricts along
the injective map $G_L\hookrightarrow G_k$. In the finite
prime-to-$p$ fibre functor, equality of the two sections identifies
the Galois actions on the fibres. The fixed lift of $P$ therefore
gives a fixed lift of $Q$; choosing the corresponding paths gives
equal section classes upstairs. Equivalently this is the section
lifting step in R1, before abelianization. The geometric group
upstairs is the open subgroup with its full prime-to-$p$ topology:
cores of finite-index subgroups have prime-to-$p$ index because
the geometric Galois closure does. Thus the equality is in the
full geometrically prime-to-$p$ group of $D$, not merely a
quotient. The two lifts remain distinct because their images are
distinct. This contradicts $D\in\mathcal N$.

For the second assertion, a morphism of curves induces a
continuous map of fundamental extensions over the same Galois
group. It carries conjugate sections to conjugate sections.
Every target therefore identifies the two images of a hypothetical
collision. Joint separation contradicts their distinctness.
For a target with geometric $p$-rank zero, the Abel map and R1
give the required target injectivity directly.

These operations can be iterated in the tower. They do not
assert that every curve is obtained from the given seeds.
Indeed, the universal point-separating-cover statement of A5 is
equivalent to A1: this proof gives one implication, and, assuming
A1, the identity cover and identity morphism give the other.
The later AG40--AG42 prove A1, so the general
point-separating-cover assertion of A5 follows by this
identity-cover argument. The asserted universal correspondence
to a supersingular curve is a stronger sufficient construction,
whose existence is a separate question.
\end{proof}

\begin{claim}{AG.3}{Proved; the exact Frobenius-orbit formula}
Let $C$ be a smooth proper genus-$g$ curve over an algebraically
closed field $\Omega$ of characteristic $p>0$, and let $L$ be a
line bundle of exact order $n>1$, with $(n,p)=1$. Let
$D\to C$ be the associated connected cyclic etale cover. For
$i\in\mathbf Z/n\mathbf Z$ put $V_i=H^1(C,L^i)$, with a fixed
trivialization of $L^n$. Absolute Frobenius gives $p$-semilinear
maps $F_i:V_i\to V_{pi}$. For an orbit $O$ of multiplication
by $p$ on $(\mathbf Z/n\mathbf Z)^\times$, choose $i\in O$
and let $r_O$ be the eventual rank of the iterates of
$F^{|O|}:V_i\to V_i$. Then
\[
                  f_p(P_n(D))=\sum_O |O|r_O.
\]
In particular the primitive part has $p$-rank zero if and only
if every such orbit composite is nilpotent. Here
$\dim_\Omega V_i=g-1$ for every nontrivial $i$.
\end{claim}
\begin{proof}
The finite etale algebra of the cyclic cover gives
\[
 H^1(D,\mathcal O_D)=\bigoplus_{i\in\mathbf Z/n\mathbf Z}V_i,
\]
and taking $p$th powers sends index $i$ to $pi$. Changing the
chosen trivialization only changes the displayed identifications
and does not change stable ranks. Riemann--Roch and
$H^0(C,L^i)=0$ give $\dim V_i=g-1$ for $i\ne0$.
For a finite-dimensional space with a Frobenius-semilinear
endomorphism, the stable image is the maximal subspace on which
Frobenius is bijective. Over $\Omega$, its dimension equals the
$\mathbf F_p$-dimension of the fixed space; this follows by
solving the semilinear equation $F(v)=v$ on the bijective part.
The Artin--Schreier sequence identifies that fixed space with
$H^1_{\mathrm{et}}(D,\mathbf F_p)$, whose dimension is $f_p(D)$.
On a character orbit the stable image has the same dimension
$r_O$ at each index, giving the factor $|O|$.

Since $n$ is invertible in characteristic $p$, the character
projectors split the cohomology. The endomorphism
$\prod_{\ell\mid n}(1-\sigma^{n/\ell})$ is nonzero precisely
on primitive characters. Its image therefore has exactly the
primitive stable ranks in the displayed formula. A descending
sequence of images of a semilinear endomorphism stabilizes in
finite dimension; its stable image is zero exactly when some
iterate is zero. This proves the criterion.

In genus two and characteristic two, the locally exact
differential sequence identifies $\ker F_i$ with
$H^0(C,\mathcal B\otimes L^i)$. Each $V_i$ is a line, so an
orbit composite is nilpotent precisely when one arrow is zero.
This is the theta-orbit condition in R5 and R8. For $g\ge3$,
noninvertibility of a single arrow is weaker than nilpotence
of the orbit composite.
\end{proof}

\begin{claim}{AG.4}{Refuted; nonordinariness alone is not an annihilation argument}
The implication from a nonordinary character-cohomology block
to absence of surviving $p$-primary vectors is false. The failure
persists for a nonabelian group and a vector with free deck orbit
in the abstract module.
\end{claim}
\begin{proof}
Already the semilinear map
$(x,y)\mapsto(x^p,0)$ on $\overline{\mathbf F}_p^2$ has rank
one and stable rank one. It is noninvertible but is not
nilpotent. For a group example take
$G=\mathbf F_8^\times\rtimes\operatorname{Gal}(\mathbf F_8/\mathbf F_2)$
acting diagonally on $W=\mathbf F_8^2$. The vector $(1,\alpha)$,
where $\alpha^3+\alpha+1=0$, has trivial stabilizer: a group
element $(u,b)$ fixing it has $u=1$ from the first coordinate,
then $\alpha^{2^b}=\alpha$, whence $b=0$.
On $W\oplus\mathbf F_8$ take the Frobenius-semilinear map
which is the identity on the $\mathbf F_2$-structure of $W$
and zero on the last summand. It is $G$-equivariant, has
dimension nine, stable rank six, and retains the fixed vector
with free $G$-orbit. The program
\texttt{code/geometric\_checks.py} verifies the group law,
all 9261 associativity triples, the free orbit, and the ranks.
This is a refutation of the proposed linear weakening, not a
curve with a collision. The correct condition is stable-rank
zero on the relevant summand, as in AG.3.
\end{proof}

\begin{claim}{AG.5}{Proved; a small integral group-ring obstruction}
Let $\Omega$ be algebraically closed of characteristic $p$, and
let $\pi:D\to C$ be a connected finite etale Galois cover of
smooth proper curves, with $g(C)\ge2$ and group $G$ of order
$m$ prime to $p$. Let
\[
 a=\sum_{g\in G}a_g g\in\mathbf Z[G]\setminus\{0\},
 \qquad d=\sum_g|a_g|.
\]
Assume $(d,m)=1$ and $m>(d-1)^2$. If $G$ is nonabelian,
then for any points $P',Q'\in D(\Omega)$ lying above distinct
points of $C$ one has
\[
                         a[Q'-P']\ne0\quad\hbox{in }J_D(\Omega).
\]
\end{claim}
\begin{proof}
Suppose the class is zero. Since the two free $G$-orbits are
disjoint, the divisor $\sum a_g(gQ'-gP')$ is a nonzero
principal divisor with positive and negative degrees both
equal to $d$. Let $f\in\Omega(D)$ have that divisor, and put
\[
 M=\Omega(\gamma f:\gamma\in G),\qquad
 e=[\Omega(D):M].
\]
Let $Z$ be the smooth proper curve with function field $M$.
Every generator has degree $d$. Poonen's Lemma 2.2 and
Corollary 2.3, with Proposition 2.1
\cite{AG_Poonen2007}, give
\[
                 e\mid d,\qquad g(Z)\le(d/e-1)^2.
\]
For clarity, the induction proving the bound is as follows.
If the current field has index $d/s$ and genus at most
$(s-1)^2$, adjoining another degree-$d$ function with relative
degree $b$ gives index $d/(bs)$ and genus at most
\[
 b(s-1)^2+(b-1)(bs-1)
   =(bs-1)^2-bs(b-1)(s-1)\le(bs-1)^2.
\]
This is Castelnuovo--Severi for the two generating subfields.
The perfectness hypothesis holds because $\Omega$ is
algebraically closed. The result allows arbitrary rational
functions, so no unproved separability assertion is needed.

Since $g(D)=m(g(C)-1)+1\ge m+1>(d-1)^2$, we have $e>1$.
The field $M$ is $G$-stable. The kernel of the action on $M$
has order dividing both $e$ and $m$, so the action on $Z$ is
faithful. It is free: if a point stabilizer contained a subgroup
of prime order $\ell$, that subgroup would act freely on the
degree-$e$ fibre divisor in $D$, partitioning its points into
orbits of size $\ell$ with constant multiplicity. This would
force $\ell\mid e$, contradicting $(e,m)=1$.

Consequently $Z\to Z/G$ is etale and
$g(Z)-1=m(g(Z/G)-1)$. The numerical bound excludes
$g(Z)\ge m+1$ and genus zero is impossible in that equality;
hence $g(Z)=1$. Choose an origin on $Z$. An automorphism is
$x\mapsto t+u(x)$ with $u$ an origin-preserving automorphism.
If $u\ne1$, the nonzero endomorphism $1-u$ of the elliptic
curve is surjective, so this automorphism has a fixed point.
Every element of a free action is therefore a translation.
Translations commute, contradicting that $G$ is nonabelian.
\end{proof}

\begin{claim}{AG.6}{Proved; rational and geometric nonabelian criteria}
In AG.5 suppose all cover data are defined over a finite field
$L$, $G$ is constant over $L$, and $P'\in D(L)$ lies above
$P\in C(L)$. Put $B_a=\operatorname{im}(a:J_D\to J_D)$,
an abelian subvariety over $L$. If $B_a(L)[p]=0$, then $P$
collides with no other point of $C(L)$. The same conclusion
holds over all finite fields of definition if $f_p(B_a)=0$
geometrically. Thus an odd nonabelian group can use, for
example, $a=(1-\delta)(1-\epsilon)$ whenever the four support
elements are distinct and $m>9$.
\end{claim}
\begin{proof}
A collision lifts by R1 to a rational nonzero $p$-primary
class $T=[Q'-P']$. Its image $aT$ is a rational $p$-primary
point of $B_a$, hence zero under the rational hypothesis.
This contradicts AG.5. In the geometric version one first
defines the cover, its automorphisms, and the lift over a
finite extension; $f_p(B_a)=0$ ensures that the torsion
vanishing persists over that extension. For the displayed
four-term element, $d=4$ and an odd group has $(d,m)=1$.

More generally, for any collection of integral endomorphisms
$a_j$, a joint map $D\to\prod_j B_{a_j}$ separating the two
fibres gives a collision obstruction if every target has no
rational $p$-torsion over the field where the lifts are rational.
AG.5 is a concrete condition that supplies this separation
for one endomorphism. It does not supply existence of a
cover with the required torsion vanishing.
\end{proof}

\begin{claim}{AG.7}{Proved; the Conway--McKay--Trojan test is outside the category}
The polynomial covers $t=x^{24}+x$ in characteristic two and
$t=x^{12}+x$ in characteristic three from
Conway--McKay--Trojan \cite{AG_CMT2008} are finite etale over
the affine line but are wildly ramified on their proper
compactifications. Their root-space representations do not
give instances of the prime-to-$p$ etale criteria above.
\end{claim}
\begin{proof}
Both affine derivatives are one. At infinity the ramification
indices are respectively 24 and 12, divisible by the
characteristic. A Galois closure therefore has nontrivial
wild inertia and its group has order divisible by $p$.
The source and base of the non-Galois polynomial maps are
projective lines, outside the genus hypothesis as well.

For the characteristic-two polynomial put
\[
 L(x)=x^{2048}+t^{64}x^{512}+t^8x^{16}+t^{16}x^8+x.
\]
Reduction modulo $x^{24}+x+t$ gives $L(x)=t+t^{24}$.
Consequently the roots satisfy the additive polynomial
$L(x)^2+(t+t^{24})L(x)$. It is separable of degree $2^{12}$,
because its coefficient of $x$ is the nonzero rational
function $t+t^{24}$. Its root space has dimension 12 over
$\mathbf F_2$, and the Galois action on the subspace spanned
by the original roots is faithful. The script
\texttt{code/geometric\_checks.py} verifies the polynomial
remainders and all additive exponents exactly. None of these
linear spaces is identified with the $p$-torsion of a
Jacobian. In particular faithfulness of this representation
cannot substitute for the geometric point-difference
condition in AG.5.
\end{proof}

\begin{claim}{AG.8}{Known, checked; the orbifold theorem covers the required class}
The pro-$\mathscr S$ finite-subgroup theorem used in AP-A.29l
applies to $\mathscr S$ equal to all finite groups of order
prime to $p$. Thus this particular cited step in R4 has the
required scope.
\end{claim}
\begin{proof}
Boggi--Zalesskii Definition 1.1 defines a class by closure under
subgroups, quotients, and extensions. Their Lemma 3.1 assumes
a hyperbolic Fuchsian group $\Lambda$ with a torsion-free
normal subgroup $H$ and $\Lambda/H\in\mathscr S$. It gives
the malnormality of the nontrivial cyclic cone groups and
containment of every finite nontrivial subgroup of the
pro-$\mathscr S$ completion in a conjugate cone group
\cite{AG_BoggiZalesskii2019}. All three closure conditions
hold for prime-to-$p$ groups.

For a prime-to-$p$ Galois map $C\to B$ with $g(C)\ge2$,
the orbifold presentation has cone orders equal to the
ramification indices. Its surjection to the deck group is
injective on each cone group, so its kernel is the
torsion-free surface group of the topological covering.
Its orbifold Euler characteristic is
$(2-2g(C))/|G|<0$. These are precisely the lemma's
hypotheses. We have read the proof, which uses the
$\mathscr S$-solenoid and acyclicity at primes occurring
in $\mathscr S$; no prime-$p$ acyclicity is needed.
This verifies the completion comparison, not every other
step of R4 by mere citation.
\end{proof}

\begin{claim}{AG.9}{Open for supersingular correspondences; nontrivial primitive-cover existence is refuted}
The universal supersingular-correspondence assertion of A5
is not established here. The cited geometric existence
theorems do not imply the universal primitive-zero cover
assertion A6; the later explicit argument AG.32--AG.33
refutes its degree-$n>1$ version in genus two and
characteristic two. The literal version permitting the
identity cover remains open for curves of positive $p$-rank.
\end{claim}
\begin{proof}[Exact obstruction retained]
The nilpotence requirement is AG.3. Raynaud's Theorem 2
provides a solvable prime-to-$p$ cover with a nonordinary
representation \cite{AG_Raynaud2000}; noninvertibility
does not imply nilpotence, as AG.11 proves on an actual cover. The theta torsion
density of Yang Proposition 2.3 and Theorem 2.4
\cite{AG_Yang2016}, and the related Pop--Saidi propositions
\cite{AG_PopSaidi2003}, do not prescribe a torsion order for
which every primitive multiplication-$p$ orbit is met.
In genus two, characteristic two, a single orbit can be
killed while other primitive orbits survive. In higher
genus even one orbit can retain positive stable rank after
some Frobenius arrows become singular.

Ozman--Pries Theorem 1.1 concerns generic curves and ordinary
Pryms; Example 9.4 excludes degree-three witnesses for
nonordinary genus-two curves in characteristic two
\cite{AG_OzmanPries2019}. Generic counterexamples over a
field of positive transcendence degree do not give a
counterexample over $\overline{\mathbf F}_p$: the latter
has only countably many moduli points, and a countable
union of proper closed loci can contain all of them.
No countability argument is used in the reverse direction.

Bogomolov--Tschinkel Theorem 1.7 assumes a hyperelliptic
source and $p\ge5$; their Proposition 1.8 assumes
characteristic different from two and three
\cite{AG_BogomolovTschinkel2002}. Poonen Theorem 1.7 has
a prime-to-$p$ Galois action with quotient genus at most
two, together with one of its additional conditions
\cite{AG_Poonen2005}. These hypotheses exclude a general
curve, and an arbitrary etale cover in those statements
must additionally be checked for prime-to-$p$ Galois
closure before applying the collision lifting argument.
The searches and subsequent attempts are recorded in the
source audit; these citations are not substitutes for the
missing existence proof.
\end{proof}

\begin{claim}{AG.10}{Proved; explicit cyclic zero-rank families in characteristics two and three}
Let $p=2$ or $3$, and let $n>1$ be odd and prime to $p$, with
$(p-1)(n-1)\ge4$. Suppose that
\[
                    (p-1)n\mid p^r+1
\]
for some positive integer $r$. The smooth proper curve
\[
                         C_{p,n}:y^p-y=x^n
\]
has a connected cyclic etale cover of degree $n$ whose entire
Jacobian, and therefore its primitive conductor-$n$ factor,
has $p$-rank zero. For $p=2$ the cover is the Fermat curve
$u^n+z^n=1$; for $p=3$ it is the smooth model of
$u^n=z^{2n}-1$. In both cases the map is
\[
                         x=uz,\qquad y=z^n.
\]
Thus these are actual cases of A6 in both small characteristics,
without a separability or Galois-closure qualification left over.
\end{claim}
\begin{proof}
The equation for $C_{p,n}$ has one point at infinity. At that
point $x$ has pole order $p$ and $y$ pole order $n$:
the valuation equation gives $p\,v(y)=n\,v(x)$, and the
Artin--Schreier extension has degree at most $p$, forcing
degree and ramification index $p$. Its different exponent is
$(p-1)(n+1)$, hence its genus is $(p-1)(n-1)/2$.
At $P=(0,0)$, the variable $x$ is a parameter and $y$ has
order $n$. Consequently
\[
                          \operatorname{div}(y)=n(P-\infty).
\]
The class $[P-\infty]$ has exact order $n$. Indeed, if its
order were a proper divisor $d<n$ of $n$, a function with
only pole $d\infty$ would exist. Every function regular away
from infinity is a sum of monomials $x^a y^b$ with $a\ge0$
and $0\le b<p$. Their pole orders $pa+nb$ are pairwise
distinct when grouped with fixed $b$ and then compared
modulo $p$. A pole order less than $n$ must therefore be
divisible by $p$, whereas $d\mid n$ is prime to $p$.

Adjoining $z$ with $z^n=y$ thus gives a connected cyclic
degree-$n$ extension. It is unramified because every
valuation of $y$ is divisible by $n$ and $n$ is invertible;
locally a unit has an $n$th root in the strict henselization.
Putting $u=x/z$ gives
\[
                           u^n=z^{(p-1)n}-1.
\]
The constant deck group $\mu_n$ acts by
$(u,z)\mapsto(\zeta^{-1}u,\zeta z)$ after a finite field
extension containing $\mu_n$. This is a cyclic geometric
Galois cover of order prime to $p$, so its own Galois closure
already has the required order.

We give the zero-rank argument to avoid relying on an
unstated supersingularity criterion for Fermat curves. Put
$Q=p^r$. The degree-$Q+1$ Fermat curve is isomorphic over
the algebraic closure to
\[
                        w^Q-w=t^{Q+1}.
\]
To see this, choose $\beta\in\mathbf F_{Q^2}\setminus\mathbf F_Q$
and $\gamma\in\mathbf F_{Q^2}$ with $\gamma^{Q+1}=-1$.
The substitutions
$U=W+\beta Z$ and $V=\gamma(W+\beta^Q Z)$ give
\[
 U^{Q+1}+V^{Q+1}
       =(\beta-\beta^Q)(W^QZ-WZ^Q).
\]
Rescaling the remaining coordinate gives the displayed
Artin--Schreier form. Its projection to the $t$-line has
additive group $\mathbf F_Q$ and is etale away from infinity.
The valuation relation at infinity, with
$\gcd(Q,Q+1)=1$, proves irreducibility, degree $Q$, and
total ramification there. Crew's Corollary 1.8
\cite{AG_Crew1984}, for a finite $p$-group acting on a
smooth proper curve over an algebraically closed field,
states the Deuring--Shafarevich formula
\[
 f_p(H)-1=|G|\bigl(f_p(H/G)-1\bigr)
                       +\sum_{R\in H}(e_R-1).
\]
Here it gives $f_p(H)-1=-Q+(Q-1)=-1$.

Set $N=(p-1)n$. Since $N\mid Q+1$, the coordinate-power
map from the degree-$Q+1$ Fermat curve to the degree-$N$
Fermat curve is finite separable of degree
$((Q+1)/N)^2$, prime to $p$. The latter Fermat curve maps
to our cover by
\[
                     (w,z)\longmapsto (w^{p-1},z)
\]
in the equation $w^N=z^N-1$. Its degree is $p-1$, also
prime to $p$. Pullback followed by norm is multiplication
by the degree; on geometric $p$-torsion this is invertible.
Zero $p$-rank therefore descends through both maps.
Every abelian subvariety of this zero-rank Jacobian has
zero $p$-rank, including the primitive factor.

The restrictions describe infinite families: one may take
$n=2^r+1$ for $p=2$, $r\ge2$, and
$n=(3^{2s}+1)/2$ for $p=3$, $s\ge1$. They do not cover
arbitrary curves of those genera.
\end{proof}

\begin{claim}{AG.11}{Refuted on an actual curve cover: nonordinary primitive blocks need not be nilpotent}
In characteristic two, let
\[
 C:y^2+y=x^7,\qquad
 D:u^7+z^7=1,\qquad
 (x,y)=(uz,z^7).
\]
Then $D\to C$ is connected cyclic etale of degree seven,
$g(C)=3$, $g(D)=15$, and $f_2(C)=0$, $f_2(D)=6$.
Each nontrivial deck-character space has coherent
dimension two and stable Frobenius dimension one. In
particular both primitive doubling-orbit composites are
noninvertible, but neither is nilpotent. The primitive
factor has dimension twelve and two-rank six.
\end{claim}
\begin{proof}
The construction and connectedness were proved in AG.10
before its congruence hypothesis was used. Alternatively,
the action
$(u:z:w)\mapsto(\zeta^{-1}u:\zeta z:w)$ on the Fermat
curve is free: in the affine chart a nonidentity element
cannot fix both nonzero coordinates, and at infinity
both $u,z$ are nonzero and fixedness requires
$\zeta^{-1}=\zeta$, impossible for nontrivial seventh
roots. Its invariants $x=uz$, $y=z^7$ give precisely the
displayed quotient. The genus formulas follow either
from these equations or from etale Riemann--Hurwitz.

We compute the Cartier operator directly. A basis of
$H^0(D,\Omega_D^1)$ is
\[
 \omega_{a,b}=u^{a-1}z^{b-7}\,du,\qquad
      a,b\ge1,\quad a+b<7.
\]
This is the usual adjunction basis
$u^i z^j\,du/z^6$, $i,j\ge0$, $i+j\le4$;
there are fifteen forms. For the separating variable $u$,
Cartier sends $(A^2+B^2u)\,du$ to $B\,du$.
Using $z=(1+u^7)/z^6$ gives
\[
 \mathcal C(\omega_{a,b})=
 \begin{cases}
 0,&a,b\text{ both odd},\\
 \omega_{a',b'},&\text{otherwise},
 \end{cases}
 \qquad
 a'\equiv4a,\quad b'\equiv4b\pmod7,
\]
where $a',b'$ are in $\{1,\ldots,6\}$. In the second
case their sum is less than seven. For example, if
$a,b$ are both even the target is $(a/2,b/2)$;
if $a$ is odd and $b$ even it is $((a+7)/2,b/2)$;
the remaining case is $(a/2,(b+7)/2)$.

The stable basis consists exactly of
\[
             (1,2),(2,4),(4,1),\qquad
             (2,1),(4,2),(1,4),
\]
which form two Cartier cycles. All other basis vectors
eventually map to zero. The deck character of
$\omega_{a,b}$ is $b-a$ modulo seven. The invariant
space has dimension three and stable rank zero; every
nonzero character has dimension two and stable rank one.
Cartier and coherent Frobenius are dual, with opposite
character conventions, so they have the same orbitwise
stable ranks. The Cartier fixed-space description of
geometric two-torsion, equivalently the Artin--Schreier
description in AG.3, now gives all stated two-ranks.

The exact graph of all fifteen basis vectors and both
multiplicity tables are generated and checked by
\path{code/fermat_cyclic_checks.py}, with its output
attached. This counterexample concerns the proposed
weakening in L8. It does not assert that $C$ has no other
cyclic cover satisfying A6.
\end{proof}

\begin{claim}{AG.12}{Proved; a primitive-only zero-rank example}
In characteristic two, the cyclic etale degree-$35$ cover
\[
 u^{35}+z^{35}=1\ \longrightarrow\
       C:y^2+y=x^{35},\qquad (x,y)=(uz,z^{35}),
\]
has $g(C)=17$, $g(D)=561$ and $f_2(D)=6$, while its
primitive conductor-$35$ factor has dimension $384$ and
two-rank zero.
\end{claim}
\begin{proof}
The connected etale construction is again the one proved in
AG.10. The Cartier basis and its calculation in AG.11 apply
with every occurrence of seven replaced by any odd $n$.
Explicitly, for $a,b\ge1$ and $a+b<n$ the operator is zero
if both indices are odd, and otherwise replaces them by
their least positive halves modulo $n$. The deck weight
is $b-a$ modulo $n$. This also proves the finite formula
\[
 f_2(P_n)=\#\left\{(a,b):
 \begin{array}{l}
 1\le a,b<n,\quad \gcd(b-a,n)=1,\\
 {}[2^j a]_n+[2^j b]_n<n
       \text{ for every }0\le j<\operatorname{ord}_n(2)
 \end{array}\right\},
\]
where $[\,\cdot\,]_n$ is the representative in
$\{1,\ldots,n-1\}$. A Cartier basis vector belongs to the
stable image exactly when its entire finite halving orbit
stays in the holomorphic basis.

For $n=35$, the complete graph of the $561$ basis vectors
has stable set
\[
 (5,10),(10,20),(20,5),\qquad
 (10,5),(20,10),(5,20).
\]
The attached program verifies every transition and iterates
the entire image set until it stabilizes; its output includes
the full graph, so the finite assertion is reproducible.
Every surviving weight is divisible by five and is therefore
nonprimitive modulo $35$. AG.3 gives primitive two-rank
zero. Its dimension is
$\varphi(35)(g(C)-1)=24\cdot16=384$.
This example also distinguishes primitive vanishing from
vanishing of the complete relative factor. It does not reach
a pair beyond A3: all surviving characters have conductor
seven, so the intermediate cyclic degree-five cover has
zero full new two-rank. After defining this cover and a
chosen lift over a finite field, it is itself an L1 witness.
\end{proof}

\begin{claim}{AG.13}{Proved; the elliptic rational lever in characteristics two and three}
Let $E/L$ be an elliptic curve over a finite field of
characteristic two. Then $E(L)[2]=0$ if and only if $E$ is
supersingular. Thus an ordinary elliptic target does not
weaken the geometric zero-rank requirement in L5 when $p=2$.
In characteristic three, if $E$ is ordinary and $E(L)[3]=0$,
this vanishing persists over every odd-degree extension of
$L$ and fails over every even-degree extension.
\end{claim}
\begin{proof}
For an ordinary elliptic curve in characteristic $p$, the
geometric group of $p$-torsion points is cyclic of order $p$.
For $p=2$ its unique nonzero point is fixed by every
automorphism of the algebraic closure over $L$, so it is
rational. A supersingular elliptic curve has no nonzero
geometric $p$-torsion points. For $p=3$, Frobenius on the
one-dimensional $\mathbf F_3$-space is either $1$ or $-1$.
The stated rational vanishing selects $-1$, whose $d$th
power is $1$ exactly for even $d$.
\end{proof}

\begin{claim}{AG.14}{Proved; a genus-lowering closure step in characteristic two}
Let $C/\overline{\mathbf F}_2$ be a hyperelliptic curve of
genus $g\ge3$. There is a connected cyclic etale cover
$D\to C$ of degree three and a smooth proper curve $B$
of genus $g-1$ with two morphisms $D\to B$ which jointly
separate geometric points. Consequently, if every genus-$(g-1)$
curve has no collision over its finite fields of definition,
then $C$ has no collision. In particular universal A1 in
genus two and characteristic two would imply A1 for every
hyperelliptic genus-three curve in that characteristic.
\end{claim}
\begin{proof}
Choose a nontrivial order-three character of the geometric
abelianized fundamental group of $C$ and let $D\to C$ be
its connected cyclic cover, with generator $\sigma$.
Choose a hyperelliptic ramification point $W$ and a lift
$W'$ to $D$. The hyperelliptic involution $\iota$ acts by
$-1$ on $J_C$, so it preserves the character kernel and
acts by inversion on the cyclic quotient. It lifts to an
automorphism $\tau$ fixing $W'$. Its square is a deck
transformation fixing $W'$, hence is the identity.
Thus
\[
            \tau^2=1,\qquad \tau\sigma\tau=\sigma^{-1}.
\]
Let $B=D/\langle\tau\rangle$. On the three-point fibre
above each hyperelliptic ramification point, $\tau$ is an
affine inversion and has exactly one fixed point. It has
no fixed points above unramified points of $C\to\mathbf P^1$.
Because $D\to C$ is etale, the completed local rings at
these fixed points, together with the involution, identify
with the completed local rings downstairs. Hence the
different for $D\to B$ has the same total degree as the
different for $C\to\mathbf P^1$, namely $2g+2$.
Etale Riemann--Hurwitz gives $g(D)=3g-2$, and the
separable degree-two quotient formula then gives
\[
 6g-6=2(2g(B)-2)+(2g+2),\qquad g(B)=g-1.
\]
This computation retains the wild different in characteristic
two; tame ramification counting would be incorrect here.

Let $\phi:D\to B$ be the quotient map. The maps $\phi$
and $\phi\circ\sigma$ jointly separate points. Indeed, if
$R\ne S$ have the same image under both, then
$S=\tau R=\sigma^{-1}\tau\sigma R$. The product of
these two distinct involutions is a nonidentity power of
$\sigma$, which would fix $R$, contradicting the etaleness
of the cyclic cover. Every object and map descends after
a finite extension of a finite field. Apply AG.2 to the
two maps and then to $D\to C$.

For $g=2$ the same construction has an elliptic target.
It only supplies the desired vanishing when that elliptic
target has two-rank zero, as AG.13 explains; the reduction
does not by itself settle the genus-two base case.
\end{proof}

\begin{claim}{AG.15}{Proved; exact small-prime theta tests in a supersingular genus-two family}
Over $\overline{\mathbf F}_2$ let
\[
                      C_a:y^2+y=x^5+a x^3.
\]
Every genus-two curve of two-rank zero is geometrically
isomorphic to some $C_a$.
Its Raynaud theta divisor, based at its point $\infty$,
is $\Theta_a=\{[R-\infty]:R\in C_a\}$.
It has no nonzero point of order three. It has a point
of exact order five if and only if $a=0$, and a point
of exact order seven if and only if $a^5=1$.
Consequently:
\begin{enumerate}
\item for $a=0$ a degree-five cyclic etale cover has
      zero primitive two-rank;
\item for $a^5=1$ a degree-seven cyclic etale cover has
      zero primitive two-rank, while every degree-three
      or degree-five cyclic cover is new-ordinary;
\item if $a\ne0$ and $a^5\ne1$, every geometric cyclic
      cover of degree three, five or seven is new-ordinary.
\end{enumerate}
These are geometric assertions. The last assertion alone
does not exclude rational L1 witnesses of those degrees
over a particular field.
\end{claim}
\begin{proof}
For the normal-form assertion, the genus-two hyperelliptic
map is a separable quadratic map to $\mathbf P^1$, hence
has an Artin--Schreier equation. Put its rational function
in reduced form, with odd pole orders. The
Deuring--Shafarevich formula makes the two-rank one less
than the number of poles. Thus there is one pole, which
Riemann--Hurwitz forces to have order five. Move it to
infinity, scale $x$, and remove the constant and even
powers by changing $y$; the equation becomes
$y^2+y=x^5+a x^3+b x$. Translation $x=u+r$ followed
by changes of $y$ by a polynomial of degree two removes
the coefficient of $u$ whenever
\[
                   r^{16}+a^4r^8+a^2r^2+r=b^4.
\]
This monic polynomial has a root over the algebraic
closure. The coefficient of $u^3$ remains $a$,
proving the assertion.

The Artin--Schreier map has a single pole, of order five,
so the genus is two and the two-rank is zero. The different
at infinity has order six; hence $\operatorname{div}(dx)=2\infty$.
Since $dx$ is locally exact, its half-divisor identifies
the Raynaud bundle with $\mathcal O(\infty)$, giving
the asserted theta origin.

Every function regular away from infinity is
$A(x)+B(x)y$, and $x,y$ have pole orders two and five.
If $[R-\infty]$ has exact odd order $\ell$, it is nonzero,
$R$ is affine, and a function $F$ has
$\operatorname{div}(F)=\ell R-\ell\infty$.
There is no function with unique pole order three, proving
the first assertion.

For $\ell=5$, normalize $F=A+y$ with $\deg A\le2$.
Its hyperelliptic norm is
\[
             A^2+A+x^5+a x^3=(x+r)^5,\qquad r=x(R).
\]
The coefficient of $x^3$ forces $a=0$. Conversely,
when $a=0$, the divisor of $y$ is
$5(0,0)-5\infty$, so that point gives exact order five.

For $\ell=7$, write
\[
 F=A+(x+d)y,\qquad
 A=c_3x^3+c_2x^2+c_1x+c_0.
\]
The leading coefficient of $x y$ is nonzero because the
pole order is seven, so this normalization loses no
possible function. The norm condition is
\[
 A^2+(x+d)A+(x+d)^2(x^5+a x^3)=(x+r)^7.
\]
Write $a=t^2$ and $r=s^2$, uniquely in the algebraic
closure. Comparing coefficients from degree six down to
degree three gives
\[
 c_3=s,\quad d=s^2+t,\quad
 c_2=s^8+t^2s^4+s^3+st+t^4,
\]
and
\[
             E_1=s^{16}+t^4s^8+t^2s^2+s+t^8=0.
\]
The coefficients of degrees two and one give
$c_1^2+c_1=r^5+c_2d$ and $c_0=r^6+c_1d$.
Substituting these in the constant coefficient gives
\[
             E_2=s^{24}+t^2s^{10}+s^9+t^6s^2+t^4s+t^7=0.
\]
The complete elimination is the polynomial identity
\[
                     E_2+(s^8+t^4)E_1=t^7(t^5+1).
\]
It is verified symbolically in
\path{code/supersingular_small_primes.py}; the output also
checks the resultant
$\operatorname{Res}_s(E_1,E_2)=t^{112}(t^5+1)^{16}$.

If $t\ne0$, these equations have a solution precisely
when $t^5=1$: necessity follows from the displayed
identity, and sufficiency follows by taking any root of
the monic degree-sixteen polynomial $E_1$ and then solving
the Artin--Schreier equation for $c_1$. Moreover
$(x+d)|_{x=r}=t\ne0$. The norm has its only finite zero
at $r$, but $F$ cannot vanish at both of the two points
over $r$, because its two values differ by $r+d=t$.
Its divisor is therefore exactly $7R-7\infty$ for one
of those points. This proves existence of a point of
order seven when $a^5=1$.

If $t=0$, the coefficient equations instead give $d=r$.
The norm condition then forces $A(r)=0$, so $F$ vanishes
at both hyperelliptic conjugate points above $r$. They
are distinct, since the affine Artin--Schreier map is
etale. Such an $F$ cannot have divisor $7R-7\infty$.
This excludes the extraneous zero of the resultant at $a=0$.

Finally theta is symmetric: the hyperelliptic involution
sends $[R-\infty]$ to its negative. For degree five,
two generates $(\mathbf Z/5)^\times$, so one theta point
kills the sole primitive doubling orbit. For degree
seven its two doubling orbits are exchanged by negation.
A theta point and its negative therefore meet both
orbits. AG.3 proves zero primitive rank in the asserted
cases. Absence of theta at one of these prime orders
makes every relevant Frobenius arrow invertible, proving
the stated new-ordinariness assertions.
\end{proof}

\begin{claim}{AG.16}{Proved; the degree-eleven theta locus in the same family}
For $C_a$ in AG.15, the Raynaud theta divisor contains
a point of exact order eleven if and only if
\[
                           a^{35}+a^{10}+1=0.
\]
At each such parameter there is a connected cyclic
etale degree-eleven cover with zero full new two-rank.
At every other parameter all degree-eleven cyclic
covers are new-ordinary.
\end{claim}
\begin{proof}
At an affine point, translate its $x$-coordinate to
$u=0$ and its $y$-coordinate to zero. Removing the
even powers by adding a quadratic polynomial to $y$
gives the local and global affine equation
\[
                        Y^2+Y=u^5+a u^3+c u.
\]
The pole orders of $u,Y$ at infinity are still two
and five. Every $c$ can occur for fixed $a$: if the
translated point has original coordinate $r$, the
coefficient calculation is
$c^4=r^{16}+a^4r^8+a^2r^2+r$, a monic degree-sixteen
polynomial in $r$. Its value can be any prescribed
element over the algebraic closure.

On the branch with $Y(0)=0$, the formal expansion to
the necessary order is
\[
 Y=cu+c^2u^2+a u^3+c^4u^4+u^5
                       +a^2u^6+c^8u^8+u^{10}+O(u^{11}).
\]
A function with sole pole order eleven has the form
$A(u)+B(u)Y$, with $\deg A\le5$ and
$B=u^3+b_2u^2+b_1u+b_0$. Choose $A$ to cancel terms
of degree at most five. Writing $v=c^4$, the necessary
and sufficient conditions for vanishing to order at
least eleven are the five coefficients
\[
 \begin{aligned}
 b_0a^2+b_1+b_2v+a&=0,&
 b_1a^2+b_2+v&=0,\\
 b_0v^2+b_2a^2+1&=0,&
 b_1v^2+a^2&=0,&
 b_0+b_2v^2&=0.
 \end{aligned}
\]
If $v=0$, the fourth equation gives $a=0$ and the
third becomes $1=0$. Thus $v\ne0$. Eliminate
$b_1=a^2/v^2$, $b_2=a^4/v^2+v$ and
$b_0=a^4+v^3$. The remaining equations are
\[
 \begin{aligned}
 E_1&=v^7+a^4v^4+a^2v^3+v^2+a^6=0,\\
 E_2&=a^2v^5+v^4+(a^6+a)v^2+a^4v+a^2=0.
 \end{aligned}
\]
The exact resultant is
\[
              \operatorname{Res}_v(E_1,E_2)
                        =a^4(a^{35}+a^{10}+1).
\]
The program \path{code/supersingular_eleven.py}
verifies all five substitutions and computes this
resultant over $\mathbf F_2[a]$, with output attached.
When $a=0$, $E_2=v^4$ excludes every nonzero $v$.
When $a\ne0$, both displayed leading degrees are
unchanged, so resultant zero is equivalent to a common
root. The constant coefficient $a^6$ of $E_1$
ensures that any such root is nonzero.

For each solution, the constructed function has a
pole of exactly order eleven, and a zero of order
at least eleven at the selected point. Thus its
divisor is precisely eleven times the difference
between that point and infinity. The point is
distinct from infinity, so its theta class has
exact prime order eleven. Conversely any such
theta point supplies the coefficient equations
above. This proves the equivalence, including
the extraneous resultant factor at $a=0$.
Since two has order ten modulo eleven, AG.3 gives
the final two-rank conclusions.
\end{proof}

\begin{claim}{AG.17}{Proved; even deck groups in odd characteristic}
Let $\Omega$ be algebraically closed of odd characteristic,
and let $D\to C$ be a connected etale Galois cover of
smooth proper curves with $g(C)\ge2$ and geometric group
$G$ of order $m$ prime to the characteristic. Suppose
$m\nmid4$. For points $P',Q'\in D(\Omega)$ in distinct
fibres, the class $[Q'-P']$ has trivial stabilizer in $G$.
In particular this holds for every nonabelian $G$,
including groups of even order.

Consequently, over a finite field of odd characteristic,
the full relative class-number criterion L1 holds for a
pointed geometrically Galois etale cover with such a group.
The deck group need not be constant: the geometric average
is the rational pullback--norm projector. This supplies
an even-group version of the nonabelian criterion, without
asserting that its existence reaches a pair outside A3.
\end{claim}
\begin{proof}
If $1\ne\delta\in G$ fixes $[Q'-P']$, then
\[
           \delta Q'+P'-\delta P'-Q'
\]
is principal. Its four points are distinct, since the
deck action is free and the two original points lie in
distinct fibres. Thus its function has degree two.
Since $g(D)\ge2$, this is the hyperelliptic map. It is
separable because the characteristic is odd.

The hyperelliptic pencil is unique, so every automorphism
of $D$ preserves its ramification-point set $W$.
There are exactly $2g(D)+2$ such geometric points in odd
characteristic. The action of $G$ on $W$ is free, because
$D\to C$ is etale. Hence
\[
 m\mid\#W=2g(D)+2=2m\bigl(g(C)-1\bigr)+4,
\]
contradicting $m\nmid4$. This proves the stabilizer assertion.

For the rational criterion, a collision lifts to a rational
$p$-primary point-difference class $T$ on $D$ by L0,
or by its pointed nonconstant-deck form AC.18. The
prime-to-$p$ isogeny projector
$m^{-1}\pi^*\pi_*$ splits its invariant and new parts.
If $p$ does not divide $h_L(D)/h_L(C)$, the rational new
part is zero. Thus $T$ is fixed by every geometric deck
element, contrary to the stabilizer assertion. The
pullback--norm operator is defined over $L$, so the same
argument covers nonconstant geometric deck groups.

The divisibility restriction in the geometric argument is
real: a hyperelliptic curve can carry a free group of
order two or four. No claim about arbitrary point-difference
stabilizers in those degrees is being inferred. For odd
groups the Feit--Thompson reduction AC.20 applies; an even
group need not have an odd-order abelian quotient, so that
specific reduction does not extend automatically.
\end{proof}

\begin{claim}{AG.18}{Proved; primitive-only vanishing already in genus two}
For the characteristic-two family $C_a$ in AG.15,
the Raynaud theta divisor contains a point of exact
order nine if and only if $a^5=1$.
At each such parameter the associated connected
cyclic etale degree-nine cover has genus ten and
two-rank two. Its primitive conductor-nine factor
has dimension six and two-rank zero, whereas its
degree-three intermediate cover is new-ordinary.
Thus primitive-only vanishing can occur with genus-two
base, not just in the genus-seventeen example AG.12.
\end{claim}
\begin{proof}
Use the local equation
$Y^2+Y=u^5+a u^3+c u$ at a selected affine point as
in AG.16; every value of $c$ can occur for fixed $a$.
A function with sole pole order nine has the form
\[
 F=A(u)+(u^2+b_1u+b_0)Y,\qquad \deg A\le4.
\]
Choose $A$ to cancel the coefficients through degree
four. With $v=c^4$, vanishing through degree eight
is precisely the following complete system:
\[
 \begin{aligned}
 b_0+b_1v+a&=0,& b_0a^2+b_1+v&=0,\\
 b_1a^2+1&=0,& b_0v^2+a^2&=0.
 \end{aligned}
\]
The third equation excludes $a=0$ and gives
$b_1=a^{-2}$. The first then gives
$b_0=v/a^2+a$. The remaining two equations become
\[
                  a^5+1=0,
       \qquad v^3+a^3v^2+a^4=0.
\]
These exact substitutions are independently checked in
\path{code/supersingular_nine.py}. If $a^5=1$, the
second polynomial is monic and has a root over the
algebraic closure. Construct $c$, the selected point,
and $F$ from that root. The function has pole exactly
nine at infinity and vanishes at least nine times at
the selected point. Its divisor is therefore exactly
$9R-9\infty$. The class has order nine: it is nonzero,
and order three is impossible by AG.15. This also
proves the necessity of the condition from any
order-nine theta point.

Multiplication by two is transitive on the six primitive
characters modulo nine. Their coherent eigenspaces are
one-dimensional; the theta point makes an arrow zero,
so AG.3 gives zero primitive rank. The sole remaining
nontrivial conductor is three. Its character subgroup
does not meet theta by AG.15, so the degree-three
intermediate new factor is ordinary of dimension two.
The base has two-rank zero. Hence the covering curve
has two-rank exactly two, and Riemann--Hurwitz gives
$g(D)=9(g(C_a)-1)+1=10$.

These parameters also admit the degree-seven whole-new
zero-rank covers of AG.15. Thus, as with AG.12, this
example separates conditions on the displayed cover;
it does not place the base curve outside the A3 class.
\end{proof}

\begin{claim}{AG.19}{Proved; the full-relative criterion for every deck order}
Let $D\to C$ be a connected etale geometrically Galois
cover of smooth proper curves over a finite field $L$
of characteristic $p$, with $g(C)\ge2$, nontrivial
geometric group $G$ of order prime to $p$, and a rational
point $P'$ above $P\in C(L)$. If
\[
                       p\nmid h_L(D)/h_L(C),
\]
then $P$ collides with no other point of $C(L)$.
There is no odd-order restriction and no constant-deck
restriction. For $|G|>2$ the proof is elementary once
the rational projector is established. The degree-two
case uses the verified odd-characteristic hyperelliptic
case of R4.
\end{claim}
\begin{proof}
We first prove the exact geometric fact needed here.
If a finite group acts freely on a smooth proper curve
$D$ of genus at least two, and $P',Q'$ lie in distinct
orbits, then
\[
             \#\operatorname{Stab}_G[Q'-P']\le2.       \tag{AG.19.1}
\]
Indeed, a nonidentity stabilizer element $\delta$ gives
the principal degree-two divisor of AG.17. The resulting
map to $\mathbf P^1$ is separable: the only possible
inseparable case is a purely inseparable extension of a
rational function field of degree two in characteristic
two, and that extension is again rational, contrary to
$g(D)\ge2$. Thus $D$ has its unique hyperelliptic
involution $\iota$. The unequal but linearly equivalent
degree-two divisors
$\delta Q'+P'$ and $\delta P'+Q'$ belong to its pencil.
Consequently
\[
                   \delta Q'=\iota P',\qquad
                   \delta P'=\iota Q'.
\]
There is at most one such nonidentity $\delta$, because
the action is free and its value on $Q'$ has now been
specified. This proves (AG.19.1) in every characteristic.

If $P,Q$ collide, the pointing lifts them to rational
$P',Q'$ with a nonzero $p$-primary difference $T$ and
equal full prime-to-$p$ section classes upstairs. As in
AC.18, the rational projector
$|G|^{-1}\pi^*\pi_*$ and the hypothesis kill the new
part of $T$. Thus the whole geometric group fixes $T$.
For $|G|>2$ this contradicts (AG.19.1).

If $|G|=2$, the characteristic is odd. The same
degree-two divisor argument makes $D$ hyperelliptic.
The lifted points still collide on $D$ by the full
subgroup-topology statement in AG.2. But R4, specifically
AP-A.17--A.19 with the induction of AP-A.15, proves
injectivity for every hyperelliptic curve in odd
characteristic over every finite extension. This is the
remaining contradiction.

This extends the sufficient criterion; it does not turn
an even-order witness into an odd-order witness. The
existence assertion A3 and its odd-group reductions
therefore retain their own quantifiers.
\end{proof}

\begin{claim}{AG.20}{Proved; an ordinary elliptic-pullback barrier in characteristic two}
Let $C$ be an ordinary smooth proper curve of genus at
least two over $\overline{\mathbf F}_2$, with a separable
degree-two map $C\to E$ to an ordinary elliptic curve.
For every connected odd-degree etale isogeny $E'\to E$,
the pullback $D=C\times_E E'$ is ordinary. In particular
no such $D$ admits a nonconstant morphism to a
positive-genus curve of two-rank zero, or a nonzero
homomorphism from its Jacobian onto an abelian variety
of two-rank zero.

This applies to both elliptic quotients of the ordinary
genus-two curves in AC.25. Thus those isogeny-pullback
families supply neither the rational new-part vanishing
of AC.24 nor the geometric zero-rank targets in L5.
It is a restriction on these specified families, not on
all odd-degree covers of $C$.
\end{claim}
\begin{proof}
Write $n=\deg(E'\to E)$ and let $r$ be the number
of geometric ramification points of $C\to E$.
The quadratic extension is ramified, because
Riemann--Hurwitz and $g(C)\ge2$ exclude an unramified
cover of an elliptic curve. It is linearly disjoint from
the odd-degree isogeny, so $D$ is connected.
Since $D\to C$ is etale of degree $n$,
\[
                      g(D)-1=n\bigl(g(C)-1\bigr).
\]
The elliptic curve $E'$ is ordinary because it is
isogenous to $E$. The quadratic map $D\to E'$ has
exactly $nr$ geometric ramification points; etale base
change preserves their local quadratic extensions.
The Deuring--Shafarevich formula
\citep[Corollary~1.8]{AG_Crew1984} gives
\[
 f_2(C)-1=2(f_2(E)-1)+r=r,
 \qquad f_2(D)-1=2(f_2(E')-1)+nr=nr.
\]
Ordinariness of $C$ makes $r=g(C)-1$, so these
equalities prove $f_2(D)=g(D)$.

Every abelian quotient of an ordinary abelian variety
is ordinary: Poincare reducibility expresses the source
up to isogeny as the quotient times a complementary
abelian variety, and both dimension and two-rank are
additive and isogeny invariant. A nonconstant morphism
of positive-genus curves gives a surjective Jacobian
pushforward. Its target therefore cannot have positive
dimension and two-rank zero. This proves the assertion.
\end{proof}

\begin{claim}{AG.21}{Proved; the complete degree-thirteen theta locus}
For the characteristic-two family $C_a$ of AG.15, put
$U=a^5$. Its Raynaud theta divisor contains a point
of exact order thirteen if and only if
\[
 (U+1)(U^3+U+1)(U^4+U+1)(U^6+U^5+1)=0.              \tag{AG.21.1}
\]
Each such parameter admits a connected cyclic etale
degree-thirteen cover with zero full new two-rank.
At every other parameter all degree-thirteen covers
are new-ordinary. In particular every degree-thirteen
cover of $C_0:y^2+y=x^5$ has ordinary new part.
\end{claim}
\begin{proof}
Use the translated local equation of AG.16,
$Y^2+Y=u^5+a u^3+c u$, and write $v=c^4$.
A function with pole exactly thirteen at infinity has
the form
\[
 F=A(u)+(u^4+b_3u^3+b_2u^2+b_1u+b_0)Y,
                  \qquad \deg A\le6.
\]
The polynomial $A$ cancels terms through degree six.
The remaining coefficients through degree twelve give
exactly the following six equations:
\[
 \begin{aligned}
 b_1a^2+b_2+b_3v+a&=0,&
 b_0v^2+b_2a^2+b_3+v&=0,\\
 b_1v^2+b_3a^2+1&=0,&
 b_0+b_2v^2+a^2&=0,\\
 b_1+b_3v^2&=0,&
 b_0a^4+b_2+v^2&=0.
 \end{aligned}
\]
The fifth and fourth equations give
$b_1=b_3v^2$ and $b_0=b_2v^2+a^2$.
Put $w=v^4+a^2$. The third equation gives
$b_3w=1$, so $w\ne0$. The first gives
$b_2=a+(a^2v^2+v)/w$. Substitution in the second
now gives $aw^2=1$, in particular $a\ne0$.

Write $a=t^8$ with $t\ne0$, uniquely in the algebraic
closure. The preceding equations force
\[
 \begin{gathered}
 v=t^4+t^{-1},\qquad b_3=t^4,\qquad b_1=t^{12}+t^2,\\
 b_2=t^{28}+t^{18}+t^3,
 \qquad b_0=t^{36}+t^{11}+t.
 \end{gathered}
\]
The last coefficient equation, multiplied by $t^2$,
is $P(t^5)=0$, where
\[
 P(U)=U^{14}+U^9+U^7+U^6+U^4+U^2+U+1.
\]
Since $a^5=(t^5)^8$ and $P$ has coefficients in
$\mathbf F_2$, this is equivalent to $P(a^5)=0$.
Direct multiplication gives the factorization in
(AG.21.1). The exact program
\path{code/supersingular_thirteen.py} checks all six
substitutions and the factorization.

Conversely, a root of (AG.21.1) has $a\ne0$; take
its unique eighth root $t$ and the displayed coefficients.
All six equations hold. The local parameter $c$ and
point on $C_a$ exist by AG.16. The resulting function
has pole exactly thirteen and zero at least thirteen
at that point, so its divisor is $13R-13\infty$.
Its nonzero class has exact prime order thirteen.
Finally two generates $(\mathbf Z/13)^\times$,
so the one-dimensional Frobenius-orbit criterion AG.3
gives the last assertions.
\end{proof}

\begin{claim}{AG.22}{Refuted; scalar base Frobenius does not identify the rational and geometric tests}
Let $C_0/\mathbf F_{16}$ be $y^2+y=x^5$, with
$P=\infty$. The Frobenius of $J_{C_0}$ is the scalar
$[-4]$. Over $L=\mathbf F_{4096}$ there exists a
constant cyclic degree-thirteen $P$-split cover whose
relative class number is odd, although its entire new
factor is ordinary of dimension twelve.
Moreover $L$ is the first extension of $\mathbf F_{16}$
over which any constant cyclic degree-thirteen cover
can occur. Thus the distinction already occurs at the
minimal constant-deck field under scalar base Frobenius.
\end{claim}
\begin{proof}
The exact point counts are
\[
            \#C_0(\mathbf F_{16})=33,
             \qquad \#C_0(\mathbf F_{4096})=4353.
\]
Both are independently enumerated in
\path{code/scalar_frobenius_count.py} by the absolute
trace criterion for $y^2+y=x^5$; the complete output
is attached. The first achieves the upper Weil bound,
so all four Frobenius eigenvalues are $-4$.
Frobenius is semisimple on a prime-to-two Tate module,
and that realization is injective on endomorphisms
\citep[Sections~1--2, pp.~134--139]{AG_Tate1966}.
Consequently the endomorphism itself is $[-4]$.
Its cube is $[-64]$, which also yields the second count
from the trace formula, and
\[
 J_{C_0}(L)=J_{C_0}[65](\overline{\mathbf F}_2)
                      \simeq(\mathbf Z/65)^4.
\]
Set $H=13J_{C_0}(L)=J_{C_0}[5]$, so $|H|=625$,
and put $A=J_{C_0}(L)/H\simeq(\mathbf Z/13)^4$.
Write $N_b$ for the number of $R\in C_0(L)$ with
$[R-P]$ in the coset $b\in A$.
The zero coset has $N_0=33$: $[5][R-P]=0$ is
equivalent to the class being fixed by $[-4]$, hence,
by the Abel embedding based at $P$, to
$R\in C_0(\mathbf F_{16})$.

Suppose every nontrivial quotient $A\to\mathbf Z/13$
failed the pointed rational criterion. Since two is
primitive modulo thirteen, failure at one nontrivial
character forces failure at all its nontrivial powers.
The genus-two formula in AC.5 therefore gives
\[
        \sum_{b\in A}N_b\psi(b)=1\quad\hbox{in }
        \overline{\mathbf F}_2
        \qquad(1\ne\psi\in\widehat A).
\]
The trivial character has the same value because
$4353$ is odd. Fourier inversion, valid since $|A|$
is odd, gives $N_b\equiv0\pmod2$ for every $b\ne0$.
Every one of the $4353-33=4320$ points in nonzero
cosets must therefore share its coset with another point.

For any nonzero $h\in J_{C_0}(L)$ there are at most
two ordered pairs of distinct curve points $(R,S)$
with $[R-S]=h$. Indeed an equality of two such
differences is an equality of effective degree-two
divisor classes; the unique hyperelliptic pencil is
the only possibility for unequal effective divisors,
giving the two exchanged ordered pairs. Therefore
the total number of ordered distinct-point pairs in
the same $H$-coset is at most
$2(|H|-1)=1248$. This contradicts $4320>1248$.
Some prime cyclic quotient is a pointed rational
witness. This is the specified application of the
simultaneous-quotient bound in AC.32; the criteria
agent's exact certificate also supplies a concrete
Mumford divisor and its full $L$-polynomials.

AG.21 shows that theta has no thirteen-torsion when
$a=0$, so every such cover has ordinary new factor,
of dimension twelve. Finally $-4$ has order three
modulo thirteen. The scalar Frobenius on the geometric
abelianization has a nonzero constant cyclic quotient
only after an extension of degree divisible by three.
This proves the minimal-field assertion.
\end{proof}

\subsection{Exact plane-quartic character computation}

\begin{claim}{AG.23}{Proved; a finite divisor algorithm for pointed plane quartics}
Let $C/\mathbf F_q$ be a smooth plane quartic, let
$P\in C(\mathbf F_q)$, and let $K$ be its canonical
divisor class, represented by a line section. From the closed
points of degrees one, two and three, their local expansions,
and linear algebra over $\mathbf F_q$, one can compute the
entire group $J_C(\mathbf F_q)$ and every pointed cyclic
prime-order character. The algorithm below gives the exact
twisted $L$-polynomial for every such character; it does not
replace the Jacobian group law by the numerical zeta polynomial.
\end{claim}
\begin{proof}
Put $B=3P$. Every degree-three divisor class has an effective
representative, by Riemann--Roch. Moreover an effective
degree-three divisor $D$ satisfies
\[
             h^0(D)=1+h^0(K-D).
\]
The degree-one class $K-D$ is effective precisely when it is
the class of a point $R\in C(\mathbf F_q)$. In this case
$D\sim K-R$ and $h^0(D)=2$; otherwise $h^0(D)=1$.
Distinct points give distinct degree-one classes, since a
function with divisor $R-S$ would have degree one on a
genus-three curve. Thus all effective degree-three divisors
represent different classes, except that, for each rational
point $R$, the $q+1$ members of $|K-R|$ represent one class.
This is an exact and exhaustive equivalence test.

The test is implemented using lines. A nonzero line section
vanishing along $D$ exists precisely in the special case;
its residual intersection is the point $R$. No two independent
line sections can vanish along a degree-three divisor on a
smooth plane quartic: their common zero in the plane is a
single point, and at most one line is tangent there. Thus
the computed residual point is unambiguous. A unique
representative for every resulting class identifies
$\operatorname{Pic}^3(C)(\mathbf F_q)$ with
$J_C(\mathbf F_q)$ by $D\mapsto[D-B]$.

Here is the group law, including nonreduced divisors and
nonrational closed points. Given effective representatives
$D,E$ of degree three, choose a nonzero cubic section
vanishing along $D+E+B$. Its existence follows from
\[
 \deg(3K-D-E-B)=3,\qquad h^0(3K-D-E-B)\ge1.
\]
All sections of $3K$ are restrictions of homogeneous cubics:
both spaces have dimension ten, and a nonzero cubic cannot
vanish identically on a smooth quartic. The residual divisor
$R$ has degree three and
\[
                         R\sim3K-D-E-B.
\]
Choose a second cubic vanishing along $R+2B$, with residual
divisor $S$ of degree three. Then
\[
 S\sim3K-R-2B\sim D+E-B,
 \qquad [S-B]=[D-B]+[E-B].
\]
In particular the result is independent of both cubic choices.
All the residual closed points have degree at most three,
so the initial finite list contains every point that can occur.

For completeness the local multiplicity conditions are exact
linear conditions, not evaluations at the support alone. At
each closed point of degree $d$, choose an affine chart and
a coordinate with nonzero partial derivative. Take the other
coordinate as a uniformizer and solve the curve equation
recursively as a formal series through order $12/d$.
The nonzero derivative makes every coefficient uniquely
determined. A multiplicity $m$ imposes the vanishing of
the first $m$ coefficients, each in $\mathbf F_{q^d}$,
or exactly $md$ linear conditions over $\mathbf F_q$.
Intersection orders are read from the same series. B\'ezout's
degree bound is four for a line and twelve for a cubic;
the program checks that all residual degrees are exactly
one or three as appropriate.

For every prime $\ell\ne\operatorname{char}(\mathbf F_q)$
dividing the computed group order, form the image
$\ell J_C(\mathbf F_q)$, its cosets, and a basis of the
elementary abelian quotient. The nonzero linear forms on
that quotient, modulo multiplication by
$\mathbf F_\ell^\times$, give exactly every prime cyclic
quotient, once each. The pointed Lang construction of AC.2
gives a geometrically connected constant cyclic cover with
a rational lift above $P$ for each quotient.

Let $\chi$ denote a selected character and put
$\kappa=[K-4P]$. Its polynomial is
\[
 \begin{split}
 L_P(\chi,T)={}&1+c_1(\chi)T+c_2(\chi)T^2
       +q\chi(\kappa)c_1(\chi^{-1})T^3
       +q^2\chi(\kappa)T^4,\\
 c_i(\chi)={}&\sum_{D\ge0,\,\deg D=i}\chi([D-iP])
                                      \quad(i=1,2).
 \end{split}                                                   \tag{AG.23.1}
\]
The first two sums are directly enumerable from the closed
points already listed. The degree-three coefficient follows
also without a functional-equation convention: every
degree-three class has one effective representative, except
$K-R$, which has $q+1$. The uniform contribution cancels
for a nontrivial character; the extra contribution is
$q\sum_R\chi([K-R-3P])$. Similarly every degree-four
class has $q+1$ effective representatives except $K$, which
has $q^2+q+1$. The extra contribution is
$q^2\chi(\kappa)$. For larger degrees Riemann--Roch makes
the number of effective representatives constant in every
class, so all subsequent coefficients vanish. This proves
(AG.23.1) and its exact normalization at $P$.

Finally evaluate at $T=1$ and take the norm from
$\mathbf Q(\zeta_\ell)$. This is the relative class number.
Its reduction modulo the characteristic is computed exactly
as a polynomial resultant with
$1+X+\cdots+X^{\ell-1}$; nonvanishing is also checked by
the independent Euclidean gcd test. Every conjugate
$L_P(\chi^i,T)$ is given by the substitution
$\zeta_\ell\mapsto\zeta_\ell^i$ in the retained coefficient
vectors. These steps compute characters of the actual
finite Jacobian, including the chosen base-point twist.
\end{proof}

\begin{claim}{AG.24}{Proved by exact computation; every pointed plane quartic over $\mathbf F_3$}
There are exactly $3001$ isomorphism classes of pointed
smooth plane quartics $(C,P)$ over $\mathbf F_3$.
For every one of them the program
\path{code/plane_quartic_F3.cpp} computes all constant
odd-prime cyclic pointed tests over $\mathbf F_3$.
It retains the exact $L$-polynomials for all $2692$ tested
quotients, including failures. There is a successful
base-field witness for $2368$ pointed classes.
For $637$ classes a degree-five witness succeeds; five is
therefore their least admissible odd prime over all finite
extensions. For the other $2364$ classes this calculation
does not determine the all-extension minimum.

The $633$ classes without a base-field witness are not
asserted to lack witnesses after extension. Likewise,
the smallest successful prime in the base-field table
is not called the all-extension minimum when it exceeds
five. The $\mathbf F_5$ quartic calculation is given in
AG.25--AG.26, with the same distinction between base-field
and all-extension minima.
\end{claim}
\begin{proof}
Move $P$ to $(0:0:1)$ and its tangent to $Y=0$.
After scaling the defining equation, every pointed quartic
has the normalized form
\[
 F=YZ^3+(aX^2+bXY+cY^2)Z^2
       +(dX^3+eX^2Y+fXY^2+gY^3)Z+Q_4(X,Y),
\]
with twelve arbitrary coefficients in $\mathbf F_3$.
There are $3^{12}=531441$ such forms. The complete
stabilizer of the point and tangent in
$\operatorname{PGL}_3(\mathbf F_3)$ consists of
\[
 (X,Y,Z)\longmapsto(AX+BY,CY,DX+EY+Z),
 \quad A,C\in\mathbf F_3^\times,
 \quad B,D,E\in\mathbf F_3,
\]
and has order $108$. The program applies every one of
these transformations, rescales the $YZ^3$ coefficient
to one, and marks every orbit. Since a smooth plane
quartic is canonically embedded, every isomorphism of
pointed curves is induced by a projective linear
transformation. Thus these are exactly the pointed
isomorphism classes, including classes with automorphisms.
Pointless quartics contribute no pointed pair; each rational
point on every other quartic is covered up to pointed
isomorphism by this table.

Smoothness is decided over the algebraic closure. The three
partial derivatives of $F$ are homogeneous cubics. As
$4$ is invertible in $\mathbf F_3$, their common
projective zero locus is exactly the singular locus.
They have no projective common zero if and only if
\[
  S_4^{\oplus3}\longrightarrow S_7,
  \qquad(A,B,C)\longmapsto AF_X+BF_Y+CF_Z
\]
is onto, where $S=\mathbf F_3[X,Y,Z]$.
Indeed, in the absence of a projective zero the three
cubics form a regular sequence, whose quotient has
Hilbert series $(1+T+T^2)^3$ and hence vanishing
degree-seven part. Conversely surjectivity excludes a
common projective zero by evaluation. This is a rank
$36$ test on a $36$ by $45$ matrix. The complete
enumeration gives $315360$ smooth normalized forms,
partitioned into the $3001$ displayed orbits.

For each representative, AG.23 is implemented with
$\mathbf F_9=\mathbf F_3[t]/(t^2+1)$ and
$\mathbf F_{27}=\mathbf F_3[t]/(t^3+2t+1)$.
The field tables check an inverse for every nonzero
element. Closed points are full Frobenius orbits,
with projective coordinates normalized by their first
nonzero coordinate. The coefficient order for a quartic
is $X^iY^jZ^{4-i-j}$, first $i=0,\ldots,4$, then
$j=0,\ldots,4-i$; all files use that order.

The program independently counts the curve over
$\mathbf F_3,\mathbf F_9,\mathbf F_{27}$, obtains
the degree-six zeta numerator by Newton identities,
and checks that its value at one equals the number
of computed degree-three divisor classes. It also
checks that the number of effective degree-three
divisors is $h(C)+3\#C(\mathbf F_3)$, that every
special class has exactly four representatives,
that the two-cubic group law has the correct identity,
and that the group order annihilates every class.
Associativity is checked on a deterministic set of
triples for each curve. Its validity in general follows
from the divisor-class calculation in AG.23.

For every character the degree-three coefficient is
additionally enumerated directly. As a vector indexed
by $\mathbf Z/\ell$, its difference from the reciprocal
formula in (AG.23.1) is exactly
$(h(C)/\ell)(1+X+\cdots+X^{\ell-1})$; the program
checks every coefficient of this identity. Resultant
and gcd tests agree for all $2692$ quotients.
The separate program \path{code/review_quartic_norms.py}
also computes the full positive integer resultant for
every quotient. All $2692$ exact relative class numbers
are retained in
\path{code/A4_quartic_F3_exact_norms.csv}, and each has
the same residue modulo three as the finite-field
algorithm. This independently checks the resultant sign,
including every failed test.
The complete audited run gives the counts in the
claim. The JSON lines retain all closed points,
all Jacobian representatives, quotient bases,
coordinates of every class in each quotient,
character vectors, exact cyclotomic coefficient
vectors and relative class numbers modulo three.
The companion CSV records every test and its result.
Since three is the characteristic and the criterion
here retains odd degree, five is the smallest allowed
prime; hence the $637$ degree-five successes have
the stated global minimum.
\end{proof}

\begin{claim}{AG.25}{Proved; exhaustive small slices for the $\mathbf F_5$ quartic computation}
Every pointed smooth plane quartic over $\mathbf F_5$
has a normalized representative of precisely one of the
following three types. Write $u,v$ for the coefficients
of $X^4,X^3Y$, respectively, and retain the notation
of AG.24 for the other coefficients:
\[
 \begin{array}{c|c|c|c}
 \text{Contact}&\text{Fixed coefficients}
       &\text{Free}&\text{Weight}\\ \hline
 2&a=1,\ b=c=d=0&8&2\\
 3&a=0,\ d=1,\ b=c=e=0&7&3\\
 4&a=d=0,\ u=1,\ b=c=v=0&6&4.
 \end{array}
\]
In a row of weight $w$, two normalized forms are pointed
isomorphic exactly when related by
$(X,Y,Z)\mapsto(AX,A^wY,Z)$ with
$A\in\mathbf F_5^\times$, followed by division by $A^w$.
Thus $5^8+5^7+5^6=484375$ forms, each under a group of
order four, give an exhaustive census. The exact geometric
smoothness test of AG.24 gives $94283$ smooth pointed
isomorphism classes, corresponding to $187512000$ of
the original $5^{12}$ normalized coefficient forms.
\end{claim}
\begin{proof}
The tangent section is
\[
                  F(X,0,Z)=aX^2Z^2+dX^3Z+uX^4.
\]
If $a=d=u=0$, the line $Y=0$ is a component, impossible
for a smooth plane quartic. The three displayed cases
therefore exhaust all possibilities and their contact
orders distinguish them intrinsically.

Apply the general point-and-tangent stabilizer from AG.24,
now over $\mathbf F_5$, and divide the transformed equation
by $C$. Its first relevant coefficients are
\[
 \begin{aligned}
 a'&=aA^2/C,&
 b'&=3D+2aAB/C+bA,\\
 c'&=3E+aB^2/C+bB+cC,&
 d'&=2aA^2D/C+dA^3/C.
 \end{aligned}                                                \tag{AG.25.1}
\]
If $a\ne0$, choose, for any nonzero $A$,
\[
 C=aA^2,\qquad D=-dA/(2a),\qquad
 B=-A(3D+bA)/2,
\]
and then the unique $E$ making $c'=0$. This gives
the first row.

If $a=0$ and $d\ne0$, take
\[
 C=dA^3,\qquad D=-bA/3.
\]
The transformed coefficient of $X^2YZ$ is then
\[
                   e'=3D^2+2bAD+3B/A+eA^2.
\]
Choose the unique $B$ making this zero, followed by
the unique $E$ making $c'=0$. This gives the second row.

Finally let $a=d=0$ and $u\ne0$. Take
$C=uA^4$ and $D=-bA/3$. The transformed coefficient
of $X^3Y$ is
\[
                 v'=D^3+bAD^2+eA^2D+4B/A+vA^3.
\]
There is a unique $B$ making it zero, and then a unique
$E$ making $c'=0$. This gives the third row. Every
division used here is legitimate in characteristic five.

Conversely, a stabilizer element carrying one form of a
fixed row to another must have $C=A^w$ and successively
$D=B=E=0$, by exactly the displayed coefficients.
The residual group is therefore precisely the stated
diagonal group. Its action on a free coefficient of
$X^iY^jZ^{4-i-j}$ is multiplication by
$A^{i+wj-w}$, which is what the program enumerates.

The full stabilizer has order
$(5-1)^2\,5^3=2000$. Each normalized slice orbit
of size $s$ consequently corresponds to $500s$
forms in the original coefficient space. Enumerating
the three slices, marking every diagonal orbit, and
performing the algebraic-closure smoothness rank test
gives the stated $94283$ classes and weighted count.
These are finite exact calculations retained in
\path{code/plane_quartic_F5_census.csv}; the subsequent
Jacobian tests use the same representatives.
\end{proof}


\begin{claim}{AG.26}{Proved by exact computation; every pointed plane quartic over $\mathbf F_5$}
For all $94283$ pointed smooth plane-quartic isomorphism
classes over $\mathbf F_5$, the program
\path{code/plane_quartic_F5.cpp} computes every constant
odd-prime cyclic pointed test over the base field.
There are $134422$ such quotient tests, and their exact
cyclotomic $L$-polynomials are retained, including failures.
A successful base-field witness exists for $89906$ classes;
$4377$ classes have no such base-field witness.
For $39917$ classes a degree-three witness succeeds.
Consequently three is their least admissible odd prime
over all finite extensions. The all-extension least prime
for each of the other $54366$ classes remains undetermined
by this calculation, even when its base-field table has
a larger successful prime.
\end{claim}
\begin{proof}
The exhaustive representatives and the algebraic-closure
smoothness test are those of AG.25. Apply the exact group
and character algorithm AG.23 to each representative,
with the explicit fields
\[
 \mathbf F_{25}=\mathbf F_5[t]/(t^2+2),\qquad
 \mathbf F_{125}=\mathbf F_5[t]/(t^3+t+1).
\]
Both displayed polynomials are irreducible over
$\mathbf F_5$: the quadratic has no root, and direct
substitution of $0,1,2,3,4$ gives no root for the cubic.
The finite-field tables additionally check every nonzero
inverse. The initial projective closed-point lists have
$31$, $310$ and $5240$ points of degrees one, two and three,
respectively, as follows either from their explicit Frobenius
orbits or from the counts of $\mathbf P^2$ over the three
fields. Every local multiplicity, cubic complement,
Jacobian operation and character value is computed over
these finite fields, with no floating-point arithmetic.

The implementation verifies, for every curve, that the
computed number of Jacobian classes agrees with the zeta
polynomial from the three point counts. Writing
$s_i=5^i+1-\#C(\mathbf F_{5^i})$, its first coefficients
are
\[
 a_1=-s_1,\qquad
 a_2=(s_1^2-s_2)/2,\qquad
 a_3=(-s_3-a_1s_2-a_2s_1)/3,
\]
and
\[
                   h(C)=126+26a_1+6a_2+a_3.
\]
The divisions by two and three here are divisions in the
integers; the program checks their exactness. It also checks
that there are $h(C)+5\#C(\mathbf F_5)$ effective
degree-three divisors, exactly six representatives for each
special class, the correct group identity, annihilation of
every class by $h(C)$, and the same deterministic
associativity checks as in AG.24. The proof of the group
law itself is AG.23.

Every odd prime divisor of $h(C)$ except the characteristic
five is used. For each prime the actual quotient
$J_C(\mathbf F_5)/\ell J_C(\mathbf F_5)$ is computed,
and every linear form up to scalar is tested. The direct
degree-three coefficient enumeration agrees coefficient by
coefficient with the reciprocal formula (AG.23.1), with
exact uniform difference $(h(C)/\ell)\Phi_\ell$.
The resultant and Euclidean-gcd tests agree on all
$134422$ characters. The independent integer-arithmetic
program \path{code/review_quartic_norms.py} additionally
computes the full positive relative class number for every
one of these tests. All $134422$ integer values agree
modulo five with the finite-field algorithm, including every
failed test; they are retained in
\path{code/A4_quartic_F5_exact_norms.csv}. Thus the
computed norms are checked as integers, as well as by
two independent finite-field algorithms.
All coefficient vectors, Jacobian
representatives, quotient coordinates, base-point shifts,
and both outcomes are retained in
\path{code/plane_quartic_F5_tests.jsonl}; the companion
CSV lists every quotient and test outcome.

The run is partitioned into four disjoint index classes
modulo four. The merging program
\path{code/launch_quartic_F5.py} checks that their census
rows are exactly the original $94283$ representatives,
with no omissions or repetitions, before accepting the
aggregate. The complete exact output is
\path{code/plane_quartic_F5_summary.json}. It gives the
counts in the claim, together with the distribution of
least successful base-field primes. Since three is the
smallest allowed odd prime in characteristic five, each
degree-three success establishes the all-extension minimum
without any need to exclude smaller extensions or primes.
The other records retain their explicitly limited scope.
\end{proof}


\begin{claim}{AG.27}{Proved; universal degree obstructions for geometric primitive vanishing}
Let $C/\overline{\mathbf F}_2$ have genus two and let
$n>1$ be odd. Put
\[
 U_n=\langle2,-1\rangle\subset(\mathbf Z/n)^\times,
 \qquad r_n=\frac{2\varphi(n)}{|U_n|}.
\]
If a connected cyclic etale degree-$n$ cover of $C$
has primitive conductor-$n$ factor of two-rank zero, then
\[
                   r_n(r_n-1)\le2(n-1).                 \tag{AG.27.1}
\]
If $C$ is nonordinary, the stronger necessary condition is
\[
                   r_n(r_n+1)\le2(n-1).                 \tag{AG.27.2}
\]
Consequently degree $127$ and degree $257$ cannot have
primitive new two-rank zero for any genus-two base in
characteristic two. More generally, no degree
$n=2^d-1$ with $d\ge48$ can have this property.
These restrictions concern geometric two-rank, and assert
no corresponding exclusion of the rational class-number
criterion. They also do not exclude the possibility that
every curve admits some other degree as in A6.
\end{claim}
\begin{proof}
In characteristic two the Raynaud bundle is a degree-one
theta characteristic $B$ on the Frobenius twist of $C$.
Its theta divisor is the embedded curve
\[
 \Theta=\{L\in J_C:h^0(B\otimes L)>0\}=C-B.
\]
Since $B^{\otimes2}\simeq\omega_C$, Serre duality shows
$\Theta=-\Theta$. The harmless Frobenius twist does not
change any geometric torsion or intersection assertion
used here. The point zero belongs to $\Theta$ if and
only if $C$ is nonordinary: its multiplicity space
$H^0(B)$ is the kernel of Frobenius on $H^1(\mathcal O_C)$,
which vanishes exactly in the ordinary case.

Let $\Gamma\subset J_C$ be the cyclic order-$n$ group
of character line bundles for the cover. Every nontrivial
coherent character space has dimension one. By AG.3,
primitive rank zero requires each orbit of multiplication
by two on the primitive elements of $\Gamma$ to meet
$\Theta$. Write $f=\operatorname{ord}_n(2)$. If
$-1\in\langle2\rangle$, each such orbit must contain
at least two theta points, since a nonzero point of odd
order is distinct from its negative. There are
$\varphi(n)/f$ orbits, giving at least
$2\varphi(n)/f=r_n$ primitive theta points. If
$-1\notin\langle2\rangle$, the doubling orbits are
paired by negation. Each pair contributes at least two
points, or $\varphi(n)/f=r_n$ points in total.

For any nonzero $h\in J_C$, there are at most two
ordered pairs of distinct points of $\Theta$ with
difference $h$. This is the genus-two difference bound
proved in AG.22: unequal effective degree-two divisors
in one class belong to the unique hyperelliptic pencil,
which supplies at most the exchanged pair. Equivalently,
distinct translates of the theta curve have intersection
number two. Thus any $m$ distinct theta points belonging
to $\Gamma$ satisfy
\[
                         m(m-1)\le2(n-1).
\]
Use $m\ge r_n$ for (AG.27.1). In the nonordinary case,
the additional point zero is distinct from all primitive
points, so $m\ge r_n+1$ gives (AG.27.2).

For $127$, two has order seven and its subgroup does
not contain $-1$. Hence $|U_{127}|=14$ and $r_{127}=18$,
whereas $18\cdot17=306>252$. For $257$, two has
order sixteen and contains $-1$; therefore $r_{257}=32$
and $32\cdot31=992>512$. All group computations,
the necessary inequalities, and all odd degrees through
$2001$ are checked in
\path{code/primitive_degree_obstructions.py}, with its
complete exact output attached.

For the infinite family, an elementary bound valid for
odd $n$ is
\[
                       \varphi(n)\ge\frac23 n^{2/3}.
\]
Indeed, for every prime $t\ge5$ one has
$t-1\ge t^{2/3}$, and this gives
$\varphi(t^e)\ge(t^e)^{2/3}$. The possible factor
$3^e$ costs at most the constant $2/3$, since
$2\cdot3^{e-1}\ge(2/3)3^{2e/3}$ for $e\ge1$.
Multiply these prime-power bounds.
If $n=2^d-1$, then $|U_n|\le2d$, and consequently
\[
 r_n\ge\frac{\varphi(n)}d
      \ge\frac{2n^{2/3}}{3d}.
\]
For $d\ge48$ we have $n^{1/6}\ge3d$: it suffices
that $2^{(d-1)/6}\ge3d$, whose seed at $d=48$
and subsequent monotonicity follow from the exact integer
inequalities $2^{47}\ge144^6$ and
$2\cdot48^6>49^6$. The ratio
$2^{(d-1)/6}/d$ increases thereafter.
It follows that $r_n\ge2\sqrt n$, contradicting
(AG.27.1). This proves the stated infinite family of
geometrically excluded degrees.
\end{proof}


\begin{claim}{AG.28}{Proved; a linear rational-exponent criterion at a fixed field}
Let $L=\mathbf F_q$ have characteristic $p$, let $C/L$
be smooth proper and geometrically connected with $g(C)\ge2$,
and let $P\in C(L)$. Let $\pi:D\to C$ be a connected
finite etale Galois cover with constant group of order $m$
prime to $p$, and choose $P'\in D(L)$ over $P$.
Write $p^e$ for the exponent of $J_D(L)[p^\infty]$,
with $p^e=1$ when this group is zero. If
\[
                       (q-1)p^e<2(m-1),                \tag{AG.28.1}
\]
then $P$ collides with no other point of $C(L)$.
The group need not be nonabelian, and its order need not
be odd. More precisely, any collision forces
\[
 \operatorname{ord}[Q'-P']\ge\frac{2(m-1)}{q-1}
\]
for the rational lifts given by the collision.

Consequently, a tower of constant pointed prime-to-$p$
Galois covers over one fixed $L$, of degrees $m_i\to\infty$
and rational $p$-primary exponents $b_i=o(m_i)$, would
exclude every collision with $P$. Existence of such a tower
for every pointed curve is not established here.
\end{claim}
\begin{proof}
Suppose $P$ and a distinct $Q\in C(L)$ collide.
AR.1 gives a rational lift $Q'\in D(L)$ and a
nonzero $p$-primary class $T=[Q'-P']$ in $J_D(L)$.
Since the deck group is constant and the action is free,
every translate of $P'$ and of $Q'$ is rational.
The two fibres are disjoint, so $D(L)$ contains at least
$2m$ points.

Put $b=\operatorname{ord}T$, which divides $p^e$.
Over the algebraic closure there is a rational function
$f$ with
\[
                         \operatorname{div}(f)=bQ'-bP'.
\]
It can be chosen over $L$. Indeed the divisor is rational,
so every Galois conjugate of $f$ differs from it by a
constant. These constants form a multiplicative cocycle
in a finite extension over which $f$ is defined; Hilbert 90
rescales $f$ to an $L$-rational function with the same
divisor. The corresponding map $D\to\mathbf P^1_L$
has degree $b$, a single geometric point above zero,
and a single geometric point above infinity. Each of the
other $q-1$ rational values has at most $b$ rational
preimages, whether or not the map is separable. Therefore
\[
                        2m\le\#D(L)\le2+(q-1)b.
\]
This gives the sharper order bound and then
$2(m-1)\le(q-1)p^e$, contrary to (AG.28.1).

In the tower assertion $q$ is fixed, so $b_i=o(m_i)$
eventually implies (AG.28.1). This is a conditional
consequence, not a proof of the required growth hypothesis.
Enlarging the constant field changes both $q$ and the
rational torsion exponent; neither is held fixed silently.
\end{proof}


\begin{claim}{AG.29}{Proved; cyclotomic exponent bounds and further quartic certificates}
Let $D\to C$ be a constant pointed cyclic etale cover
of prime degree $\ell\ne p$ over $L=\mathbf F_q$,
with $g(C)\ge2$. Put
\[
 h=\#J_C(L),\quad h_{\rm rel}=\#J_D(L)/h,
 \quad f=\operatorname{ord}_\ell(p),\quad v=v_p(h_{\rm rel}).
\]
Choose a faithful character $\chi$ and write
$a=L_P(\chi,1)\in\mathbf Z[\zeta_\ell]$ as a polynomial
in $\zeta_\ell$. Let $b$ be the degree of its gcd with
$\Phi_\ell$ modulo $p$. Define
\[
 t=\begin{cases}
 0,&v=0,\\
 1+(v-b)/f,&v>0.
 \end{cases}
 \qquad B=p^{\max\{v_p(h),t\}}.                       \tag{AG.29.1}
\]
Then $t$ is a nonnegative integer and
\[
             \exp J_D(L)[p^\infty]\mid B.
\]
Consequently $(q-1)B<2(\ell-1)$ is a sufficient
no-collision criterion at $P$ by AG.28.
The weaker bound with $t=v/f$ is valid as well.

Applying this bound to the complete, already certified
quartic polynomials and integer norms of AG.24 and AG.26
gives the following exact base-field conclusions:
\[
 \begin{array}{c|r|r|r|r}
 q&\text{Pointed classes}&\text{L1 successes}
   &\text{Additional successes}&\text{Total certified}\\ \hline
 3&3001&2368&33&2401\\
 5&94283&89906&123&90029.
 \end{array}
\]
Every additional class has failed every base-field
constant odd-prime L1 test, yet has a successful exponent
certificate. There are $52$ individual covers over
$\mathbf F_3$ and $272$ over $\mathbf F_5$ which fail
L1 and satisfy this exponent criterion. This is a strict
enlargement of the verified sufficient tests over the
fixed base field. It makes no assertion that those
pointed curves lack an L1 witness after extension.
\end{claim}
\begin{proof}
Let $P_D=\ker(J_D\to J_C)^0$ be the connected Prym.
The pullback--norm isogeny
$J_C\times P_D\to J_D$ has kernel killed by a power
of $\ell$, so it induces an isomorphism on $p$-divisible
groups over $L$. In particular its rational $p$-primary
group splits into the base and Prym parts. The base part
has exponent dividing $p^{v_p(h)}$.

Let $\sigma$ generate the deck group. The relation
\[
 \Phi_\ell(\sigma)=1+\sigma+\cdots+\sigma^{\ell-1}
                       =\pi^*\pi_*
\]
vanishes on $P_D$. Therefore the full integer ring
$\mathcal O_K=\mathbf Z[\zeta_\ell]$, with
$K=\mathbf Q(\zeta_\ell)$, acts integrally on $P_D$,
not merely up to isogeny. Frobenius $F$ commutes with
this action because the deck group is constant over $L$.
Its $K$-linear characteristic polynomial $f_F(Z)$
has degree $2g(C)-2$ and coefficients in $\mathcal O_K$;
its conjugates are the reciprocal character polynomials
$Z^{2g(C)-2}L_P(\chi^i,Z^{-1})$. One may choose the
embedding of $K$ so the identity character corresponds
to the displayed $\chi$; reversing the deck generator
only permutes these embeddings. Integrality follows
also because its roots are Frobenius algebraic integers.

Cayley--Hamilton on an auxiliary prime-to-$p\ell$
Tate module, followed by faithfulness of that realization,
gives the endomorphism identity $f_F(F)=0$ on $P_D$.
On $P_D(L)[p^\infty]$ the endomorphism $F$ is one,
so $f_F(1)=a$ annihilates this finite group. Moreover
\[
                      N_{K/\mathbf Q}(a)=h_{\rm rel}.
\]
These statements use the actual Frobenius endomorphism
and the integral cyclotomic action; there is no inference
of an integral group structure from the numerical zeta
polynomial alone.

Because $p\ne\ell$, the algebra
\[
 \mathcal O_K\otimes\mathbf Z_p=\prod_{i=1}^{(\ell-1)/f}O_i
\]
is a product of unramified discrete valuation rings,
each of residue degree $f$. Write $t_i=v_p(a_i)$ for
the image of $a$ in $O_i$. The corresponding factor
of the finite Prym group is annihilated by $p^{t_i}$:
$a_i$ is $p^{t_i}$ times a unit of $O_i$.
Thus its exponent divides $p^{\max_i t_i}$, and
\[
                       v=f\sum_i t_i.
\]
Exactly $b/f$ of the integers $t_i$ are positive.
Indeed $\Phi_\ell$ modulo $p$ is square-free, with
irreducible factors all of degree $f$, and its common
factors with $a$ are precisely the residue factors
where $a_i$ is zero. If $v=0$, every $t_i$ is zero.
Otherwise
\[
 \max_i t_i\le\sum_i t_i-\bigl(b/f-1\bigr)
                         =1+(v-b)/f=t.
\]
This proves (AG.29.1), the divisibility assertion,
and the simpler bound $\max_i t_i\le v/f$.

The exact script \path{code/quartic_rational_exponent.py}
reads the retained Jacobian orders, mod-$p$ gcd degrees,
and independently computed integer relative class numbers.
It verifies $f\mid v$, $f\mid b$, $v\ge b$, and
$v=0\Longleftrightarrow b=0$ for every quotient,
then applies (AG.29.1) and the strict inequality of AG.28.
The complete successful certificate tables are
\path{code/plane_quartic_F3_rational_exponent_certificates.csv}
and
\path{code/plane_quartic_F5_rational_exponent_certificates.csv},
with companion JSON summaries. Each certificate records
its pointed model, prime quotient character, both exact
class numbers, $f,v,b$, the exponent upper bound,
and the two sides of the strict inequality. Comparing
with every L1 test for the same pointed model gives
exactly the additional counts in the table. The weaker
norm-only scan and the initially tried total-class-number
bound are also retained, with their outputs, so that no
failed or superseded attempt is suppressed.
\end{proof}


\begin{claim}{AG.30}{Proved; the geometric limitation of two ordinary elliptic targets}
Let $\Omega$ be algebraically closed, let $D/\Omega$
be a smooth proper curve of genus at least two, let $A$
be an abelian surface, and let $u:A\to A'$ be an isogeny
with a nonzero point in its geometric kernel.
For every morphism $\phi:D\to A$, the composite
$u\phi$ fails to be injective on geometric points.
In particular, if $A=E_1\times E_2$ is a product of
ordinary elliptic curves in characteristic two, then
no joint map
\[
 D\longrightarrow E_1\times E_2
       \xrightarrow{[2^{e_1}]\times[2^{e_2}]}
                    E_1\times E_2,
 \qquad e_1,e_2\ge1,
\]
is geometrically point-separating.

Thus replacing each of two ordinary elliptic targets by
multiplication by its rational two-primary exponent cannot
supply a geometrically point-separating family for L6.
This does not exclude separation on rational points over
a specified finite field: the geometric intersection points
in the proof may only appear after extension, when the
rational exponents themselves can increase.
\end{claim}
\begin{proof}
If $\phi$ already identifies two geometric points, so does
$u\phi$. Suppose therefore that $\phi$ is injective on
geometric points. It is nonconstant and finite onto its
irreducible image curve $Z\subset A$. Its function-field
extension over the normalization $\widetilde Z$ is purely
inseparable: a nontrivial separable degree would give
multiple geometric points in a general fibre. A finite
purely inseparable map between smooth curves over a perfect
field is a power of relative Frobenius followed by an
isomorphism. Hence
\[
                g(\widetilde Z)=g(D)\ge2.
\]
The image may be singular, but its arithmetic genus is
at least this normalization genus. Adjunction on the smooth
abelian surface, whose canonical bundle is trivial, gives
\[
                   Z^2=2p_a(Z)-2>0.
\]
In fact $Z$ is ample. Every irreducible curve $Y$ on an
abelian surface has $Y^2\ge0$, by intersecting it with a
distinct translate. If $Z\cdot Y=0$, the Hodge index
theorem and $Z^2>0$ would force $Y^2<0$, since $Y$
is not numerically zero (an ample divisor has positive
intersection with $Y$). This is impossible. Thus $Z$
meets every curve positively, and the surface
Nakai--Moishezon criterion proves ampleness. Only the
already established inequality $Z^2>0$ is needed below.

Choose $0\ne t\in\ker u(\Omega)$. If $Z=Z-t$,
any point $z\in Z$ gives two distinct points $z,z+t$
on $Z$. Otherwise the two curves $Z$ and $Z-t$ have
intersection number $Z^2>0$, because translations preserve
numerical equivalence. They therefore meet; a point in
their intersection again gives distinct $z,z+t\in Z$.
Choose preimages of these two points in $D$. Their images
under $u\phi$ agree, contradicting injectivity.

For ordinary elliptic curves in characteristic two,
$E_i[2](\Omega)$ is nontrivial, so the displayed product
isogeny has a nonzero geometric kernel. Moreover the unique
nonzero geometric two-torsion point of each ordinary
elliptic curve over a finite field is rational; hence its
rational two-primary exponent is always at least two,
as noted in AG.13. The general result applies.

Finally, merely knowing that two unmultiplied maps
$\phi_i$ separate points imposes no new torsion constraint
on a hypothetical collision. Once
$T=[Q'-P']$ is $p$-primary, every induced Jacobian
homomorphism automatically sends $T$ to a $p$-primary
point of its target. If target exponents $b_i$ are used,
it is the maps $[b_i]\phi_i$, rather than the original
$\phi_i$, whose point separation would contradict the
collision. The theorem explains exactly why this geometric
replacement fails for two ordinary elliptic targets.
\end{proof}


\begin{claim}{AG.31}{Proved; a limit of normal-chain tests with nonconstant deck Frobenius}
There is a finite odd solvable group $G$, a finite
$\mathbf F_2[G]$-module $V$, and an automorphism $F$ of $V$
semilinear for an automorphism of $G$, such that
\[
 (V^H)^F\supsetneq(V^K)^F
 \quad\text{for every proper inclusion }H<K
       \text{ of normal subgroups of }G,
\]
but no $F$-fixed vector has trivial $G$-stabilizer.
Thus the constant-deck simultaneous-vector assertion cannot
be extended to nonconstant Frobenius using only these
normal-chain tests, even if they are imposed over every
finite constant-field extension.
This is an abstract module counterexample, without a claim
that the module is realized by a curve or its Jacobian.
\end{claim}
\begin{proof}
Let $N=(\mathbf F_{25},+)$ and let $c$ of order three act
on $N$ by multiplication by $\omega$, where
$\omega^2+\omega+1=0$ over $\mathbf F_5$.
The polynomial is irreducible over $\mathbf F_5$,
so $N$ is an irreducible two-dimensional $\mathbf F_5$
module for $\langle c\rangle$. Put
\[
             G=N\rtimes\langle c\rangle,
             \qquad |G|=75.
\]
Its normal subgroups are exactly $1,N,G$.
Indeed, for a normal subgroup $L$, the intersection
$L\cap N$ is either $1$ or $N$ by irreducibility.
In the first case $[L,N]\subset L\cap N=1$,
so $L\subset C_G(N)=N$ and $L=1$.
In the second case the quotient $G/N$ has prime order,
so $L=N$ or $G$.

Take the left regular module $V=\mathbf F_2[G]$,
choose $0\ne t\in N$, and set $F=\rho(t)$,
where $\rho$ denotes the left regular representation.
Then
\[
 F\rho(g)F^{-1}=\rho(tgt^{-1}),
\]
so $F$ is semilinear for the inner deck automorphism
$g\mapsto tgt^{-1}$. For every subgroup $H$,
\[
 (V^H)^F=V^{\langle H,t\rangle},\qquad
 \dim_{\mathbf F_2}(V^H)^F
       =\frac{75}{|\langle H,t\rangle|}.
\]
The dimensions for $H=1,N,G$ are respectively
$15,3,1$. Every required normal-chain inclusion
is therefore strict. This remains true for every power $F^a$:
the dimensions are $15,3,1$ when $5\nmid a$, and $75,3,1$
when $5\mid a$. Nevertheless every vector in
$V^F$ is fixed by the nonidentity element $t$.
The indicator vector of the subgroup $\langle t\rangle$
has stabilizer exactly that subgroup, so the minimum
stabilizer order on $V^F$ is exactly five.

The missing intermediate test is explicit:
$1\triangleleft\langle t\rangle$ is stable under the
deck Frobenius, and
\[
            (V^1)^F=(V^{\langle t\rangle})^F.
\]
Here $\langle t\rangle$ is not normal in $G$.
There is also a missing test with a $G$-normal upper group:
if $A_0\subset N$ is an $\mathbf F_5$-line not containing $t$,
then $A_0\triangleleft N$, both are $F$-stable, and
$\dim(V^{A_0})^F=\dim(V^N)^F=3$.
Consequently a hypothesis imposed on all relevant
intermediate relative covers can detect this obstruction,
even though the normal-subgroup-only hypothesis cannot.
The standard-library certificate
\path{code/nonconstant_normal_chain_counterexample.py}
enumerates the complete normal-subgroup lattice,
checks the semilinear identity and all fixed dimensions,
and computes the displayed vector's exact stabilizer.
Its full output is retained in the companion JSON file.
\end{proof}


% AG.32--AG.33: full direct proof; independently approved by root and criteria.
\begin{claim}{AG.32}{Proved; an effective conductor bound in the scalar supersingular case}
Let $C/\mathbf F_{2^{2r}}$ be a genus-two curve with two-rank zero,
with rational Weierstrass point $P$, and suppose Jacobian Frobenius
is $[\varepsilon2^r]$, where $\varepsilon\in\{1,-1\}$.
If a nontrivial connected cyclic \'etale cover of odd degree $n$
has primitive part of two-rank zero, put
\[
 f=\operatorname{ord}_n(2),\qquad d=\gcd(f,r),\qquad
 M=\frac{\varphi(n)}d.
\]
Then
\[
             M(M+1)\le 2(n-1),\qquad n<10r^3.       \tag{AG.32.1}
\]
Consequently the geometric primitive-cover question for this
fixed scalar-Frobenius curve has only finitely many possible
conductors, with the displayed explicit bound.
\end{claim}
\begin{proof}
The Raynaud theta divisor is the symmetric Abel curve
$\Theta=C-P$, and contains zero. In characteristic two its
nontrivial character cohomology has dimension $g-1=1$.
The orbit criterion AG.3 therefore says that primitive
rank zero requires at least one theta point in every
multiplication-by-two orbit on the primitive characters
of the cyclic subgroup of order $n$.
There are $\varphi(n)/f$ such orbits, each of size $f$.

Frobenius preserves $\Theta$, and symmetry removes the sign
$\varepsilon$. Thus multiplication by $2^r$ preserves its
geometric points. A theta point in one of the required
orbits brings with it its entire multiplication-by-$2^r$
orbit, of size $f/d$. The primitive theta intersection
therefore has at least $M=\varphi(n)/d$ points.
Together with zero there are at least $M+1$ points of
$\Theta$ in the cyclic subgroup. For any nonzero difference
there are at most two ordered pairs of distinct points of
a genus-two Abel curve, since two distinct translates of
its principal theta divisor have intersection number two.
The exact difference count used in AG.27 now gives
$M(M+1)\le2(n-1)$.

For completeness, the elementary uniform estimate
\[
                  \varphi(n)\ge\frac23 n^{5/6}
                       \quad(n\text{ odd})             \tag{AG.32.2}
\]
follows prime-power by prime-power. For primes $\ell\ge5$,
$(\ell-1)^6\ge\ell^5$: it holds at five, and
$(x-1)^6/x^5$ increases for $x>1$. Hence
$\varphi(\ell^a)^6/\ell^{5a}\ge1$ for every $a\ge1$.
The only remaining prime, three, contributes a factor at
least $64/243$, which is larger than $(2/3)^6$.
Multiplying proves (AG.32.2). Since $d\le r$, the first
inequality of (AG.32.1) implies
\[
 \frac{4}{9r^2}n^{5/3}\le M^2<2n,
 \qquad
 n<\left(\frac92\right)^{3/2}r^3<10r^3.
\]
The strict inequality is preserved because $n>1$.
\end{proof}

\begin{claim}{AG.33}{Refuted: universal nontrivial primitive-cover existence}
\label{clm:AGprimitivecounterexample}
There are supersingular genus-two curves over
$\mathbf F_{2^{128}}$ for which every nontrivial connected
cyclic \'etale cover of odd degree has primitive part of
positive two-rank. Thus the version of the existence statement
A6 requiring degree $n>1$ is false, already for genus two
in characteristic two.

A completely specified family containing these examples is
\[
 C_t:\qquad y^2+y=
       \left(\frac{u^2+t^2u+t}{t^3+1}\right)^3,
 \qquad
 t\in\mathbf F_{2^{128}}\setminus\mathbf F_4,       \tag{AG.33.1}
\]
where $C_t$ means the smooth proper model. The proof gives
an explicit finite incidence set of excluded parameters
and a positive exact lower bound for its complement.
The program certifies this nonemptiness and the displayed
algebraic identities; it does not claim to have enumerated
one particular parameter in the complement.

All these curves themselves satisfy A1: their Jacobians
have geometric two-rank zero. The refutation concerns
existence of nontrivial primitive-zero covers, and gives no
counterexample to A1 or to the arithmetic criterion A3.
The prompt does not explicitly require $n>1$. With its
convention $P_1=J_C$, the identity cover is a witness for
every curve constructed here. Therefore these examples
do not refute the literal degree-one-permitted version
of A6; that version remains open for curves of positive
$p$-rank.
\end{claim}
\begin{proof}
We first establish the family and its torsion incidence
bound directly. It is the characteristic-two pencil
studied by Moret--Bailly in \emph{Polarisations de degr\'e 4
sur les surfaces ab\'eliennes}, the theorem, parts~1(c) and~3
\citep{AG_MoretBailly1979}; the explicit calculations below
supply the required properties independently.

\emph{1. Explicit genus-two curves and elliptic maps.}
Let $k$ be a finite field containing $\mathbf F_4$, put
$E:y^2+y=x^3$, $A=E\times E$, and take
$t\in k\setminus\mathbf F_4$. Write $d=t^3+1$, so
$t\ne0$ and $d\ne0$. If $X=u^2+t^2u+t$, expansion gives
\[
 X^3=u^6+t^2u^5+(t+t^4)u^4+t^6u^3
                +(t^5+t^2)u^2+t^4u+t^3.
\]
Over the perfect field $k$, successive Artin--Schreier
changes remove the even powers $u^6,u^4,u^2$. The highest
remaining odd power is $u^5$, with nonzero coefficient
$t^2/d^3$. Thus (AG.33.1) has genus two and two-rank zero,
and has exactly one point $P$ above infinity; that point
is $k$-rational and is its unique Weierstrass point.
In particular $\omega_C\simeq\mathcal O_C(2P)$.

Define two morphisms to $E$ by
\[
\begin{aligned}
 x_1&=\frac{u^2+t^2u+t}{d}, & y_1&=y,\\
 x_2&=\frac{tu^2+u+t^2}{d}, & y_2&=y+\frac{u^3+1}{d}.
\end{aligned}                                                   \tag{AG.33.2}
\]
The first elliptic equation is the defining equation of
$C_t$. For the second, put $Z=tu^2+u+t^2$ and $Y=u^3+1$.
The polynomial identity
\[
                 X^3+Z^3=dY^2+d^2Y
\]
proves $y_2^2+y_2=x_2^3$. Properness extends both maps
across $P$, and both take $P$ to the elliptic origin.
The joint map $i_t:C_t\to A$ is birational onto its image,
because
\[
                         u=x_2+t x_1,
                         \qquad y=y_1.                \tag{AG.33.3}
\]

\emph{2. The pencil and its only geometric base point.}
On $A$ consider
\[
 f=\frac{y_2+y_1+1+x_1x_2}{x_2+x_1^2},\qquad
 D_\infty=E\times\{0\}+\Gamma_{F_E},                 \tag{AG.33.4}
\]
where $F_E:E\to E$ is the degree-two Frobenius morphism.
The polar divisor of $f$ is exactly $D_\infty$.
Here is the local check. In the affine part, the denominator
vanishes on the two graphs $\Gamma_{F_E}$ and
$\Gamma_{-F_E}$. On the first graph its numerator is one,
so the pole is simple. On the second the numerator
vanishes and the factor cancels. Along $E\times\{0\}$,
the orders of $y_2$ and $x_2$ are minus three and minus two,
so again the pole is simple with nonzero leading coefficient.
Along $\{0\}\times E$, the numerator has order minus three
and the denominator minus four, so $f$ vanishes rather
than having a pole. These account for every possible
polar divisor.

The two components of $D_\infty$ meet geometrically only
at the origin: the kernel of $F_E$ has only its
identity as a geometric point. At every nonzero point of
the components of $D_\infty$, the preceding nonzero pole
coefficient shows that the section $f$ of
$\mathcal O_A(D_\infty)$ does not vanish. Hence the pencil
generated by $1$ and $f$ has no nonzero geometric base point.
It has the origin as a base point, because (AG.33.2) sends
$P$ to zero for every admissible $t$.

Substitution in (AG.33.4) gives $f\circ i_t=t$.
The divisor $D_t=(f-t)+D_\infty$ is precisely the image
$i_t(C_t)$, with multiplicity one. To check that no
component is being suppressed, write $x=x_1,z=x_2$.
Eliminating $y_2$ from the numerator equation for $f=t$
and the two elliptic equations gives
\[
 (z+x^2)\bigl(z^2+t^2z+t^2x^2+x+t\bigr)=0.
\]
The first factor is the cancelled graph. The second is
exactly (AG.33.1) after $u=z+tx$, as (AG.33.2)--(AG.33.3)
show. It is irreducible because its smooth proper model
has the Artin--Schreier equation with nonzero reduced
$u^5$ coefficient. It has multiplicity one: away from
the denominator, differentiating on $A$ gives
\[
                   df=dx+
                  \frac{z^2+x+f}{z+x^2}\,dz,
\]
which is nonzero along the generic point of this component.
The boundary component $\{0\}\times E$ is a zero fibre
only for $t=0$, excluded here, and the other boundary
component is polar. Therefore $D_t\sim D_\infty$.

\emph{3. The induced isogeny and the theta origin.}
The elliptic components in $D_\infty$ have self-intersection
zero and intersect with length $\deg F_E=2$. Consequently
\[
                         D_t^2=4.                     \tag{AG.33.5}
\]
The map $i_t$, based at $P$, induces a homomorphism
$v_t:J_{C_t}\to A$. Its image is all of $A$: otherwise
the genus-two image curve would lie in a translate of an
elliptic subvariety and would itself have normalization
genus one. Thus $v_t$ is an isogeny.
For the principal Abel theta divisor $\Theta=C_t-P$,
birationality gives $(v_t)_*\Theta=D_t$.
For an isogeny of abelian surfaces, pullback and pushforward
on the rational N\'eron--Severi group satisfy
$v_t^*(v_t)_*=(\deg v_t)\operatorname{id}$, and pullback
multiplies intersection numbers by $\deg v_t$.
Using $\Theta^2=2$ therefore gives
\[
                       D_t^2=2\deg v_t.
\]
Equation (AG.33.5) proves $\deg v_t=2$.
In particular it induces an isomorphism on geometric
$n$-torsion for every odd $n$.

There is no unspecified translation of Raynaud theta in
this argument. Since $C_t$ has two-rank zero,
$J_{C_t}[2](\overline k)=0$. Its theta characteristic is
therefore unique; as $\omega_C=\mathcal O_C(2P)$, the
Raynaud theta characteristic is $\mathcal O_C(P)$.
The corresponding theta divisor is exactly $C_t-P$.
Thus a nonzero theta point killed by an odd $n$ maps to
a nonzero point of $A[n]$ on $D_t$.

\emph{4. Uniform torsion incidence and scalar Frobenius.}
Every nonzero geometric point of $A$ belongs to at most
one member of the pencil, by its single base point.
Since $\#A[n](\overline k)=n^4$ for odd $n$, there are
at most $n^4-1$ parameters over $\overline{\mathbf F}_2$
for which $\Theta_{C_t}$ contains a nonzero $n$-torsion point.
This bounds all geometric torsion points, regardless of
which extension makes them rational.

For $k=\mathbf F_{4^e}$, Frobenius on $E$ is $[(-2)^e]$:
indeed $\#E(\mathbf F_2)=3$, so its Frobenius
polynomial is $X^2+2$, and Cayley--Hamilton for the
endomorphism gives $F_E^2=[-2]$.
The isogeny $v_t$ is defined over $k$, so the same exact
scalar formula holds on $J_{C_t}$. The conductor bound
AG.32 applies with $r=e$. If a primitive-zero cyclic cover
exists, its odd conductor therefore satisfies $3\le n<10e^3$,
and its theta criterion supplies a nonzero theta point
of exact order $n$.

\emph{5. A finite, nonempty counterexample set.}
Take $e=64$, $q=4^{64}=2^{128}$, and $B=10\cdot64^3=2621440$.
Define a finite subset of $\mathbf F_q$ by
\[
 \mathcal B=\mathbf F_4\ \cup\
 \left(\mathbf F_q\cap
 \bigcup_{\substack{3\le n<B\\ n\text{ odd}}}
 \{f(a):a\in A[n](\overline{\mathbf F}_2)\setminus\{0\},
                         \ f(a)\text{ finite and defined}\}\right).
                                                               \tag{AG.33.6}
\]
Here $f(a)$ means the value of the pencil morphism
$A\setminus\{0\}\to\mathbf P^1$, including extension across
removable zeros of the displayed numerator and denominator.
This is an exact algebraic definition using the displayed
rational function on the fixed elliptic surface $E\times E$.
It specifies a finite set, without an assertion that the
large set has been enumerated. The incidence bound gives
\[
 \#\mathcal B\le4+
       \sum_{\substack{3\le n<B\\ n\text{ odd}}}(n^4-1)<4+B^5<q.
\]
Every $t\in\mathbf F_q\setminus\mathcal B$ defines a
smooth curve (AG.33.1), and no odd conductor allowed by
AG.32 can supply a primitive-zero cover. Larger conductors
are impossible by that same bound. This proves the
asserted counterexample existence over the stated finite
field, with no restriction on the field of definition of
the hypothetical cover.

The exact counting certificate actually gives
\[
 \#(\mathbf F_q\setminus\mathcal B)\ \ge\
       340282354541538070615576657943687331836>0.
\]
The standard-library program
\path{code/supersingular_primitive_counterexample.py}
verifies the polynomial identities of (AG.33.2)--(AG.33.4),
the Artin--Schreier coefficient expansion, the elementary
integer bounds, and the exact odd fourth-power sum.
Its complete output is retained in the companion JSON.
The program is a finite algebra and nonemptiness certificate,
not an enumeration of an individual parameter or a test
of infinitely many covers.

Finally, $J_{C_t}$ is isogenous to $E^2$ and hence has
geometric two-rank zero. R1 therefore already excludes a
collision on each $C_t$. The two supersingular elliptic
maps (AG.33.2) jointly separate geometric points as well:
$v_t$ has degree two and a connected kernel, since
$J_{C_t}[2](\overline k)=0$. It is therefore radicial on
geometric points, while the Abel map is an
embedding. This is consistent with AG.30, whose hypothesis
requires a nonzero geometric kernel point.
\end{proof}

\begin{claim}{AG.34}{Proved; nontrivial primitive-cover counterexamples over $\mathbf F_{2^{38}}$}
For the family $C_t$ of AG.33 over $\mathbf F_{2^{38}}$,
at least $59927733332$ parameters $t$ define curves having
no nontrivial connected cyclic \'etale cover of odd degree
with primitive two-rank zero. This counts parameters,
without asserting that the corresponding curves are
pairwise nonisomorphic or selecting one individual parameter.
The degree-one qualification at the end of AG.33 applies.
\end{claim}
\begin{proof}
There is a useful refinement of the torsion incidence bound.
For an integer $e\ge1$ let
\[
\begin{split}
 \mathcal S_e={}&\{n:\ 3\le n<10e^3,\ n\text{ odd},\
       M_e(n)(M_e(n)+1)\le2(n-1)\},\\
 M_e(n)={}&\frac{\varphi(n)}{\gcd(\operatorname{ord}_n(2),e)},\qquad
 J_4(n)=n^4\prod_{\ell\mid n}(1-\ell^{-4}).
\end{split}                                                   \tag{AG.34.1}
\]
The integer $J_4(n)$ is the number of points of exact order
$n$ in $A[n]\simeq(\mathbf Z/n\mathbf Z)^4$.
For $n\in\mathcal S_e$ define $\mathcal B_{e,n}$ to be
the parameters $t\in\mathbf F_{4^e}\setminus\mathbf F_4$
for which the pencil member $D_t$ contains at least
$M_e(n)$ distinct points of exact order $n$ in $A$.
Every nonzero point belongs to at most one pencil member.
Counting incidences therefore gives
\[
               \#\mathcal B_{e,n}
                  \le\left\lfloor\frac{J_4(n)}{M_e(n)}\right\rfloor.
                                                               \tag{AG.34.2}
\]
If $C_t$ admits a nontrivial primitive-zero cover of
conductor $n$, AG.32 puts $n$ in $\mathcal S_e$ and gives
at least $M_e(n)$ exact-order-$n$ points of its theta
divisor. The odd-torsion isomorphism induced by $v_t$
sends them injectively to points counted in
$\mathcal B_{e,n}$. Thus the explicitly defined set
\[
 \mathbf F_{4^e}\setminus
       \left(\mathbf F_4\cup
          \bigcup_{n\in\mathcal S_e}\mathcal B_{e,n}\right)
                                                               \tag{AG.34.3}
\]
consists entirely of parameters with the asserted property,
and its size is at least
\[
       4^e-4-
       \sum_{n\in\mathcal S_e}
          \left\lfloor\frac{J_4(n)}{M_e(n)}\right\rfloor.
                                                               \tag{AG.34.4}
\]
The definition of $\mathcal B_{e,n}$ uses the threshold
$M_e(n)$; the stronger bound (AG.34.2) is not being asserted
for parameters with only one theta point of order $n$.

For $e=19$, exact enumeration below the proved global
cutoff $68590$ gives precisely
\[
 \mathcal S_{19}=\{191,229,361,419,457,571,573,647,687,
                         955,1083,1257,1371\}.
\]
The certificate records all factorizations, totients,
multiplicative orders, values of $M_{19}(n)$, exact-order
point counts, and integer quotients. Its sum is
\[
 \sum_{n\in\mathcal S_{19}}
      \left\lfloor\frac{J_4(n)}{M_{19}(n)}\right\rfloor
       =214950173608.
\]
As $4^{19}=274877906944$, formula (AG.34.4) becomes
\[
             274877906944-4-214950173608
                     =59927733332>0.
\]
The root program \path{code/A6_scalar_conductor_bounds.py}
and its complete JSON output give the global enumeration.
The independent review
\path{agent_geometric/code/primitive_counterexample_e19_review.py}
factors every odd integer below the cutoff by trial division
and recomputes every relevant multiplicative order without
using the root program's sieve or cache. Its complete JSON
agrees on every $e=19$ record. These are exhaustive finite
arithmetic certificates for the displayed existence theorem,
not an enumeration of a parameter in (AG.34.3).
\end{proof}

\begin{claim}{AG.35}{Proved: geometric origin alone does not bound unit-root returns}
Consider the fixed smooth elliptic family over the multiplicative group
\[
 \mathcal E\longrightarrow\mathbf G_{m,\mathbf F_4},
 \qquad E_a:\ y^2+xy=x^3+a.                         \tag{AG.35.1}
\]
For every positive odd integer $m$, there is a closed parameter $a$
of exact degree $m$ over $\mathbf F_4$ such that
\[
 \#E_a(\mathbf F_{4^m})=4^m,
 \qquad E_a(\mathbf F_{4^m})\simeq\mathbf Z/2^{2m}\mathbf Z.
                                                        \tag{AG.35.2}
\]
If $\beta_a\in\mathbf Z_2^\times$ is its unit-root Frobenius
eigenvalue, then
\[
 \beta_a^2-\beta_a+4^m=0,
 \qquad v_2(\beta_a-1)=2m.                         \tag{AG.35.3}
\]
The parameter is a point of odd order in
$\mathbf G_m(\overline{\mathbf F}_2)$, and its Teichmuller lift
is prime-to-two torsion. Thus a single fixed geometric family,
ordinary fibers, integral cohomology of Weil weight one, and
prime-to-two Frobenius periods do not give a uniform bound for
the return unit minus one, even at prime-to-two torsion parameters.
The assertion does not put these parameters in the powers of one
fixed odd prime, and does not identify this family with the Abel
tower of AR.20. Those restrictions remain essential unresolved
inputs to that route.
\end{claim}
\begin{proof}
The discriminant and $j$-invariant of (AG.35.1) are $a$ and $1/a$,
respectively. Hence the displayed projective Weierstrass model is
a smooth elliptic scheme over $\mathbf G_m$. Each fiber is ordinary:
the geometric point $(0,\sqrt a)$ is a nonzero point of order two,
whereas an elliptic curve has at most one independent geometric
two-torsion point in characteristic two.

Fix odd $m$ and put $q=4^m$. Waterhouse's classification
\citep[Theorem~4.1, case~(1), p.~536]{AG_Waterhouse1969}
asserts that every integer $t$ with $|t|\le2\sqrt q$ and
$(t,2)=1$ is the Frobenius trace of an ordinary elliptic curve
over $\mathbf F_q$. Apply it to $t=1$. The resulting elliptic
curve $E/\mathbf F_q$ has $q$ rational points.
Since $4\mid q$, Proposition~5.1 of
\citep[p.~311]{AG_AuerTop2002} identifies $E$ over $\mathbf F_q$
with a member $E_a$ of (AG.35.1).

For completeness, the normal-form step can be seen directly.
An ordinary elliptic curve over the perfect field $\mathbf F_q$
has an equation
\[
               y^2+xy=x^3+b x^2+a,\qquad a\ne0.
\]
Its unique nonzero geometric two-torsion point is
$T=(0,\sqrt a)$. As $E(\mathbf F_q)$ is a two-group and
the geometric two-primary torsion has rank one, it is cyclic;
since its order is $q\ge4$, the point $T$ has a rational half.
The duplication formula for a point $R$ with nonzero
$x$-coordinate is
\[
                 x(2R)=x(R)^2+\frac a{x(R)^2}.
\]
Thus a half of $T$ must have $x$-coordinate $r=a^{1/4}$.
Writing its $y$-coordinate as $rz$, the equation becomes
\[
                       z^2+z=b+r+r^2.
\]
It is solvable over $\mathbf F_q$ precisely when
$\operatorname{Tr}_{\mathbf F_q/\mathbf F_2}(b)=0$.
Consequently $b=s^2+s$ for some $s\in\mathbf F_q$, and
the substitution $y\mapsto y+sx$ removes the $x^2$ term.
This proves the required isomorphism to $E_a$ and also confirms
the rational four-torsion step in the cited proposition.

The characteristic polynomial of $E_a$ is $X^2-X+q$.
Modulo two it has distinct roots zero and one. Its unit root
$\beta_a$ therefore belongs to $\mathbf Z_2^\times$, and the
other root is $q/\beta_a$, of valuation $2m$. Evaluating at one
gives
\[
              (1-\beta_a)(1-q/\beta_a)=q.
\]
The second factor is a two-adic unit, proving (AG.35.3).
The geometric two-primary torsion of an ordinary elliptic curve
is $\mathbf Q_2/\mathbf Z_2$, so its finite rational two-primary
subgroup is cyclic. Its full rational order here is $q$, which
proves the group assertion in (AG.35.2).

It remains to prove that the parameter has exact degree $m$;
the existence of a curve over $\mathbf F_{4^m}$ alone would not
establish the assertion about periods. Let
$d=[\mathbf F_4(a):\mathbf F_4]$, so $d\mid m$. The displayed
model $E_a$ is defined over $\mathbf F_{4^d}$. Write
$\gamma\in\mathbf Z_2^\times$ for its unit root over that field.
The unit root over $\mathbf F_{4^m}$ is
$\gamma^{m/d}$. Because $m/d$ is odd,
\[
 v_2(1-\gamma^{m/d})=v_2(1-\gamma);
\]
the quotient $(1-\gamma^{m/d})/(1-\gamma)$ is a sum of an
odd number of odd two-adic integers. The nonunit Frobenius root
does not contribute to the two-primary part of the rational
point count, so
\[
             v_2\bigl(\#E_a(\mathbf F_{4^d})\bigr)=2m.
\]
If $d<m$, Hasse's bound contradicts this divisibility:
\[
 4^m\le\#E_a(\mathbf F_{4^d})
       \le(2^d+1)^2<4^m.
\]
Hence $d=m$. Finally $a\in\mathbf F_{4^m}^\times$ has order
dividing the odd number $4^m-1$, and its fourth-power Frobenius
orbit has length exactly $m$. Teichmuller lifting preserves that
order and that orbit length.

The small exact certificate
\path{agent_geometric/code/geometric_unitroot_growth.py}
constructs examples for $m=1,3,5$, records the finite fields,
the parameter, its order and degree, and the point count, and
checks the unit-root congruence. This finite check illustrates
the theorem; its assertion for every odd $m$ follows from the
classification and the degree argument above.
\end{proof}

% AG.36: root and criteria independently audited.
\subsection{A geometric unit-root obstruction along a fixed prime-power sequence}

\begin{claim}{AG.36}{Proved}\label{clm:AGfixedthreepower}
For each \(n\geq1\), choose a primitive \(3^n\)-th root
\(a_n\in\overline{\mathbb F}_2\), and put
\[
 d_n=2\cdot3^{n-1},\qquad k_n=\mathbb F_{2^{d_n}},\qquad
 E_n:\ y^2+xy=x^3+a_n.
 \tag{AG.36.1}
\]
Let \(\beta_n\in\mathbb Z_2^\times\) be the unit root of
\(X^2-t_nX+2^{d_n}\), where \(t_n\) is the Frobenius trace
of \(E_n/k_n\). Then
\[
 \lim_{n\to\infty}v_2(\beta_n-1)=+\infty,
 \qquad
 \exp E_n(k_n)[2^\infty]=2^{v_2(\beta_n-1)}.
 \tag{AG.36.2}
\]
Every \(E_n\) is a fibre of the single smooth ordinary elliptic
family \(y^2+xy=x^3+a\) over \(\mathbb G_{m,\mathbb F_2}\).
The parameter \(a_n\) has exact degree \(3^{n-1}\) over
\(\mathbb F_4\), an odd degree, and has exact multiplicative
order \(3^n\). Thus a fixed geometric family, weight-one Weil
purity, integrality, ordinariness, a fixed prime-power torsion
sequence, and bounded \(2\)-part of the parameter degrees do
not imply a uniform bound for the return valuations.

The conclusion asserts divergence, without a quantitative
comparison between \(2^{v_2(\beta_n-1)}\) and a prescribed
power of \(3^n\). It neither identifies this family with the
relative factors of an Abel tower nor disproves the weaker
polynomial growth bound required by that route.
\end{claim}

\begin{proof}
We first prove a finite-cover statement. If \(Y\to
\mathbb A^1_{\mathbb F_2}\) is a finite etale torsor under a
finite abelian \(2\)-group and its fibre at \(0\) is split,
then its fibre at every primitive \(3^n\)-th root is split over
\(k_n\) for all sufficiently large \(n\).
It is enough to prove this for a cyclic quotient of order
\(2^r\), and then take the maximum of finitely many bounds.
The Artin--Schreier--Witt exact sequence and affine vanishing
of \(H^1(\mathbb A^1,W_r\mathcal O)\) give
\[
 H^1_{\mathrm{\acute et}}(\mathbb A^1_{\mathbb F_2},
                         \mathbb Z/2^r)
 =W_r(\mathbb F_2[a])/(F-1)W_r(\mathbb F_2[a]).
 \tag{AG.36.3}
\]
Here \(F\) is Witt Frobenius. These standard exact statements
are recorded explicitly in \citet[Theorem~1.2 and the following
cohomology sequence, p.~2]{AG_Sekiguchi2001}; compare the
field-theoretic treatment and Frobenius trace formula in
\citet[Theorem~3.2, Section~5.4, p.~29]{AG_KostersWan2016}.
Affine vanishing also follows directly by induction from the
filtration of \(W_r\mathcal O\) with additive-group quotients.

Choose a representative \(w(a)\) in (AG.36.3).
On \(W_r(\mathbb F_2)\), Frobenius is the identity. Since
splitting at \(0\) says that the specialized cohomology class
vanishes, it follows that \(w(0)=0\), as an actual Witt vector.
There is a finite additive expansion
\[
                 w(a)=\sum_{i=0}^{r-1}\sum_{j\in J_i}
                                  V^i[a^j],
                 \qquad J_i\subset\mathbb Z_{>0}\text{ finite}.
 \tag{AG.36.4}
\]
Repeated terms may be retained in this expression. To obtain
it, subtract the sum of the Teichmuller monomials of the first
Witt coordinate. The difference has first coordinate zero,
so is uniquely \(V\) of a shorter Witt vector with polynomial
coordinates. This shorter vector still specializes to zero
at \(a=0\). Induction proves (AG.36.4), with no constant
monomial and only finitely many exponents. Let \(D\) be the
largest exponent occurring, taking \(D=0\) if the class is
zero.

The order of \(2\) modulo \(3^n\) is \(d_n\): the order
is even and \(v_3(4^m-1)=1+v_3(m)\). Hence the absolute
Frobenius conjugates of \(a_n\) are precisely all primitive
\(3^n\)-th roots. In the full Witt ring of \(k_n\), the
Teichmuller lift \(z=[a_n]\) has exact order \(3^n\), and
\[
 \sum_{s=0}^{d_n-1}z^{j2^s}
 =\sum_{\zeta\text{ primitive of order }3^n}\zeta^j=0
 \qquad(0<j<3^{n-1}).
 \tag{AG.36.5}
\]
Indeed the last sum is the sum over all \(3^n\)-th roots
minus the sum over all \(3^{n-1}\)-th roots; both geometric
sums vanish for the stated exponents. Since \(F\) commutes
with \(V\), (AG.36.4)--(AG.36.5) show that
\[
 \operatorname{Tr}_{W_r(k_n)/W_r(\mathbb F_2)}w(a_n)=0
                       \qquad\text{if }3^{n-1}>D.
 \tag{AG.36.6}
\]
For \(F(b)-b=w(a_n)\), the action of the degree-\(d_n\)
Frobenius on \(b\) is translation by this trace, by telescoping
\(F^{d_n}(b)-b\). Thus (AG.36.6) is exactly the splitting
condition. This proves the finite-cover statement.

We now apply it to the elliptic family. Its discriminant is
\(a\); each fibre over \(\mathbb G_m\) is ordinary and
has geometric etale \(2\)-rank one. For each \(r\geq1\),
let
\[
 Q_r=\ker\bigl(V^r:E^{(2^r)}\longrightarrow E\bigr).
 \tag{AG.36.7}
\]
This is a finite etale group scheme over \(\mathbb G_m\),
with geometric fibres cyclic of order \(2^r\). Its frame
torsor has group \((\mathbb Z/2^r)^\times\), a finite
abelian \(2\)-group. We show directly that this torsor is
trivial over the local field \(K=\mathbb F_2((a))\).
This statement concerns the etale group (AG.36.7), and does
not presume rationality of the non-etale group scheme
\(E[2^r]\) over the imperfect field \(K\).

In the purely inseparable extension \(K^{1/2}\), the point
\(R_1=(0,a^{1/2})\) has exact order two. Suppose a point
\(R=(x,y)\) of order \(2^s\) is defined over
\(K^{1/2^s}\), and both coordinates have positive valuation
(the zero coordinate is allowed). The convergent series
\[
                  \lambda=\sum_{i\geq0}x^{2^i}
 \tag{AG.36.8}
\]
belongs to the same complete local field, has positive
valuation, and satisfies \(\lambda^2+\lambda=x\).
Set
\[
 u^2=y+(\lambda+1)x,\qquad v=u(\lambda+u).
 \tag{AG.36.9}
\]
These coordinates belong to \(K^{1/2^{s+1}}\) and have
positive valuation. The curve equation for \(R\) gives
\(u^4+xu^2=a\); in particular \(u\ne0\).
It follows directly that \(v^2+uv=u^3+a\), and the
elliptic doubling formulas give \(2(u,v)=R\).
Thus \((u,v)\) has exact order \(2^{s+1}\).
Induction constructs a generator of the geometric
\(2^r\)-torsion in \(E(K^{1/2^r})\).
Applying relative Frobenius \(F^r\) gives a \(K\)-rational
generator of \(Q_r\). Consequently \(Q_r/K\) is constant.

The frame torsor is therefore unramified at \(a=0\),
with trivial residue Frobenius there. Its normalization over
\(\mathbb A^1_{\mathbb F_2}\) is finite etale and has split
fibre at \(0\). Apply the finite-cover statement to this
extension. For each fixed \(r\), and all sufficiently large
\(n\), the group \(Q_r\) at \(a_n\) is constant over
\(k_n\). Relative Frobenius identifies its geometric points
Galois-equivariantly with those of \(E_n[2^r]\), so the
unit root satisfies \(\beta_n\equiv1\pmod{2^r}\).
This proves the limit in (AG.36.2).

Finally, the rational \(2\)-primary subgroup of an ordinary
elliptic curve over a finite field is cyclic. The other
Frobenius root \(2^{d_n}/\beta_n\) is divisible by two,
and therefore
\[
 v_2\#E_n(k_n)
 =v_2\bigl((1-\beta_n)
              (1-2^{d_n}/\beta_n)\bigr)
 =v_2(1-\beta_n).
\]
This proves the exponent formula. The exact parameter degree
over \(\mathbb F_4\) is \(\operatorname{ord}_{3^n}(4)
=3^{n-1}\), as asserted.
\end{proof}

\paragraph{Finite certificates and quantitative scope.}
The independent complete point-count program
\texttt{agent\_geometric/code/kloosterman\_three\_power\_test.cpp}
and its JSON output use the irreducible polynomial
\(\Phi_{3^n}(X)=X^{2\cdot3^{n-1}}+X^{3^{n-1}}+1\)
for \(n=1,2,3\). They verify, respectively,
\[
\begin{array}{c|r|r|r|r}
 n&\#k_n&\#E_n(k_n)&t_n&v_2(\beta_n-1)\\\hline
 1&4&4&1&2\\
 2&64&64&1&6\\
 3&262144&262912&-767&8
\end{array}
\]
The trace-one pattern of the first two examples already
fails at the third. The proof of divergence does not use
these finite data. Its threshold for a fixed torsion level
is the finite degree bound \(3^{n-1}>D\) in (AG.36.6),
where \(D\) comes from representatives of that level's
frame-torsor characters. No bound for the growth of these
representative degrees with the torsion level is asserted.

% Supplement to AG.34: an exact reduction for any individual-parameter test.
\subsection{Four remaining conductors in the explicit pencil}

\begin{claim}{AG.37}{Proved}\label{clm:AGsidonfour}
For every member \(C_t\) of the explicit family in
Claim~AG.33 over
\(\mathbb F_{2^{38}}\), any nontrivial odd conductor of a
primitive two-rank-zero cyclic cover belongs to
\[
                         \{191,229,457,1257\}.
                                                        \tag{AG.37.1}
\]
More precisely, its theta curve has no point of exact order
\(361,419,571,573,647,687,955,1083\), or \(1371\).
This assertion is a necessary-condition reduction; it does not select
an individual parameter from the nonempty complement in Claim~AG.34.
\end{claim}

\begin{proof}
The scalar Frobenius is \([-2^{19}]\), and the theta curve
\(\Theta=C_t-P\) is symmetric. If \(Q\in\Theta\) has odd
exact order \(n\), put
\[
 H_n=\langle-2^{19},-1\rangle\subseteq(\mathbb Z/n)^\times,
 \qquad S_n=\{0\}\cup H_n.
\]
Every point \(aQ\), \(a\in S_n\), lies on \(\Theta\).
For each nonzero \(d\in\mathbb Z/n\), the ordered pairs
\((a,b)\in S_n^2\) with \(a-b=d\) give distinct geometric
points in \(\Theta\cap(\Theta+dQ)\). The two theta curves
are distinct: translation stabilizing the divisor of a principal
polarization belongs to the trivial kernel of that polarization.
Their intersection number is \(\Theta^2=2\), so there are at
most two such pairs, even if some intersections are nonreduced.

Apply this test to the thirteen conductors already proved exhaustive
in Claim~AG.34. The complete finite calculation gives
\[
\begin{array}{r|r|r}
 n&|H_n|&\max_{d\ne0}\#\{(a,b)\in S_n^2:a-b=d\}\\\hline
 191&10&2\\229&4&2\\361&18&4\\419&22&6\\457&4&2\\
 571&6&4\\573&20&10\\647&34&6\\687&8&4\\955&40&30\\
 1083&18&4\\1257&22&2\\1371&8&4
\end{array}
\]
The standard-library program
\texttt{primitive\_counterexample\_sidon\_review.py}
constructs each subgroup by closure under its two specified generators
and enumerates every ordered pair. Its retained JSON output records
all subgroups and, for each failing difference, every representing
ordered pair. The nine failures exclude even a single theta point of
the indicated exact order. A primitive two-rank-zero conductor must
have such a theta point, by Claim~AG.32. This proves (AG.37.1).
\end{proof}


\subsection{An individually specified finite counterexample}

\begin{claim}{AG.38}{Refuted, with an explicit finite counterexample: nontrivial primitive-cover existence}
\label{clm:AGexplicitprimitivecounterexample}
Let
\[
 \begin{split}
 M(Z)&=Z^{38}+Z^{14}+Z^{10}+Z^9+Z^8+Z^5+Z^2+Z+1,\\
 k&=\mathbb F_2[Z]/(M(Z)),\qquad t=Z\bmod M(Z),
 \qquad d=t^3+1.
 \end{split}                                                    \tag{AG.38.1}
\]
The polynomial \(M\) is irreducible. Let \(C/k\) be the smooth
projective model of the affine equation
\[
             y^2+y=\left(\frac{u^2+t^2u+t}{d}\right)^3,
                                                               \tag{AG.38.2}
\]
and let \(P\) be its unique point at infinity. This is a
supersingular genus-two curve, with \(P\in C(k)\).
For every connected cyclic etale cover
\(D\to C_{\bar k}\) of odd degree \(n>1\), and every generator
\(\sigma\) of its deck group, the abelian subvariety
\[
 P_n(D)=\operatorname{im}\left(
       \prod_{\ell\mid n}(1-\sigma^{n/\ell}):J_D\longrightarrow J_D
                              \right)
\]
has positive two-rank. Thus the nontrivial-cover version of A6 is
false for this completely specified curve, independently of the
counting existence argument.

The degree-one qualification is essential: \(J_C\) has two-rank
zero, so the identity cover satisfies the literal variant that
permits \(n=1\). The example also satisfies A1 and is not a
collision or a counterexample to A1.
\end{claim}

\begin{proof}
First, the standard-library bit-polynomial program
\path{primitive_parameter_field_check.py} verifies
\[
 Z^{2^{38}}\equiv Z\pmod M,
 \qquad
 \gcd(M,Z^{2^2}-Z)=\gcd(M,Z^{2^{19}}-Z)=1.
                                                               \tag{AG.38.3}
\]
Since the prime divisors of \(38\) are \(2\) and \(19\),
the finite-field irreducibility criterion proves that \(M\) is
irreducible. Its retained output gives every successive squaring
residue, using bitwise addition, multiplication, and division in
\(\mathbb F_2[Z]\). In particular \(t\notin\mathbb F_4\),
and \(d\ne0\).

Claim~AG.33 applies to (AG.38.2). It supplies the genus,
supersingularity, unique point at infinity \(P\), which is rational,
scalar Frobenius
\([-2^{19}]\), and the pointed degree-two isogeny
\(\alpha:J_C\to E^2\), where \(E:y^2+y=x^3\). On the
embedded theta curve \(C-P\), its two coordinates have
\[
 A(u)=\frac{u^2+t^2u+t}{d},\qquad
 B(u)=\frac{tu^2+u+t^2}{d}                              \tag{AG.38.4}
\]
as their elliptic \(x\)-coordinates. Both elliptic images of a
finite point are finite, and \(P\) maps to \((0_E,0_E)\).
The theta divisor relevant to the primitive two-rank criterion
is precisely \(C-P\).

By Claims~AG.32, AG.37, and~\ref{clm:ACpacketconductors}, the
only remaining possible conductors of a primitive-zero cover are
\(191\) and \(229\). We prove the stronger assertion that
\(C-P\) has no nonzero point killed by either integer.

Here is a direct characteristic-two verification of the elliptic
torsion polynomials used in that test. On \(E\), the tangent
and chord formulas give the rational-function identities
\[
 x([2]R)=x(R)^4,
 \qquad
 x(R+S)+x(R-S)=\frac{1}{(x(R)+x(S))^2}.
                                                               \tag{AG.38.5}
\]
For the second identity, replacing \(S=(x_S,y_S)\) by
\(-S=(x_S,y_S+1)\) changes the chord slope by
\(1/(x_R+x_S)\); the other terms in the two addition formulas
cancel. The tangent slope is \(x(R)^2\), giving the first
identity. Hence the rational functions \(r_n(X)=x([n]R)\),
where \(X=x(R)\), satisfy
\[
 r_1=X,\qquad r_{2m}=r_m^4,\qquad
 r_{2m+1}=X+\frac{1}{(r_m+r_{m+1})^2}.             \tag{AG.38.6}
\]
All equalities are rational-function identities, so values where
the displayed addition formula has a denominator zero are
handled by the reduced rational function. Its finite poles are
exactly the \(x\)-coordinates of nonzero points killed by \(n\).

Independently, define polynomials in \(\mathbb F_2[X]\) by
\[
 \begin{gathered}
 \psi_1=\psi_2=\psi_4=1,\qquad \psi_3=X^4+X,\\
 \psi_{2m}=\psi_m^4,\qquad
 \psi_{2m+1}=\psi_{m+2}\psi_m^3+
                       \psi_{m-1}\psi_{m+1}^3\quad(m\ge2).
 \end{gathered}                                                 \tag{AG.38.7}
\]
The program \path{primitive_parameter_rational_check.py}
constructs (AG.38.6) using reduced numerator--denominator pairs,
without using (AG.38.7) for that construction. For both required
integers it verifies that the reduced denominator is exactly
\(\psi_n^2\). Thus no general division-polynomial formula in
characteristic different from two is being imported. The checks are
\[
\begin{array}{c|r|r|r}
 n&\deg\psi_n&\deg\operatorname{num}(r_n)&
                         \deg\operatorname{den}(r_n)\\\hline
 191&18240&36481&36480\\
 229&26220&52441&52440
\end{array}
\]
and the reduced numerators are coprime to their denominators.

Finally, exact polynomial arithmetic over the explicitly defined
field \(k\) gives polynomials \(U_n,V_n\in k[u]\) satisfying
\[
       U_n(u)\psi_n(A(u))+V_n(u)\psi_n(B(u))=1
                \qquad(n=191,229).                            \tag{AG.38.8}
\]
The producer \path{primitive_parameter_gcd.py} evaluates
(AG.38.7) directly at the quadratic polynomials (AG.38.4), computes
the extended gcd, and checks the whole identity (AG.38.8) by
polynomial multiplication. The degrees of the two evaluated
polynomials are \(36480\) and \(52440\), respectively; in
each case both Bezout coefficients have degree four less.
All four polynomials for each integer are retained as binary
coefficient arrays. Their encoding, degrees, and SHA-256 hashes
are specified in \path{primitive_parameter_gcd.json}:
coefficients are in ascending degree, each represented by an
unsigned 64-bit little-endian integer whose bit \(i\) is the
coefficient of \(t^i\). The recorded runtime dependency is
\texttt{python-flint==0.9.0}. This is a complete exact identity
certificate, not a numerical approximation or a random-point test.
An independently written loader, \texttt{A6\_individual\_curve\_verify.py},
also verifies every coefficient-array hash, degree, and encoding,
reconstructs both compositions independently, and checks both full
Bezout identities; its retained run exits successfully.

If a finite \(Q\in C(\bar k)\) had \([Q-P]\ne0\) killed
by \(n=191\) or \(229\), its image under \(\alpha\) would
be killed by \(n\) in both elliptic factors. By (AG.38.5)--
(AG.38.7), the finite value \(u(Q)\) would be a common zero of
\(\psi_n(A(u))\) and \(\psi_n(B(u))\), contradicting
(AG.38.8). The remaining point \(P\) represents the origin
and cannot supply a primitive character. Therefore neither
remaining conductor occurs, over any finite extension or over
\(\bar k\), and the claimed assertion follows.

Finally \(J_C[2^\infty](\bar k)=0\). A collision would already
give a two-primary class \([Q-P']\) on \(C\), for the two
colliding points \(P',Q\), by R1. This class is zero, and the
Abel map of a positive-genus curve is injective on geometric
points. Hence \(Q=P'\), proving the stated A1 assertion.
\end{proof}

% Independently audited by root and the F/topos agent; integrated as AG.39.
\subsection{A failure of full prime-to-\texorpdfstring{\(p\)}{p} configuration exactness}

\begin{claim}{AG.39}{Proved: the full prime-to-\(p\) configuration fiber does not inject}
\label{clm:AGconfigurationodd}
Let \(C\) be a smooth proper connected complex curve of genus
\(g\geq2\), let \(P\in C\), and put
\[
 U=\{(x,z)\in C\times C:x\ne z\},\qquad F=C-\{P\}.
\]
Use the second projection \(U\to C\) to identify \(F\) with a
fiber. For every prime \(p\), write \(\Gamma^{(p')}\) for the
inverse limit of the finite quotients of a group \(\Gamma\)
whose orders are prime to \(p\). Then
\[
 \pi_1(F)^{(p')}\longrightarrow
 \pi_1(U)^{(p')}                                     \tag{AG.39.1}
\]
is not injective. The same assertion holds for the geometric
prime-to-\(p\) etale fundamental groups of every smooth proper
curve of genus at least two in characteristic \(p\).
In particular it holds for the smooth projective model of
\[
                    C_0/\mathbb F_2:\quad y^2+y=x^5,
             \qquad P=\infty.                       \tag{AG.39.2}
\]

This refutes the full prime-to-\(p\) configuration exactness
condition proposed as a sufficient route in
\cite[Remark 3.13]{AR_ST2009}. It does not supply two coincident
point sections and does not refute A1. The proof specifies a finite
curve and finite branched-cover construction; it does not enumerate
a particular nonidentity profinite kernel element or an individual
finite monodromy matrix witnessing the failure.
\end{claim}

\begin{proof}[Independent proof for \(p=2\)]
For this proof write \((-)^{\mathrm{odd}}=(-)^{(2')}\).
The extension to every prime is proved below by a separate
congruence argument.
\emph{An odd-degree branched family.}
Choose a connected cyclic etale cover
\(a:C'\to C\) of degree three. Put \(m=3^{2g}\), and let
\(b:B\to C\) be the connected etale cover corresponding to
\[
 \ker\bigl(\pi_1(C)\longrightarrow H_1(C,\mathbb F_3)\bigr).
                                                               \tag{AG.39.3}
\]
It is Galois of degree \(m\), factors through \(C'\), and
\(b^*:H^1(C,\mathbb F_3)\to H^1(B,\mathbb F_3)\) is zero:
each homomorphism \(\pi_1(C)\to\mathbb F_3\) restricts to
zero on the subgroup in (AG.39.3).

On \(Y=C'\times B\) consider the effective divisor
\[
          D=(a\times b)^{-1}(\Delta_C).
                                                               \tag{AG.39.4}
\]
It consists of three disjoint graphs \(B\to C'\). Its first
Chern class is divisible by three in \(H^2(Y,\mathbb Z)\).
Indeed, in the Kunneth decomposition modulo three, the two
degree-two summands of the diagonal class pull back by maps of
degrees \(3\) and \(m\), and its cross summand vanishes because
\(b^*H^1(C,\mathbb F_3)=0\). The integral cohomology of a
product of compact oriented surfaces is torsion-free, so vanishing
modulo three is equivalent to integral divisibility.
The class \(c_1(D)/3\) is integral of type \((1,1)\); the
Lefschetz \((1,1)\) theorem gives a line bundle with that Chern
class. Its cube differs from \(\mathcal O_Y(D)\) by an element
of \(\operatorname{Pic}^0(Y)\). Surjectivity of multiplication
by three on this complex abelian variety therefore gives
\[
                         L^{\otimes3}\simeq\mathcal O_Y(D).
                                                               \tag{AG.39.5}
\]

Take the cyclic cover \(S\to Y\) obtained by adjoining a cube
root of the canonical section of \(L^{\otimes3}\). It is a
connected smooth projective surface. The composite
\(f:S\to B\) is smooth and proper with connected fibers:
locally along any graph the cover has the form \(w^3=t\), where
\((t,z)\) are coordinates with \(z\) on \(B\); these become
the coordinates \((w,z)\) on \(S\). Each fiber is a cyclic
triple cover of \(C'\), branched at the three distinct points
\(a^{-1}(b(z))\), each with ramification index three. Consequently
\[
 g(C')=3g-2,\quad g(B)=m(g-1)+1,\quad
                  h:=g(S_z)=9g-5.                   \tag{AG.39.6}
\]
This is the usual cyclic construction of a Kodaira fibration;
compare \cite[Section 1.3]{AG_Rollenske2010}. All the particular
divisibility and smoothness assertions needed here were proved
above.

The family is nonisotrivial. To see this without any genericity
assumption, suppose every fiber were isomorphic to a fixed curve
\(T\). There are only finitely many degree-three morphisms
\(T\to C'\): their Hom scheme is of finite type, and the tangent
space at a morphism \(q\) is
\(H^0(T,q^*T_{C'})=0\), since this line bundle has negative
degree. Thus there would be only finitely many possible branch
sets in \(C'\). The sets \(a^{-1}(b(z))\), however, vary with
the surjective map \(b:B\to C\). This is a contradiction.

\emph{The required comparison of completions.}
Let \(S^\circ\) be the complement of the ramification divisor
in \(S\). There is a tower of finite etale covers
\[
 S^\circ\longrightarrow Y-D
 \longrightarrow (C\times B)-\operatorname{graph}(b)
 \longrightarrow U.                                 \tag{AG.39.7}
\]
The first two arrows are cyclic of degree three, and the last is
the Galois base change of (AG.39.3). The Galois closure of the
composite has a three-group as its group: the Galois closure of a
tower of covers with three-group Galois groups embeds in the
corresponding iterated wreath product. In a fiber of \(U\to C\),
the cover \(S_z^\circ\to F\) has degree nine and likewise has
three-group Galois closure.

Here is the elementary topology-of-completion fact being used.
If \(H\leq E\) contains a normal subgroup of finite odd index,
then \(H^{\mathrm{odd}}\) is the open subgroup given by the
closure of \(H\) in \(E^{\mathrm{odd}}\). For any normal
odd-index subgroup \(N\) of \(H\), first intersect \(N\)
with that fixed normal subgroup of \(E\), and then take its
\(E\)-core. The resulting quotient is an extension of an odd
group by a subgroup of a finite product of odd groups. It has
odd order and its restriction to \(H\) refines the given
quotient. This proves fullness of the induced pro-odd topology.

Suppose now that (AG.39.1) were injective. Apply the preceding
fact to both the total cover and its degree-nine fiber cover.
The resulting homomorphism
\[
 \pi_1(S_z^\circ)^{\mathrm{odd}}\longrightarrow
                 \pi_1(S^\circ)^{\mathrm{odd}}        \tag{AG.39.8}
\]
is injective. The discrete fiber group is normal in the discrete
total group and the latter surjects onto \(\pi_1(B)\):
\(S^\circ\to B\) is a locally trivial punctured-surface
fibration with connected fiber. Thus its compact pro-odd closure
is normal in the completed total group. These assertions also
follow from the Fadell--Neuwirth sequence for \(U\to C\) and
the pullback covers; see \cite[Section 2.1]{AG_Bellingeri2004}.

For every odd prime \(\ell\), compactification of the punctured
fiber gives a quotient
\[
 \pi_1(S_z^\circ)^{\mathrm{odd}}\longrightarrow
                          H_1(S_z,\mathbb F_\ell).
                                                               \tag{AG.39.9}
\]
Its kernel is stable under the discrete total group, because the
family preserves the peripheral loops, and therefore under its
pro-odd closure. The conjugation action on this finite quotient
is continuous. For example this follows by testing the images
of a finite generating set of the pro-odd fiber group. Hence its
image has odd order. Inner conjugation by the fiber group acts
trivially on (AG.39.9). Since the total discrete group surjects
onto \(\pi_1(B)\), the image just obtained is precisely the
usual monodromy image on \(H_1(S_z,\mathbb F_\ell)\).
In particular, if
\[
 \Gamma=\operatorname{im}\bigl(\pi_1(B)\to
                    \operatorname{Sp}_{2h}(\mathbb Z)\bigr),
\]
then the reduction of \(\Gamma\) modulo every odd prime has
odd order. No assertion that completion is left exact for an
arbitrary surface fibration has been used.

\emph{The monodromy contradiction, in characteristic zero.}
The Feit--Thompson theorem makes each of those finite reductions
solvable \cite{AC_FeitThompson1963}. The uniform solvable-length
theorem for linear groups supplies an integer \(d(2h)\),
depending only on the dimension and not on the field, that
bounds their derived lengths; see
\cite[Section 1, printed p. 59]{AG_Newman1972}.
Every element of \(\Gamma^{(d(2h))}\) therefore becomes the
identity modulo every odd prime. An integral matrix with this
property is the identity. Thus \(\Gamma\) is solvable.

On the other hand, Deligne's theorem states that the identity
component of the Zariski closure of the monodromy of a smooth
proper complex algebraic family is semisimple
\cite[Corollary 4.2.9(a)]{AG_Deligne1971}.
The Zariski closure of a solvable group is solvable, so that
identity component must be trivial. Hence \(\Gamma\) is finite.
After a finite etale cover of \(B\) the integral local system
is constant. Its period map is then a globally defined
holomorphic map from a compact curve into Siegel upper half
space. In its bounded realization all coordinate functions are
constant. The principally polarized Jacobian is constant, and
Torelli makes the fibers of \(f\) isomorphic. This contradicts
the nonisotriviality already proved and proves (AG.39.1).

\emph{Passage to characteristic two.}
This last Hodge argument is used only over \(\mathbb C\).
Let \((C,P)\) be a pointed smooth proper curve over an
algebraically closed field of characteristic two. Choose a
finitely generated subfield over \(\mathbb F_2\) over which
\((C,P)\) is defined, and replace it by its algebraic closure
\(k_0\) inside the given field. Prime-to-two fundamental groups
of these smooth proper pairs are invariant under algebraically
closed extension, so it suffices to work over \(k_0\).
A pointed smooth proper lifting over \(W(k_0)\) exists: the
obstruction group for deforming a pointed smooth curve is
\(H^2(C,T_C(-P))=0\), and a relatively ample line bundle
algebraizes the formal lifting. The countable field \(k_0\)
has a Witt fraction field of cardinality at most that of
\(\mathbb C\); an algebraic closure of this characteristic-zero
field therefore embeds abstractly into \(\mathbb C\). This
gives the complex geometric generic fiber used below.

For this lift, both
\((\mathcal C,\mathcal P)\) and
\((\mathcal C\times\mathcal C,\Delta_{\mathcal C})\)
are smooth proper pairs with relative normal-crossing boundary.
SGA~1 XIII, Corollary~2.9 and the specialization construction in
Section~2.10 identify their maximal prime-to-two fundamental
groups on the geometric generic and special fibers
\cite{AR_SGA1}. In more detail, Corollary~2.9 makes the nonabelian
\(R^1\) for every finite group of odd order locally constant;
its base-change statement and Section~2.10 give these
specialization isomorphisms. Functoriality for
\(x\mapsto(x,\mathcal P)\) identifies the fiber-inclusion maps
as well. Injectivity in characteristic two would therefore imply
injectivity for the complex curve, which was disproved above.

For the completely specified example (AG.39.2), a lift requires
no deformation argument. The equation \(y^2+y=x^5\) over
\(\mathbb Z_2\), with the chart at infinity
\[
                         v^2+u^3v=u,
\]
gives the required smooth genus-two model and the section
\(u=v=0\). On the special fiber the finite chart is smooth
because its derivative with respect to \(y\) is one; the
infinity chart is smooth at that section because its derivative
with respect to \(u\) is a unit. Its generic fiber embeds into
\(\mathbb C\), so the same functorial specialization diagram
applies.
\end{proof}

\paragraph{Finite numerical checks and scope.}
For genus two, (AG.39.3)--(AG.39.7) have
\[
 m=81,\qquad g(C')=4,\qquad g(B)=82,
 \qquad g(S_z)=13,
\]
and the punctured fiber has three punctures and first Betti
number \(28\); its compactification has first Betti number
\(26\). The retained exact-rational program
\path{configuration_family_invariants.py} also verifies
\[
 D^2=-486,\quad K_YD=972,\quad K_Y^2=1944,
 \quad e(S)=3888,\quad K_S^2=9072,\quad
 \operatorname{sign}(S)=432.
                                                               \tag{AG.39.10}
\]
The signature is an independent numerical consistency check,
not a substitute for the monodromy argument. The new obstruction
concerns completing the \emph{whole} configuration group to
pro-odd. It does not contradict the exact sequence with full
tame base and only the fiber replaced by its prime-to-two
quotient, nor the exact single-prime pro-\(\ell\) sequences.

\paragraph{A cohomological obstruction to naive cuspidal devissage.}
There is a concrete loss of information when one tries to
iterate maximal cuspidally pro-\(\ell\) reconstruction through
ramified covers. Work over an algebraically closed field of
characteristic two, or over \(\mathbb C\), and take the
pro-odd groups throughout. Let \(C\) have genus \(g\geq2\),
let \(P\in C\), and write
\[
 F=\pi_1(C-P)^{\mathrm{odd}},\quad
 G=\pi_1(C)^{\mathrm{odd}},\quad K=\ker(F\to G),
 \quad N=\ker(K\to K^{(3)}).
\]
Here \(K^{(3)}\) denotes the maximal pro-three quotient,
so \(F/N\) is the maximal cuspidally pro-three quotient over
\(G\). Choose a cyclic etale cubic \(a:C'\to C\), and a
cyclic cubic \(D\to C'\) branched exactly at the three
points of \(a^{-1}(P)\), with ramification index three at
each. Such a cover exists because the degree-three branch
divisor has a cube root in \(\operatorname{Pic}(C')\):
choose a line bundle of degree one and correct its cube
by divisibility of \(\operatorname{Pic}^0(C')\) over the
algebraically closed field.
It has genus \(h=9g-5\).

Put \(D^\circ=D-\pi^{-1}(P)\). Its degree-nine etale cover
of \(C-P\) is visible in \(F/N\), since its Galois closure
has three-group Galois group. Let \(H\leq F/N\) be its
open subgroup, and let \(\overline H\) be the quotient of
\(H\) by the closed normal subgroup generated by its three
puncture inertia groups. Then, for every odd prime
\(\ell\ne3\),
\[
 \dim_{\mathbb F_\ell}H^1(\overline H,\mathbb F_\ell)
       =2g(C')=6g-4
       <18g-10
       =\dim_{\mathbb F_\ell}H^1(D,\mathbb F_\ell).
                                                               \tag{AG.39.11}
\]
Thus \(\overline H\) is a proper quotient of the full
group \(\pi_1(D)^{\mathrm{odd}}\), already detected on
finite abelian quotients. For \(g=2\), the two dimensions
are \(8\) and \(26\).

Here is a direct proof of the comparison, including its
surjectivity assertion. Let \(F'\leq F\) correspond to
\(C'-a^{-1}(P)\), and let
\(\chi:F'\twoheadrightarrow\mathbb Z/3\) define
\(D^\circ\). The group \(K\) is contained in \(F'\),
and \(\chi|_K\) is surjective, because a branch inertia
generator lies in \(K\) and maps nontrivially. If
\(I=\ker\chi\), this gives \(F'=IK\). Also \(N\subset I\),
and hence
\[
 1\longrightarrow (I\cap K)/N\longrightarrow
 H=I/N\longrightarrow \pi_1(C')^{\mathrm{odd}}
 \longrightarrow1.                                  \tag{AG.39.12}
\]
The left-hand term is a closed subgroup of the pro-three
group \(K/N\). Every homomorphism from \(H\) to
\(\mathbb F_\ell\) therefore factors uniquely through
the last group when \(\ell\ne3\). The puncture inertia
groups in \(H\) lie in the left-hand term, so killing them
does not change these homomorphisms. This proves the first
equality in (AG.39.11); the second is the usual surface
cohomology dimension and the computed genus of \(D\).
The natural map from \(\pi_1(D)^{\mathrm{odd}}\) to
\(\overline H\) is surjective by the construction.

This is an obstruction to the specific procedure of taking
an open subgroup of a cuspidally pro-three quotient and then
killing its puncture inertia. It is not a theorem that all
cuspidal quotients together determine insufficient information
by every possible method. In a solvable induction, the next
cuspidally pro-\(\ell\) theorem needs the \emph{full} proper
prime-to-two group of the intermediate ramified compactification.
Equation (AG.39.11) shows exactly why the preceding formal
operations do not provide that input. That ramified
compactification is also absent from the original proper
unramified tower.

% Extension of AG.39; intended for incorporation into that statement/proof,
% not a second independently numbered configuration claim.

\paragraph{The congruence lemma needed for every excluded prime.}
Let \(\Gamma\subset\operatorname{GL}_N(\mathbb Z[1/M])\) be
finitely generated. Suppose that the identity component \(G\) of its
Zariski closure over \(\mathbb Q\) is a nontrivial semisimple group.
For every prime \(p\) and every integer \(a\geq1\), there are
infinitely many rational primes \(\ell\ne p\), with \(\ell\nmid M\),
such that
\[
       p^a\mid
       \#\operatorname{im}\bigl(\Gamma\longrightarrow
                    \operatorname{GL}_N(\mathbb F_\ell)\bigr).
                                                        \tag{AG.39.13}
\]

\begin{proof}
The subgroup \(\Gamma_0=\Gamma\cap G(\mathbb Q)\) has finite index,
is finitely generated, and is Zariski dense in \(G\). Let
\[
     1\longrightarrow Z\longrightarrow \widetilde G
          \xrightarrow{\pi}G\longrightarrow1
\]
be its simply connected central covering, and put \(c=\#Z\).
The connecting map
\[
               \delta:G(\mathbb Q)\longrightarrow H^1(\mathbb Q,Z)
\]
is a group homomorphism. Indeed, for lifts of two rational elements,
the two Galois cocycles multiply because their values lie in the
central subgroup \(Z\). The abelian group \(H^1(\mathbb Q,Z)\)
is killed by \(c\). Therefore the image of the finitely generated
group \(\Gamma_0\) under \(\delta\) is finite. Its kernel
\(\Gamma_1\) has finite index in \(\Gamma_0\), and every element
of \(\Gamma_1\) has a rational lift through \(\pi\).

Choose a finite generating set of \(\Gamma_1\), choose rational
lifts of its generators, and let \(\widetilde\Gamma\) be the
group generated by those lifts. Its image is exactly \(\Gamma_1\).
Its Zariski closure is all of \(\widetilde G\): the image of that
closure under the finite morphism \(\pi\) is closed and contains
the dense subgroup \(\Gamma_1\); equality of dimensions, and
connectedness of \(\widetilde G\), finish the assertion.

Choose a faithful rational representation of \(\widetilde G\).
After enlarging a finite set of primes \(S\), all chosen lifts
are integral over \(\mathbb Z[S^{-1}]\), and \(\pi\), the
original faithful representation of \(G\), and their central
kernel spread out over this ring. We may require that these
models are smooth semisimple group schemes, that the morphism
is a central isogeny of degree \(c\), and that the embedding
of \(G\) into the original matrix group is a closed immersion.
Thus, outside \(S\), reducing \(\pi(\widetilde\gamma)\) gives
the same matrix as reducing the original element of \(\Gamma_1\).
The kernel of this comparison on finite-field points has at
most \(c\) elements.

Nori's strong-approximation theorem, in precisely its
connected and simply connected form
\cite[Theorem 5.4, printed p. 274]{AC_Nori1987}, says that the
congruence closure of \(\widetilde\Gamma\) is open in
\(\widetilde G(\prod_{\ell\notin S}\mathbb Z_\ell)\).
Consequently, after excluding finitely many more primes,
\[
       \widetilde\Gamma\bmod\ell
                      =\widetilde G(\mathbb F_\ell).
\]
Here smoothness gives surjectivity of reduction from
\(\mathbb Z_\ell\)-points, and an open subgroup of a product
contains the full product outside finitely many factors.

Choose a finite Galois number field \(K\) which splits
\(\widetilde G\) and contains the roots of unity of order
\(p^{a+v_p(c)}\). Exclude also the finitely many rational
primes needed to spread the chosen splitting isomorphism
over the ring of integers of \(K\). Infinitely many rational primes split
completely in \(K\). One may use the usual splitting-prime
theorem, or the following elementary sufficient form. If
\(K=\mathbb Q(\alpha)\) with \(\alpha\) integral and minimal
polynomial \(f\), infinitely many primes divide nonzero values
\(f(n)\). Indeed, if all such prime divisors belonged to a
finite set with product \(R\), choose \(b\) with
\(d=f(b)\ne0\). The integer polynomial
\(f(b+dRt)/d\) is congruent to one modulo \(R\) and has
nonconstant absolute value, a contradiction for sufficiently
large integers \(t\). Away from the discriminant and the
index of \(\mathbb Z[\alpha]\), a root of \(f\) modulo a
prime gives a residue-degree-one prime of \(K\). Since
\(K/\mathbb Q\) is Galois, that rational prime splits completely.

Discard the finite exceptional set, and choose any such
\(\ell\ne p\). Then \(\widetilde G_{\mathbb F_\ell}\)
is a nontrivial split semisimple group and
\(p^{a+v_p(c)}\mid\ell-1\). Its split maximal torus has
positive rank, so
\[
       v_p\bigl(\#\widetilde G(\mathbb F_\ell)\bigr)
                          \geq a+v_p(c).
\]
Its image in the original \(\operatorname{GL}_N(\mathbb F_\ell)\)
loses at most \(v_p(c)\) from this valuation, and that image is
the reduction of \(\Gamma_1\). It is a subgroup of the
reduction of \(\Gamma\). This proves (AG.39.13).
\end{proof}

\paragraph{Replacement of the pro-odd restriction in AG.39.}
Fix an arbitrary prime \(p\), and write \((-)^{(p')}\) for
completion with respect to all finite groups of order prime
to \(p\). For every smooth proper complex curve \(C\) of
genus \(g\geq2\), every \(P\in C\), and
\(U=(C\times C)-\Delta_C\), the map
\[
       \pi_1(C-\{P\})^{(p')}
                 \longrightarrow \pi_1(U)^{(p')}
                                                        \tag{AG.39.14}
\]
is not injective. The corresponding geometric prime-to-\(p\)
map is likewise noninjective for every smooth proper curve of
genus at least two over an algebraically closed field of
characteristic \(p\).

\begin{proof}
In the explicit family construction of AG.39 choose a prime
\(r\ne p\): one may take \(r=2\) for odd \(p\), and
\(r=3\) for \(p=2\). Replace its cyclic triple cover by a
cyclic degree-\(r\) etale cover \(a:C'\to C\), and take
\(b:B\to C\) to be the full mod-\(r\) homology cover,
of degree \(r^{2g}\). This kills the cross term of the
diagonal class modulo \(r\), and both remaining Kunneth
terms vanish modulo \(r\). The same integral-cohomology and
Picard-divisibility argument therefore gives
\[
       L^{\otimes r}\simeq
       \mathcal O_{C'\times B}\bigl((a\times b)^{-1}\Delta_C\bigr).
\]
The resulting cyclic cover is a smooth proper family
\(S\to B\). Each fiber is a degree-\(r\) cover of \(C'\)
branched at the \(r\) distinct points \(a^{-1}(b(z))\), and
\[
   g(C')=1+r(g-1),\qquad g(B)=1+r^{2g}(g-1),\qquad
   g(S_z)=1+r^2(g-1)+\frac{r(r-1)}2.
\]
It is nonisotrivial by the same finite-Hom argument: a fixed
source curve has only finitely many degree-\(r\) maps to the
fixed hyperbolic curve \(C'\), while their required branch
sets vary through the fibers of the surjection \(a\).

Every cover in the punctured total-space tower has an
\(r\)-group Galois group, and its composite has an
\(r\)-group Galois closure by the wreath-product argument.
The punctured fiber has the same property. An open subgroup
containing a normal subgroup of finite prime-to-\(p\) index
has its full pro-prime-to-\(p\) topology induced from the
ambient group: taking normal cores produces an extension
of a prime-to-\(p\) group by a subgroup of a product of such
groups. This uses no solvability assumption.

Consequently, if (AG.39.14) were injective, the completed
punctured-fiber group would be a closed normal subgroup of
the completed total group of \(S^\circ\to B\). For every
\(\ell\ne p\), its compactified homology
\(H_1(S_z,\mathbb F_\ell)\) is a finite quotient with
monodromy-invariant kernel. Conjugation is continuous, the
fiber acts trivially on this abelian quotient, and the
discrete total group surjects onto \(\pi_1(B)\). Therefore
every reduction modulo \(\ell\ne p\) of the integral
monodromy group
\[
       \Gamma\subset\operatorname{Sp}_{2g(S_z)}(\mathbb Z)
\]
would have order prime to \(p\). These are reductions of
the actual integral homology local system, rather than
abstract finite quotients of an unrelated group.

Deligne's semisimplicity theorem used in AG.39 makes the
connected Zariski closure of \(\Gamma\) semisimple. It is
nontrivial: a trivial identity component would make
\(\Gamma\) finite, so after a finite etale base change the
period map would be a holomorphic map from a compact curve
to the bounded realization of Siegel space. It would be
constant; Torelli would make the family isotrivial. The
congruence lemma, with \(a=1\), now contradicts the asserted
prime-to-\(p\) orders of all the reductions. This proves
(AG.39.14) over \(\mathbb C\).

For characteristic \(p\), lift the pointed smooth curve
as in AG.39 and apply the prime-to-\(p\) specialization
isomorphisms for the two smooth proper pairs
\((\mathcal C,\mathcal P)\) and
\((\mathcal C\times\mathcal C,\Delta_{\mathcal C})\).
SGA~1 XIII, Corollary~2.9 and Section~2.10 apply because
all the finite groups in question have order invertible
on the base. Functoriality identifies the two fiber maps.
Thus injectivity of the special-fiber map would imply
injectivity of its complex counterpart, which has just
been contradicted. The Hodge and period-map argument has
been used only in characteristic zero.
\end{proof}

The extension eliminates the full configuration exactness
shortcut for every excluded prime. It does not produce a
collision, and by itself does not decide A1.
The later AG40--AG42 settle A1; universal relative-witness
existence remains a separate question. The characteristic-two proof in AG.39
via Feit--Thompson and bounded solvable length remains an
independent alternative to the strong-approximation step.

\paragraph{What the same-prime quotient does retain.}
The preceding cross-prime rank loss must not be confused
with a loss of the entire zeta function. There is a general
elementary complement. Let \(F\) be a profinite group,
let \(K\triangleleft F\) be closed, fix a prime \(r\),
and put \(N=\ker(K\to K^{(r)})\). For every open
subgroup \(I\leq F\) containing \(N\), the natural map
induces an isomorphism
\[
                         I^{(r)}\simeq (I/N)^{(r)}.
                                                        \tag{AG.39.15}
\]
Here the superscript denotes the maximal pro-\(r\)
quotient. In particular, the assertion is not restricted
to open subgroups having an \(r\)-group Galois core.

To prove it, take an open normal subgroup \(J\triangleleft I\)
such that \(I/J\) is an \(r\)-group, and set
\(I_0=\operatorname{core}_F(I)\).
The group \(J\cap I_0\) is normal in \(I_0\), and
\(I_0/(J\cap I_0)\) embeds in \(I/J\).
Let \(M=\operatorname{core}_F(J\cap I_0)\).
Then \(I_0/M\) embeds in a finite product of
\(r\)-groups, so it is an \(r\)-group. Since
\(N\triangleleft F\) and \(N\subset I\), we have
\(N\subset I_0\); consequently the image of \(K\)
in \(F/I_0\) is an \(r\)-group. Its image in
\(F/M\) is an extension of an \(r\)-group by a
subgroup of \(I_0/M\), and is therefore an
\(r\)-group as well. The defining property of \(N\)
now gives \(N\subset M\subset J\). Thus every
finite \(r\)-quotient of \(I\) kills \(N\),
which proves (AG.39.15).

Apply this to the geometric punctured-curve group and
its cuspidal kernel. For every finite cover visible in
the maximal cuspidally pro-\(r\) quotient, quotienting
by the puncture inertia and then taking the maximal
pro-\(r\) quotient recovers the \emph{full} proper
pro-\(r\) group of its smooth compactification.
Indeed, formation of a maximal pro-\(r\) quotient
commutes with quotienting by the closed normal subgroup
generated by the specified inertia groups. This applies
in particular to the ramified cover \(D\) in
(AG.39.11), even though its other-prime cohomology has
been lost.

If the curves and cover are defined over a finite field
of characteristic different from \(r\), this comparison
is compatible with the arithmetic Frobenius action.
It therefore retains \(H^1(D_{\bar k},\mathbb Z_r)\)
with Frobenius, and hence the exact integral Weil
polynomial and zeta function of \(D\). Equality of the
\(r\)-adic characteristic polynomials is equality of
their integer polynomials, since
\(\mathbb Z[T]\hookrightarrow\mathbb Z_r[T]\).
Thus the obstruction in (AG.39.11) concerns the full
proper prime-to-\(p\) group needed by a subsequent
different-prime reconstruction theorem; it does not
prevent transfer of these zeta functions. No inference
from that transfer to point separation is made here.
In particular, the Weil polynomial does not by itself
specify the integral \(p\)-primary rational torsion
exponent, or supply R1-compatible point lifts on these
ramified covers.


\subsection{A conditional cuspidal criterion for point separation}

\begin{claim}{AG.40}{Proved: a conditional cuspidal separation criterion}
\label{clm:AGcuspidalcriterion}
Let \(k\) be a finite field of characteristic \(p\), let
\(C/k\) be smooth, proper and geometrically connected of
genus at least two, and let \(P,Q\in C(k)\).
Write \(\Pi_C\) for the arithmetic fundamental group with
maximal geometrically prime-to-\(p\) quotient. Fix a prime
\(r\ne p\). Let \(G_P\) and \(G_Q\) be the maximal
cuspidally pro-\(r\) quotients of the corresponding
arithmetic prime-to-\(p\) groups of \(C-P\) and \(C-Q\),
relative to their maps to \(\Pi_C\).

Suppose there is a continuous isomorphism
\[
                  \beta:G_P\xrightarrow{\sim}G_Q
\]
over \(\operatorname{id}_{\Pi_C}\) which carries the
cuspidal decomposition group at \(P\) to a conjugate of
the cuspidal decomposition group at \(Q\). Then \(P=Q\).
The isomorphism is in particular over the identity of
\(G_k\), so it preserves arithmetic Frobenius.
\end{claim}

This is an implication from the existence of the specified
isomorphism. Its proof does not use a configuration-space
fiber exactness theorem. The required construction from
conjugate proper point sections is supplied separately by
Claims~\ref{clm:AGfixedrconfiguration} and
\ref{clm:AGcuspidaltransport}.

\begin{proof}
Assume \(P\ne Q\). We give the argument in five steps.

\emph{1. The point difference is \(p\)-primary.}
Projection of the cuspidal decomposition groups to the
proper quotient shows that the sections of \(P\) and
\(Q\) in \(\Pi_C\) are conjugate. The conjugator
may be made geometric: multiply it by an element of
the target point decomposition group with the same
arithmetic projection, which does not change the
conjugated subgroup. Thus these sections are
geometrically conjugate. Passing
to the geometric abelian quotient and the Jacobian
\(J_C\), their difference has zero Kummer class for
every prime \(\ell\ne p\). The Kummer map on the
\(\ell\)-primary part of \(J_C(k)\) is injective:
the usual Kummer exact sequences, followed by inverse
limit, give the injection of the \(\ell\)-adic
completion of \(J_C(k)\) into
\(H^1(G_k,T_\ell J_C)\). Since \(J_C(k)\) is finite,
this completion is precisely its \(\ell\)-primary part.
Consequently \([P-Q]\) has \(p\)-power order. There are
an integer \(e\geq0\) and \(f\in k(C)^\times\) with
\[
                  \operatorname{div}(f)=p^e(P-Q).
                                                        \tag{AG.40.1}
\]
Indeed this divisor is principal over \(\bar k\), and
its rational generator descends by Hilbert~90, since
the divisor is \(k\)-rational.

\emph{2. Every visible finite cover has its zeta function
and its number of missing points preserved.}
For the geometric group of \(C-P\), let \(K_P\) be
the kernel over the full proper geometric prime-to-\(p\)
group and let \(N_P=\ker(K_P\to K_P^{(r)})\).
Equation (AG.39.15) applies to every open subgroup
containing \(N_P\). Thus, for every finite cover
visible in \(G_P\), killing its cusp inertia and then
taking its maximal pro-\(r\) quotient gives the full
proper pro-\(r\) group of its smooth compactification.
This construction is natural for the arithmetic action.
The same holds for \(G_Q\).

The isomorphism \(\beta\) identifies the relevant
inertia groups, because it identifies the cuspidal
decomposition groups and their geometric kernels.
It therefore identifies the compactified
\(H^1(-,\mathbb Z_r)\) representations, including
Frobenius. The integral Weil polynomials, hence the
point counts of the compactifications over every
finite constant extension, agree. The missing points
also correspond: their decomposition groups and
residue degrees are given by the cuspidal
double-cosets and their arithmetic projections in
the respective open subgroups. It follows that
corresponding \emph{open} covers have the same
number of rational points over their constant
fields and over every further finite extension.

\emph{3. Noncuspidal point sections are transported.}
Extend constants temporarily so that a chosen
noncuspidal point \(R\) is rational. Let
\(D_R^P\subset G_P\) be its decomposition subgroup.
It maps isomorphically to the absolute Galois group
of the new constant field. Every open subgroup
\(H'\subset G_Q\) containing \(\beta(D_R^P)\)
has full arithmetic projection. The source cover
for \(\beta^{-1}(H')\) has a rational noncuspidal
point over \(R\), so the target open cover has a
rational noncuspidal point by Step~2.

The finite sets of rational points of these target
covers form a directed inverse system with the
finite intersection property: any finite collection
of conditions can be checked in the cover for the
intersection of the corresponding open subgroups.
Compactness gives a compatible point of the
pro-cover corresponding to the closed subgroup
\(\beta(D_R^P)\). Choose a geometric lift in the
universal pro-cover defined by \(G_Q\). Its
stabilizer is the image in \(G_Q\) of an actual
target decomposition subgroup, and is contained in
\[
        \bigcap_{H'\supset\beta(D_R^P)}H'
                                      =\beta(D_R^P).
\]
Both groups map isomorphically onto the constant-field
Galois group, so this containment is equality.
Equivalently, \(\beta\) carries \(D_R^P\), up to
conjugacy, to the decomposition subgroup of an actual
noncuspidal target point.

For the full proper prime-to-\(p\) group, the
exceptional geometric point set is finite
\cite[Proposition~1.8(vi)]{AR_ST2009}.
Outside that set, a proper decomposition subgroup
has a unique point label. Since \(\beta\) is over
the identity of \(\Pi_C\), the target point in
Step~3 is therefore \(R\) itself whenever \(R\)
is outside this finite set and \(\{P,Q\}\).
This statement applies to every closed point after
its finite constant extension.

\emph{4. A difference cocycle turns point transport
into constant Kummer evaluations.}
Let \(M=\mathbb Z_r(1)\). Denote by
\[
  1\longrightarrow M\longrightarrow E_P
     \longrightarrow\Pi_C\longrightarrow1,
  \qquad
  1\longrightarrow M\longrightarrow E_Q
     \longrightarrow\Pi_C\longrightarrow1
                                                        \tag{AG.40.2}
\]
the maximal geometrically cuspidally central
quotients of \(G_P,G_Q\). Their geometric
extension classes are the degree-one cycle classes
of \(P\) and \(Q\) in
\(H^2(\Delta_C,M)\simeq\mathbb Z_r\).
Here is a construction avoiding any assertion
about a completed configuration or torsor fiber.
Use the Kummer cycle class of
\(\mathcal O_C(P)\) in
\(H^2(\Pi_C,M)\) to define the extension.
Its restriction to \(C-P\) is canonically
split by the canonical section of that line
bundle, giving a homomorphism from the
punctured-curve group to the extension.
The local residue maps the puncture inertia
onto the canonical generator of \(M\).
Conversely, the maximal geometrically central
cuspidal kernel is generated topologically by
this one inertia image: the full cuspidal
kernel is normally generated by the puncture
inertia, and geometric conjugation has been
killed. Hence the displayed homomorphism
identifies this central kernel with \(M\).
The same construction applies to \(Q\).
The geometric degree can also be checked directly
from the punctured surface presentation
\cite[XIII, Corollary~2.12]{AR_SGA1}: the puncture
word is the product of the \(g\) handle commutators.
After geometric cuspidal centralization it is one central
generator. Its image has infinite \(r\)-adic order:
send the first handle generators to \(1+e_{12}\)
and \(1+e_{23}\) in the upper-unitriangular
Heisenberg group over \(\mathbb Z_r\), and the
remaining handle generators to one. The puncture word
maps to \(1+e_{13}\). The resulting central relator
is the standard degree-one orientation extension.
This verifies the geometric degree-one class directly,
in addition to the cycle-class construction.
These are one-dimensional Kummer and localization
constructions, with no configuration-space fiber
injectivity input.

The isomorphism \(\beta\) induces an extension
isomorphism \(\bar\beta:E_P\to E_Q\).
Its multiplier on the kernel \(M\) is one:
it is over the identity on \(\Delta_C\), and
both degree-one extension classes are the same
generator of \(H^2(\Delta_C,M)\).
On the fiber product \(E_P\times_{\Pi_C}E_Q\)
define
\[
                 \chi(u,v)=\bar\beta(u)v^{-1}\in M.
                                                        \tag{AG.40.3}
\]
This is a continuous crossed homomorphism for the
cyclotomic action. Pull it back along the natural
map from the arithmetic prime-to-\(p\) group of
\(C-\{P,Q\}\) to that fiber product. Its residues
at \(P,Q\) are respectively \(+1,-1\), with
the canonical Kummer orientations. Therefore
\[
                         \kappa_r(f)-p^e\chi
                                                        \tag{AG.40.4}
\]
has zero residues, kills both inertia groups, and
descends to a class in \(H^1(\Pi_C,M)\).

This last cohomology group consists of constant
Kummer classes. For completeness, the
Hochschild--Serre exact sequence gives
\[
 H^1(G_k,M)\xrightarrow{\sim}H^1(\Pi_C,M),
                                                        \tag{AG.40.5}
\]
because \(H^1(\Delta_C,M)^{G_k}=0\): its
Frobenius eigenvalues have nonzero Weil weight,
so none is one. Finite-field
Kummer theory represents the class of (AG.40.4) by
some \(a\in k^\times\).

We check that the evaluation of \(\chi\) at
each point \(R\) from Step~3 is zero as a
cohomology class. Work over its finite residue
field so that \(R\) is rational. The target
proper point section satisfies
\[
                       N_{\Pi_C}(D_R)=D_R.
                                                        \tag{AG.40.6}
\]
Indeed a normalizer element can be divided by
the unique element of \(D_R\) with the same
arithmetic projection; the quotient is geometric
and centralizes \(D_R\), hence is trivial by
the standard Frobenius-weight centralizer
argument (equivalently
\cite[Lemma~1.7]{AR_ST2009}).

Choose a conjugator between
\(\beta(D_R^P)\) and the target subgroup
\(D_R^Q\). Its projection to \(\Pi_C\)
lies in the normalizer in (AG.40.6). Remove that
projection by multiplying by the corresponding
element of \(D_R^Q\). The resulting conjugator
lies in the cuspidal kernel. Its image in
\(E_Q\) lies in \(M\); conjugation by it
changes a section only by a coboundary.
Consequently (AG.40.3) restricts to a coboundary
on the actual point section of \(R\).
Evaluating (AG.40.4) thus gives
\[
       \kappa_{r,k(R)}\bigl(f(R)/a\bigr)=0
                                                        \tag{AG.40.7}
\]
for every closed \(R\) outside the fixed finite
exceptional set and \(\{P,Q\}\).

\emph{5. Constant Kummer evaluations force a
constant function.}
If \(f\) were nonconstant, its map from
\(C_{\bar k}\) to \(\mathbb P^1_{\bar k}\)
would be surjective. Choose a nontrivial
\(r\)-power root of unity \(\zeta\) such that
\(a\zeta\) avoids the finitely many images
of the excluded points. There are infinitely
many choices because \(r\ne p\).
Take a closed point \(R\) with
\(f(R)=a\zeta\). Its finite residue field
contains \(\zeta\), whose nontrivial
\(r\)-primary Kummer class in that field
contradicts (AG.40.7). Thus \(f\) is constant,
contradicting (AG.40.1) for \(P\ne Q\).
This proves the conditional criterion.
\end{proof}

\paragraph{The separate existence construction and source scope.}
Claims~\ref{clm:AGfixedrconfiguration} and
\ref{clm:AGcuspidaltransport} below supply the required
cuspidal isomorphism from conjugate proper point sections,
using a direct one-prime relative configuration argument.
They do not use the full prime-to-\(p\) configuration
fiber assertion excluded by AG39.
For clarity about the earlier literature route,
Saidi--Tamagawa's Corollary~2.3 constructs the compatible
cuspidally abelian comparison only when the chosen cusp
coordinates lie outside its exceptional sets; its
Theorem~2.7 refers to that comparison. The full
configuration-fiber assertion in the cited Mochizuki
author versions requires a quotient-convention or
source-scope reconciliation with AG39. Neither an
unverified exceptional-coordinate extension of those
statements nor that full configuration assertion is an
input to AG40--AG42.


% Integrated after independent root, criteria, and geometric audits, 2026-11-16.
% Original author draft retained in agent_F/F_fixed_r_configuration_draft.tex.
\subsection{A fixed-prime repair of configuration fiber exactness}

\begin{claim}{AG.41}{Proved: fixed-prime cuspidal configuration exactness}
\label{clm:AGfixedrconfiguration}
Let \(p\) and \(r\) be distinct primes. Let \(C\) be a
smooth proper connected complex curve of genus at least two,
let \(x\in C\), and set
\[
 U=(C\times C)-\Delta_C,\qquad V=C-\{x\},\qquad
 \Delta=\pi_1(C)^{(p')},\quad F=\pi_1(V)^{(p')},\quad
 E=\pi_1(U)^{(p')}.
                                                        \tag{AG.41.1}
\]
The fiber embedding uses the second projection \(U\to C\).
Define
\[
\begin{aligned}
 K_F&=\ker(F\to\Delta),
 &K_E&=\ker(E\to\Delta\times\Delta),\\
 F_{\mathrm{c}r}&=F/\ker(K_F\to K_F^{(r)}),
 &E_{\mathrm{c}r}&=E/\ker(K_E\to K_E^{(r)}).
\end{aligned}
                                                        \tag{AG.41.2}
\]
Then the fiber map induces an exact sequence
\[
 1\longrightarrow F_{\mathrm{c}r}
 \longrightarrow E_{\mathrm{c}r}
 \xrightarrow{\mathrm{pr}_2}\Delta\longrightarrow1.
                                                        \tag{AG.41.3}
\]
The same statement holds for geometric etale prime-to-\(p\)
groups over every algebraically closed field of characteristic
\(p\), with the corresponding characteristic relative quotients.

More precisely, every finite quotient of \(F_{\mathrm{c}r}\)
is dominated, on the fiber, by a finite quotient of
\(E_{\mathrm{c}r}\). The finite quotient of the total group
can be chosen to have a normal \(r\)-subgroup whose quotient
is a finite prime-to-\(p\) quotient of
\(\Delta\times\Delta\).
\end{claim}

The distinction from full prime-to-\(p\) exactness is essential.
Here the entire cuspidal kernel, on both sides, is replaced by
its maximal pro-\(r\) quotient for one fixed prime \(r\).
No injection of \(F\) into \(E\) is assumed or concluded.

\begin{proof}
\emph{1. A pro-\(r\) monodromy lemma without a section.}
Let \(1\to L\to H\to B\to1\) be an exact sequence of
discrete groups, where \(L\) is a free group of finite rank
at least two. Suppose the image of conjugation on
\(L/[L,L]L^r\) is an \(r\)-group. Then
\[
                         L^{(r)}\longrightarrow H^{(r)}
                    \quad\hbox{is injective}.             \tag{AG.41.4}
\]
Indeed, conjugation by \(H\) acts canonically on the
characteristic completion \(L^{(r)}\). The kernel of
\[
 \operatorname{Aut}(L^{(r)})\longrightarrow
 \operatorname{Aut}(L^{(r)}/\Phi(L^{(r)}))
                                                        \tag{AG.41.5}
\]
is pro-\(r\). Here is the finite argument, including the
topological passage. For a finite \(r\)-group \(P\), fix
an ordered basis of \(P/\Phi(P)\), of size \(d\).
There are exactly \(|\Phi(P)|^d\) tuples lifting that basis,
and every such tuple generates \(P\), by the elementary
Frattini criterion. The kernel of the automorphism action
on \(P/\Phi(P)\) acts freely on these tuples: an
automorphism fixing a generating tuple is the identity.
Its order therefore divides \(|\Phi(P)|^d\), so it is an
\(r\)-group. The iterated Frattini subgroups of the
finitely generated pro-\(r\) group \(L^{(r)}\) are
characteristic, open, and cofinal: their images in every
finite \(r\)-quotient eventually vanish. Applying the
finite argument to these quotients proves (AG.41.5).

It follows that the closure of the conjugation image of
\(H\) in \(\operatorname{Aut}(L^{(r)})\) is pro-\(r\).
The homomorphism from \(H\) to that closure extends to
\(H^{(r)}\). On \(L^{(r)}\) its composite is the inner
automorphism map. A nonabelian free pro-\(r\) group has
trivial center, so this composite is injective; see, for
example, \cite[Theorem~5.1, proof Case~I, and
Corollary~5.3]{FF_KochloukovaZalesskii2011}.
This proves (AG.41.4). In particular no splitting of
\(H\to B\), topological section, or lift of an outer
action is needed.

The compact embedding in (AG.41.4) induces the full
pro-\(r\) topology on the fiber. Explicitly, if
\(W\triangleleft L^{(r)}\) is open, there is an open
normal subgroup \(N\triangleleft H^{(r)}\) with
\(N\cap L^{(r)}\subset W\). Thus the finite quotient
\(H^{(r)}/N\) dominates the prescribed quotient of the
fiber; equality with that quotient is unnecessary.

\emph{2. Finite proper covers make the monodromy an
\(r\)-group modulo \(r\).}
Let \(a:C'\to C\) be a connected finite Galois etale
cover with group \(A\) of order prime to \(p\). Form
\[
 Y=(C'\times C')-(a\times a)^{-1}(\Delta_C).
                                                        \tag{AG.41.6}
\]
The map \(Y\to U\) is Galois with group \(A\times A\).
The second projection \(Y\to C'\) has fiber
\(V'=C'-a^{-1}(x)\), after choosing a point over \(x\)
in the second factor. Its removed divisor is the disjoint
union of the graphs of all deck transformations of \(a\).
Thus every puncture is labeled throughout this family.
The topological fibration is locally trivial; since the
closed surface \(C'\) is aspherical, its homotopy sequence
gives
\[
 1\longrightarrow\pi_1(V')\longrightarrow\pi_1(Y)
 \longrightarrow\pi_1(C')\longrightarrow1.
                                                        \tag{AG.41.7}
\]
The fiber group is free of rank
\(2g(C')+|A|-1\), in particular of rank at least two.

In homology with coefficients in \(\mathbb F_r\), the
boundary loops span a subspace \(W\) and give an exact
sequence
\[
 0\longrightarrow W\longrightarrow H_1(V',\mathbb F_r)
 \longrightarrow H_1(C',\mathbb F_r)\longrightarrow0.
                                                        \tag{AG.41.8}
\]
The boundary labels and complex orientations make monodromy
the identity on \(W\). The compactified family is the
product \(C'\times C'\to C'\), so monodromy is also
the identity on the last term. Therefore its image lies
in the additive shear group
\(\operatorname{Hom}_{\mathbb F_r}(H_1(C',\mathbb F_r),W)\),
a finite \(r\)-group. Fiber conjugations already act
trivially on homology. Step~1 consequently proves
\[
             \pi_1(V')^{(r)}\hookrightarrow\pi_1(Y)^{(r)}.
                                                        \tag{AG.41.9}
\]

\emph{3. Every finite relative fiber quotient is detected.}
Take a continuous surjection
\(q:F_{\mathrm{c}r}\to Q\), with \(Q\) finite, and
put \(R=q(K_F^{(r)})\). This is a normal \(r\)-subgroup
of \(Q\); the quotient \(A=Q/R\) is a finite quotient
of \(\Delta\). Use its Galois proper cover \(C'\to C\)
in Step~2. On discrete fundamental groups write
\(\Gamma=\pi_1(U)\), \(\Gamma_0=\pi_1(Y)\), and
\(\Lambda=\pi_1(V)\), \(\Lambda_0=\pi_1(V')\).
Then \(\Gamma_0\triangleleft\Gamma\),
\(\Gamma/\Gamma_0=A\times A\), and
\(\Lambda\cap\Gamma_0=\Lambda_0\).
The restriction of \(q\) to \(\Lambda_0\) has image
in \(R\), so factors through a finite \(r\)-quotient
of \(\Lambda_0^{(r)}\).

By (AG.41.9) and the final paragraph of Step~1 there is
a finite \(r\)-quotient of \(\Gamma_0\), with kernel
\(J\triangleleft\Gamma_0\), such that
\[
                         J\cap\Lambda_0\subset\ker(q).
                                                        \tag{AG.41.10}
\]
Take its normal core in \(\Gamma\):
\[
                  J^*=\bigcap_{g\in\Gamma/\Gamma_0}gJg^{-1}.
                                                        \tag{AG.41.11}
\]
This finite intersection is legitimate because \(J\) is
normal in \(\Gamma_0\). The quotient
\(\Gamma_0/J^*\) embeds in a finite product of
\(r\)-groups, hence is an \(r\)-group. Consequently
\(\Gamma/J^*\) is a finite prime-to-\(p\) group with
normal \(r\)-subgroup \(\Gamma_0/J^*\) and quotient
\(A\times A\). Moreover
\[
                 J^*\cap\Lambda\subset\ker(q),
                                                        \tag{AG.41.12}
\]
because \(J^*\subset\Gamma_0\) and (AG.41.10) applies.
Thus this total-space quotient dominates \(q\) on the
fiber. The inclusion of finite kernels continues to hold
after prime-to-\(p\) completion: it can be checked in
the joint finite image of the two quotient maps, in which
the discrete group has dense, hence surjective, image.

The quotient \(\Gamma\to\Gamma/J^*\) factors through
\(E_{\mathrm{c}r}\), not merely through \(E\).
Indeed its image of \(K_E\) lies in
\(\Gamma_0/J^*\), since \(\Gamma/\Gamma_0\) comes
from the proper product cover. That image is therefore an
\(r\)-group, and the kernel of \(K_E\to K_E^{(r)}\)
is killed. This proves the claimed cofinal finite-quotient
statement.

\emph{4. Injectivity and exactness.}
The fiber map factors through \(F_{\mathrm{c}r}\),
because it takes \(K_F\) to the pro-\(r\) kernel in
\(E_{\mathrm{c}r}\). Every finite quotient of that fiber
group is dominated by one from Step~3, so its kernel is
trivial. This proves the first arrow of (AG.41.3).

The ordinary topological fibration
\(C-\{x\}\to U\to C\) is right exact on fundamental
groups. Right exactness survives maximal prime-to-\(p\)
completion: quotienting \(E\) by the closed image of
the fiber has exactly the finite prime-to-\(p\) quotients
of \(\pi_1(C)\), hence is \(\Delta\). The subgroup
killed to form \(E_{\mathrm{c}r}\) is contained in
this fiber kernel. The image of the compact group
\(F_{\mathrm{c}r}\) is closed and is therefore exactly
the kernel of \(\mathrm{pr}_2\). This proves (AG.41.3).

\emph{5. Characteristic-\(p\) specialization.}
Spread out the pointed curve over a finitely generated
subfield, pass to a countable algebraic closure, and lift
the smooth pointed proper curve to a mixed-characteristic
complete discrete valuation ring. Curve deformations are
unobstructed, and proper formal lifts algebraize. An
algebraic closure of the fraction field embeds abstractly
in \(\mathbb C\). Both
\((\mathcal C,\mathcal x)\) and
\((\mathcal C\times\mathcal C,\Delta_{\mathcal C})\)
are smooth proper pairs with relative normal-crossing
boundary. SGA~1 XIII, Corollary~2.9 and Section~2.10 give
functorial isomorphisms of their geometric prime-to-\(p\)
groups \cite{AR_SGA1}; this is the same comparison used
in Claim~\ref{clm:AGconfigurationodd}.
These isomorphisms identify \(F,E,\Delta\), their proper
quotients and the fiber maps. They consequently identify
the characteristic relative quotients (AG.41.2).
The complex result therefore proves the asserted
characteristic-\(p\) result.
\end{proof}

\paragraph{Arithmetic consequence, with the section specified.}
Let \(C/k\) instead be a smooth proper geometrically
connected curve of genus at least two over a finite field
of characteristic \(p\). Write \(\Pi_C\) for its
arithmetic geometrically prime-to-\(p\) group and let
\(\mathcal E\) be the corresponding group of
\(U=(C\times C)-\Delta_C\). Form its cuspidally
pro-\(r\) quotient relative to
\[
                \Pi_C\times_{G_k}\Pi_C.
                                                        \tag{AG.41.13}
\]
The kernel relative to this product is geometric, so this
quotient has geometric subgroup \(E_{\mathrm{c}r}\)
from (AG.41.2). If \(P\in C(k)\) and
\(s_P:G_k\to\Pi_C\) is its actual point section,
the natural map from the arithmetic cuspidally pro-\(r\)
group of \(C-P\) identifies it with
\[
       \mathcal E_{\mathrm{c}r}
                  \times_{\Pi_C,\mathrm{pr}_2,s_P}G_k.
                                                        \tag{AG.41.14}
\]
Indeed on geometric kernels this is (AG.41.3), and on
arithmetic quotients it is the identity of \(G_k\).
The elementary extension comparison therefore makes it
an isomorphism. Since a section is injective, the right
side can equivalently be viewed as the inverse image of
\(s_P(G_k)\) under the second projection. Choices of
basepoint paths give the usual geometric inner ambiguity;
no section of the configuration fibration was postulated.

\paragraph{The boundary comparison.}
Claim~\ref{clm:AGcuspidaltransport} identifies the actual
puncture decomposition group with the intersection with
the global diagonal boundary decomposition group. This
is a separate tame local calculation; it is not silently
included in the geometric fiber theorem.

% Integrated after independent root, criteria, and geometric audits, 2026-11-16.
% Original author draft retained in agent_F/F_fixed_r_boundary_draft.tex.
\subsection{The diagonal boundary and transport of cuspidal groups}

Retain the arithmetic notation of
\ref{clm:AGfixedrconfiguration}: \(C/k\) is smooth proper
of genus at least two over a finite field of characteristic
\(p\), \(r\ne p\), and
\(\mathcal E_{\mathrm{c}r}\) is the cuspidally pro-\(r\)
quotient of the geometrically prime-to-\(p\) group of
\(U=(C\times C)-\Delta_C\), relative to
\(\Pi_C\times_{G_k}\Pi_C\).

\begin{claim}{AG.42}{Proved: diagonal boundary and cuspidal transport}
\label{clm:AGcuspidaltransport}
There is a diagonal boundary decomposition subgroup
\(\mathcal D\subset\mathcal E_{\mathrm{c}r}\) with
an exact sequence
\[
 1\longrightarrow\mathbb Z_r(1)\longrightarrow\mathcal D
 \longrightarrow\Pi_C\longrightarrow1,
                                                        \tag{AG.42.1}
\]
whose map to the proper product is the diagonal map on
the last term. For every rational point \(P\), choose
compatible geometric paths through this boundary and put
\(D_P=s_P(G_k)\subset\Pi_C\) and
\(\mathcal E_P=\mathrm{pr}_2^{-1}(D_P)\).
The fixed-\(r\) configuration theorem identifies
\(\mathcal E_P\) with the arithmetic cuspidally
pro-\(r\) group of \(C-P\), and under this
identification its puncture decomposition group is
\[
                           \mathcal D\cap\mathcal E_P.
                                                        \tag{AG.42.2}
\]

If two rational points \(P,Q\) have geometrically conjugate
proper sections in \(\Pi_C\), there is an isomorphism
between the actual cuspidally pro-\(r\) groups of
\(C-P\) and \(C-Q\), over the identity of
\(\Pi_C\), which carries the first puncture decomposition
group to a conjugate of the second. Consequently
\(P=Q\): rational point sections of the full proper
geometrically prime-to-\(p\) fundamental group are
injective up to geometric conjugacy.
\end{claim}

\begin{proof}
\emph{Construction of the boundary group.}
We give the local tame-cover construction explicitly to
isolate it from any full configuration-fiber assertion.
Let \(X=C\times C\), \(D=\Delta_C\), and
\(N=\mathcal O_X(D)|_D\), the normal line bundle.
For each positive integer \(n\) prime to \(p\), form
the root stack \(X_n=\sqrt[n]{(X,D)}\). Locally where
\(D\) has parameter \(t\), this is the quotient of
\(\operatorname{Spec}(A[z]/(z^n-t))\) by \(\mu_n\).
The reduced inverse image \(D_n\) of \(D\) is the
\(\mu_n\)-gerbe of \(n\)-th roots of \(N\).

A finite etale cover of \(U\) with prime-to-\(p\)
geometric monodromy has tame inertia along \(D\).
Take \(n\) divisible by its finitely many inertia
orders. The absolute Abhyankar lemma, SGA~1 XIII,
Proposition~5.2, says that after the local substitution
\(z^n=t\), its normalization extends uniquely to a
finite etale cover of the whole Kummer chart
\cite{AR_SGA1}. The cited statement already includes
extension over the entire chart; its proof applies
purity after the codimension-one tame calculation.
The unique extensions and their deck
descent data glue to a finite etale cover of \(X_n\).
Restriction to \(D_n\) gives a finite etale cover of
that gerbe. All constructions commute with increasing
\(n\), with field extension, and with pullback
transverse to \(D\). They therefore define the boundary
homomorphism to the tame fundamental group of \(U\),
and hence to its geometrically prime-to-\(p\) quotient.

Here are the exact group assertions needed from this
construction. The fundamental group of the gerbe \(D_n\)
maps onto \(\pi_1(D)\), and its kernel is the image
of \(\mu_n\). This is the elementary gerbe exact
sequence: etale locally on \(D\) the gerbe is
\(B\mu_n\); a cover on which the stabilizer acts
trivially descends to \(D\), and covers of \(D\)
pull back fully faithfully. The inertia is normal,
with the cyclotomic descent action. No injection of
\(\mu_n\) is required here. Passing to compatible
root orders gives the tame boundary group \(\mathcal T\)
and a right exact sequence
\[
 \widehat{\mathbb Z}^{(p')}(1)\longrightarrow\mathcal T
              \longrightarrow\pi_1(D)\longrightarrow1.
                                                        \tag{AG.42.3}
\]
The exactness follows either from the same descent
criterion or by inverse limits of the compact gerbe
sequences; the transition maps are surjective, as is
seen in the local maps \(\mu_m\to\mu_n\) for
\(n\mid m\).

Now take the maximal geometrically prime-to-\(p\)
quotients. On geometric groups, right exactness of
completion shows that quotienting the completed
boundary group by the image of inertia gives
\(\Delta_C\). On arithmetic groups the quotients
over \(G_k\) are unchanged. Thus
\[
 \widehat{\mathbb Z}^{(p')}(1)\longrightarrow
 \mathcal T^{\mathrm{geom},p'}\longrightarrow
 \Pi_C\longrightarrow1                             \tag{AG.42.4}
\]
is right exact. This step explicitly rules out an extra
kernel arising from the discarded geometric base kernel.

Let \(\mathcal D\) be the image of
\(\mathcal T^{\mathrm{geom},p'}\) in
\(\mathcal E_{\mathrm{c}r}\). Its projection to the
proper product factors through the injective diagonal
\(\Pi_C\to\Pi_C\times_{G_k}\Pi_C\).
Consequently the kernel of \(\mathcal D\to\Pi_C\)
is exactly the image of the inertia in (AG.42.4).
Since it is contained in the cuspidal pro-\(r\) kernel,
this inertia image is a quotient of \(\mathbb Z_r(1)\).

This quotient is the full \(\mathbb Z_r(1)\).
To verify this without using a global boundary theorem,
take a geometric fiber \(C-\{P\}\). Its cuspidal
pro-\(r\) group maps onto its maximal pro-\(r\)
fundamental group, which is free on the usual
\(2g\) surface generators. Its puncture inertia
generator is, up to inversion, the product
\(\prod_{i=1}^g[a_i,b_i]\). Map this free group to
the pro-\(r\) upper-unitriangular Heisenberg group
over \(\mathbb Z_r\) by
\[
 a_1\longmapsto 1+e_{12},\qquad
 b_1\longmapsto 1+e_{23},\qquad
 a_i,b_i\longmapsto1\quad(i>1).
                                                        \tag{AG.42.5}
\]
The puncture generator maps to \(1+e_{13}\), of
infinite \(r\)-adic order. Thus its \(r\)-inertia
is \(\mathbb Z_r(1)\). The fixed-\(r\) fiber
injection carries this inertia injectively into
\(\mathcal E_{\mathrm{c}r}\). In the local Kummer
charts the fiber meets the diagonal transversely
with multiplicity one, so this is exactly the
diagonal inertia in (AG.42.4), with its ordinary
Kummer generator. This proves (AG.42.1).

Finally, the same transverse local chart sends the
actual puncture decomposition group of \(C-P\)
into \(\mathcal D\). It lies over \(D_P\), and
it maps isomorphically on inertia and on \(D_P\).
The inverse image of \(D_P\) in (AG.42.1) is an
extension of those same two groups. The resulting
map is therefore an isomorphism onto that inverse
image, proving (AG.42.2). Compatible paths can be chosen
simultaneously using the common boundary gerbes;
other choices give the permitted geometric inner
conjugacies.

\emph{Transport from conjugate proper point sections.}
Suppose \(P,Q\in C(k)\) have geometrically conjugate
proper sections. Choose
\(\gamma\in\Delta_C\) with
\(\gamma D_P\gamma^{-1}=D_Q\).
By (AG.42.1), choose a geometric
\(d\in\mathcal D\) mapping to \(\gamma\).
Conjugation by \(d\) identifies
\(\mathcal E_P\) with \(\mathcal E_Q\), and by
(AG.42.2) carries their puncture decomposition groups
onto each other. On the first proper coordinate
it induces \(\operatorname{Inn}(\gamma)\).

The first projection
\(\mathcal E_Q\to\Pi_C\) is surjective, and is
surjective on geometric kernels: it is the proper
quotient of the punctured curve group, by the
fixed-\(r\) theorem. Choose a geometric
\(e\in\mathcal E_Q\) mapping to \(\gamma^{-1}\)
under that first projection. Then
\[
            \beta=\operatorname{Inn}(e)\circ
                         \operatorname{Inn}(d):
                 \mathcal E_P\xrightarrow{\sim}\mathcal E_Q
                                                        \tag{AG.42.6}
\]
is over the identity of \(\Pi_C\), is over the
identity of \(G_k\), and carries the first puncture
decomposition group to a conjugate of the second
inside \(\mathcal E_Q\). These are exactly the
hypotheses of Claim~\ref{clm:AGcuspidalcriterion}.
The construction uses the one-prime relative fiber
comparison of Claim~\ref{clm:AGfixedrconfiguration}
and the tame boundary comparison just proved; it does
not invoke full prime-to-\(p\) configuration fiber
exactness. Claim~\ref{clm:AGcuspidalcriterion} now
implies \(P=Q\), as asserted. The argument applies
unchanged after every finite extension of the constant
field.
\end{proof}

